% Complex Numbers — Pure Mathematics with Mechanics Upper Sixth — Cours signé OnBuch+
% Official MINESEC syllabus. Compiler avec : tectonic PMM02-complex-numbers.tex
\newcommand{\DOCMATIERE}{Pure Mathematics with Mechanics}
\newcommand{\DOCNIVEAU}{Upper Sixth}
\newcommand{\DOCMODULE}{Upper Sixth — Pure mathematics with mechanics}
\newcommand{\DOCLECON}{2}
\newcommand{\DOCTITRE}{Complex Numbers}
% OnBuch+ — préambule « cours signés » ENRICHI (compilé avec tectonic / XeLaTeX)
% Système visuel « Pop » d'OnBuch+ : fond crème (jamais blanc), bordures encre
% épaisses, ombres « dures » décalées, titres Archivo Black en capitales.
% Version MATHÉMATIQUES : graphes (pgfplots/tikz), boîtes pédagogiques
% (définition, propriété, méthode, attention, le savais-tu, exercice résolu,
% exercice, TP, bilan), page de garde et signature OnBuch+ ancrée en bas.
%
% Macros attendues dans main.tex AVANT \input{preamble} :
%   \DOCMATIERE  \DOCNIVEAU  \DOCMODULE  \DOCLECON (numéro)  \DOCTITRE
\documentclass[11pt]{article}

\usepackage{fontspec}
\usepackage[english]{babel}
\usepackage[a4paper,margin=2cm,top=3.4cm,bottom=2.8cm,headheight=1.4cm,footskip=1.3cm]{geometry}
\usepackage{amsmath,amssymb,amsfonts}
\usepackage{mathtools} % \xmapsto, \coloneqq... souvent utilisés par les modèles
\usepackage{cancel} % simplifications barrées (bilans de forces)
% \abs{z} (module, valeur absolue) et \norm{u} (norme), utilisés par les modèles
\providecommand{\abs}{}\let\abs\relax\DeclarePairedDelimiter{\abs}{\lvert}{\rvert}
\providecommand{\norm}{}\let\norm\relax\DeclarePairedDelimiter{\norm}{\lVert}{\rVert}
\usepackage[table]{xcolor} % [table] : fournit aussi \rowcolors (alternance de lignes)
\usepackage{pagecolor}
\usepackage{tikz}
\usetikzlibrary{arrows.meta,positioning,calc,shapes.geometric,shapes.misc,shapes.arrows,decorations.pathmorphing,decorations.pathreplacing,decorations.markings,patterns,shadows,mindmap,trees,fit,backgrounds,matrix,intersections,angles,quotes}
\usepackage{pgfplots}
\usepackage[version=4]{mhchem} % équations nucléaires \ce{^{235}_{92}U}
\usepackage[european,siunitx]{circuitikz} % schémas électriques
\pgfplotsset{compat=1.18}
\usepgfplotslibrary{fillbetween,groupplots} % aires entre courbes (calcul intégral)
\usepackage{enumitem}
\usepackage{fancyhdr}
% [most] charge toutes les bibliothèques tcolorbox (listings, minted, raster,
% documentation...) et gonfle inutilement la mémoire TeX sur un document long
% (~300 boîtes) : on ne charge que ce qui est réellement utilisé.
\usepackage[skins,breakable]{tcolorbox}
\usepackage{graphicx}
\usepackage{colortbl}
\usepackage{booktabs}
\usepackage{tabularx}
\usepackage{diagbox} % case d'en-tête diagonale des tableaux à double entrée
% Colonnes extensibles centrée / gauche / droite (C, L, R), souvent utilisées par les modèles
\newcolumntype{C}{>{\centering\arraybackslash}X}
\newcolumntype{L}{>{\raggedright\arraybackslash}X}
\newcolumntype{R}{>{\raggedleft\arraybackslash}X}
\usepackage{array}
\usepackage{multicol}
\usepackage{siunitx}
% parse-numbers=false : accepte aussi \SI{2,5 \times 10^{-3}}{...}
\sisetup{locale=UK,output-decimal-marker={.},group-minimum-digits=4,parse-numbers=false}
\DeclareSIUnit\ohm{\ensuremath{\symup{\Omega}}} % U+2126 absent de la police
\DeclareSIUnit\division{div} % réglages d'oscilloscope (ms/div, V/div)
\DeclareSIUnit\tr{tr} % tours (vitesse de rotation, ex. \SI{1500}{\tr\per\minute})
\usepackage{qrcode}
% \degree : commande fréquemment utilisée par les modèles pour un angle ou une
% température en dehors de \SI{}{} (non fournie par les paquets chargés).
\providecommand{\degree}{^{\circ}}
% \qed (fin de démonstration), utilisé par les modèles sans amsthm
\providecommand{\qed}{\hfill$\square$}
% \barre : double barre des valeurs interdites dans un tableau de variations (tkz-tab)
\providecommand{\barre}{\ensuremath{\|}}
% environnement proof (amsthm non chargé) utilisé par les modèles
\ifdefined\proof\else\newenvironment{proof}[1][Proof]{\par\noindent\textit{#1.}\ }{\qed\par}\fi
\usepackage{lastpage}
\usepackage{hyperref}
\hypersetup{hidelinks}

% ============================================================
% Palette « Pop » OnBuch+ — mêmes hex que la classe P (plus_ui.dart)
% ============================================================
\definecolor{poporange}{HTML}{F59321}
\definecolor{popdark}{HTML}{E07A0C}
\definecolor{popcream}{HTML}{FFF6E8}
\definecolor{popink}{HTML}{1C1714}
\definecolor{poppurple}{HTML}{7A5AE0}
\definecolor{popgreen}{HTML}{2E9E6A}
\definecolor{poppink}{HTML}{E0457B}
\definecolor{popblue}{HTML}{2F7DE1}
\definecolor{popgold}{HTML}{C98A12}
\definecolor{popmuted}{HTML}{8A7F76}
\definecolor{popline}{HTML}{EADFD2}
\definecolor{popgray}{HTML}{8A7F76}
\colorlet{popgrey}{popgray}
% Alias tolérants (le modèle nomme parfois les couleurs différemment)
\colorlet{popwhite}{white}
\colorlet{popblack}{popink}
\colorlet{popred}{poppink}
\colorlet{popyellow}{popgold}
% Teintes claires pour fonds de tableaux / zones
\colorlet{popblueL}{popblue!10!white}
\colorlet{poporangeL}{poporange!14!white}
\colorlet{popgreenL}{popgreen!12!white}
\colorlet{poppurpleL}{poppurple!12!white}
\colorlet{poppinkL}{poppink!12!white}
\colorlet{popgoldL}{popgold!14!white}

\pagecolor{popcream}

% ============================================================
% Typographie
% ============================================================
\defaultfontfeatures{Ligatures=TeX}
\setmainfont{PlusJakartaSans-Medium.ttf}[Path=../../../fonts/,BoldFont=PlusJakartaSans-Bold.ttf]
% Maths en sans-serif (Fira Math) avec chiffres et lettres droites de Plus
% Jakarta, pour que les formules se fondent dans le texte.
\usepackage[math-style=ISO,bold-style=ISO]{unicode-math}
\setmathfont{FiraMath-Regular.otf}[Path=../../../fonts/]
\setmathfont{PlusJakartaSans-Medium.ttf}[Path=../../../fonts/,range={up/num,up/{Latin,latin},\mathup}]
% (icomma retiré : décimales avec point en anglais)
% Caractères mathématiques Unicode absents des polices de texte (Plus Jakarta,
% Archivo Black) : rendus par leur équivalent mathématique, en texte comme en
% formule — sinon ils s'affichent en carré vide (ex. « Arithmétique dans ℤ »).
\usepackage{newunicodechar}
\newunicodechar{ℝ}{\ensuremath{\mathbb{R}}}
\newunicodechar{ℤ}{\ensuremath{\mathbb{Z}}}
\newunicodechar{ℂ}{\ensuremath{\mathbb{C}}}
\newunicodechar{ℕ}{\ensuremath{\mathbb{N}}}
\newunicodechar{ℚ}{\ensuremath{\mathbb{Q}}}
\newunicodechar{𝔻}{\ensuremath{\mathbb{D}}}
\newunicodechar{α}{\ensuremath{\alpha}}
\newunicodechar{β}{\ensuremath{\beta}}
\newunicodechar{γ}{\ensuremath{\gamma}}
\newunicodechar{δ}{\ensuremath{\delta}}
\newunicodechar{ε}{\ensuremath{\varepsilon}}
\newunicodechar{θ}{\ensuremath{\theta}}
\newunicodechar{λ}{\ensuremath{\lambda}}
\newunicodechar{μ}{\ensuremath{\mu}}
\newunicodechar{σ}{\ensuremath{\sigma}}
\newunicodechar{τ}{\ensuremath{\tau}}
\newunicodechar{φ}{\ensuremath{\varphi}}
\newunicodechar{ψ}{\ensuremath{\psi}}
\newunicodechar{ω}{\ensuremath{\omega}}
\newunicodechar{Σ}{\ensuremath{\Sigma}}
% majuscules produites par \MakeUppercase dans les titres (α-aminés -> Α)
\newunicodechar{Α}{\ensuremath{\alpha}}
\newunicodechar{Β}{\ensuremath{\beta}}
\newunicodechar{Γ}{\ensuremath{\Gamma}}
\newunicodechar{Δ}{\ensuremath{\Delta}}
\newunicodechar{Ε}{\ensuremath{\varepsilon}}
\newunicodechar{Θ}{\ensuremath{\Theta}}
\newunicodechar{Λ}{\ensuremath{\Lambda}}
\newunicodechar{Μ}{\ensuremath{\mu}}
\newunicodechar{Τ}{\ensuremath{\tau}}
\newunicodechar{Φ}{\ensuremath{\Phi}}
\newunicodechar{Ψ}{\ensuremath{\Psi}}
\newunicodechar{Ω}{\ensuremath{\Omega}}
\newunicodechar{↦}{\ensuremath{\mapsto}}
\newunicodechar{∈}{\ensuremath{\in}}
\newunicodechar{∉}{\ensuremath{\notin}}
\newunicodechar{∧}{\ensuremath{\wedge}}
\newunicodechar{⇔}{\ensuremath{\Leftrightarrow}}
\newunicodechar{⟺}{\ensuremath{\Longleftrightarrow}}
\newunicodechar{⇒}{\ensuremath{\Rightarrow}}
\newunicodechar{⟹}{\ensuremath{\Longrightarrow}}
\newunicodechar{∘}{\ensuremath{\circ}}
\newunicodechar{∪}{\ensuremath{\cup}}
\newunicodechar{∩}{\ensuremath{\cap}}
\newunicodechar{≡}{\ensuremath{\equiv}}
\newunicodechar{⊂}{\ensuremath{\subset}}
\newunicodechar{⊥}{\ensuremath{\perp}}
\newunicodechar{∃}{\ensuremath{\exists}}
\newunicodechar{∀}{\ensuremath{\forall}}
\newunicodechar{∛}{\ensuremath{\sqrt[3]{\,}}}
\newunicodechar{∜}{\ensuremath{\sqrt[4]{\,}}}
\newunicodechar{⃗}{\ensuremath{{}^{\to}}}
\newunicodechar{ⁿ}{\ifmmode^{n}\else\textsuperscript{\ensuremath{n}}\fi}
\newunicodechar{ˣ}{\ifmmode^{x}\else\textsuperscript{\ensuremath{x}}\fi}
\newunicodechar{ᵇ}{\ifmmode^{b}\else\textsuperscript{\ensuremath{b}}\fi}
\newunicodechar{ʳ}{\ifmmode^{r}\else\textsuperscript{\ensuremath{r}}\fi}
\newunicodechar{ᵘ}{\ifmmode^{u}\else\textsuperscript{\ensuremath{u}}\fi}
\newunicodechar{ʸ}{\ifmmode^{y}\else\textsuperscript{\ensuremath{y}}\fi}
\newunicodechar{ᵃ}{\ifmmode^{a}\else\textsuperscript{\ensuremath{a}}\fi}
\newunicodechar{ᵉ}{\ifmmode^{e}\else\textsuperscript{\ensuremath{e}}\fi}
\newunicodechar{ᵏ}{\ifmmode^{k}\else\textsuperscript{\ensuremath{k}}\fi}
\newunicodechar{ᵖ}{\ifmmode^{p}\else\textsuperscript{\ensuremath{p}}\fi}
\newunicodechar{ᴺ}{\ifmmode^{N}\else\textsuperscript{\ensuremath{N}}\fi}
\newunicodechar{ᶜ}{\ifmmode^{c}\else\textsuperscript{\ensuremath{c}}\fi}
\newunicodechar{ʰ}{\ifmmode^{h}\else\textsuperscript{\ensuremath{h}}\fi}
\newunicodechar{ᵠ}{\ifmmode^{q}\else\textsuperscript{\ensuremath{q}}\fi}
\newunicodechar{ₙ}{\ifmmode_{n}\else\textsubscript{\ensuremath{n}}\fi}
\newunicodechar{ₐ}{\ifmmode_{a}\else\textsubscript{\ensuremath{a}}\fi}
\newunicodechar{ᵢ}{\ifmmode_{i}\else\textsubscript{\ensuremath{i}}\fi}
\newunicodechar{ᵣ}{\ifmmode_{r}\else\textsubscript{\ensuremath{r}}\fi}
\newunicodechar{ₖ}{\ifmmode_{k}\else\textsubscript{\ensuremath{k}}\fi}
\newunicodechar{ᵦ}{\ifmmode_{\beta}\else\textsubscript{\ensuremath{\beta}}\fi}
\newunicodechar{ₓ}{\ifmmode_{x}\else\textsubscript{\ensuremath{x}}\fi}
\newfontfamily\popdisplay{ArchivoBlack-Regular.ttf}[Path=../../../fonts/]
\newfontfamily\poplogo{SpaceGrotesk-Bold.ttf}[Path=../../../fonts/]
\newfontfamily\popsemi{PlusJakartaSans-SemiBold.ttf}[Path=../../../fonts/]
\newfontfamily\popxbold{PlusJakartaSans-ExtraBold.ttf}[Path=../../../fonts/]
\newfontfamily\popxboldls{PlusJakartaSans-ExtraBold.ttf}[Path=../../../fonts/,LetterSpace=3.0]

\setlist[enumerate]{leftmargin=1.7em,itemsep=5pt,topsep=4pt}
\setlist[itemize]{leftmargin=1.7em,itemsep=4pt,topsep=4pt,label={\color{poporange}\textbullet}}
\setlength{\parindent}{0pt}
\setlength{\parskip}{5pt}
\renewcommand{\arraystretch}{1.25}

% pgfplots : style Pop par défaut
\pgfplotsset{
  every axis/.append style={axis line style={popink,line width=1pt},tick style={popink},
    grid style={popline,line width=0.5pt},label style={font=\small},tick label style={font=\footnotesize}},
  popcurve/.style={poporange,line width=1.8pt,no markers,smooth},
}
% Même style disponible dans un \draw TikZ simple (hors pgfplots)
\tikzset{popcurve/.style={poporange,line width=1.8pt}}
\tikzset{popgrid/.style={popline,very thin}}
\colorlet{popcurve}{poporange} % le nom du style est parfois utilisé comme couleur

% ============================================================
% Pill / squiggle
% ============================================================
\newcommand{\poppill}[3]{%
  \tikz[baseline=-0.55ex]\node[draw=popink,line width=1.2pt,rounded corners=2.6pt,%
    fill=#2,inner xsep=6.5pt,inner ysep=3pt]{%
    \popxboldls\fontsize{7}{7}\selectfont\color{#3}\MakeUppercase{#1}};%
}
\newcommand{\popsquiggle}[1]{%
  \tikz[baseline=-0.4ex]{
    \draw[#1,line width=2.6pt,line cap=round]
      plot [smooth] coordinates {(0,0.09) (0.34,0.01) (0.68,0.21) (1.02,0.01) (1.36,0.10)};
  }%
}
% Mot-clé surligné (marqueur orange sous le texte)
\newcommand{\cle}[1]{{\popxbold\color{popdark}#1}}

% ============================================================
% En-tête / pied
% ============================================================
\pagestyle{fancy}
\fancyhf{}
\renewcommand{\headrulewidth}{0pt}
\renewcommand{\footrulewidth}{0pt}
\lhead{\raisebox{-0.15ex}{\poplogo\fontsize{13}{13}\selectfont\color{popink}OnBuch\color{poporange}+}}
\rhead{\poppill{\DOCMATIERE\ \textbullet\ \DOCNIVEAU\ \textbullet\ Chapter \DOCLECON}{white}{popink}}
\lfoot{\popxbold\fontsize{7.6}{7.6}\selectfont\color{popmuted}\MakeUppercase{Signed course OnBuch+}\ \textbullet\ {\popsemi\DOCTITRE}}
\rfoot{\tikz[baseline=-0.6ex]\node[rounded corners=8.5pt,fill=popink,minimum height=17pt,minimum width=17pt,inner xsep=5pt,inner sep=0]{\popxbold\fontsize{8}{8}\selectfont\color{popcream}\thepage\,/\,\pageref*{LastPage}};}

% ============================================================
% Titres
% ============================================================
\newcommand{\chaptitre}[2]{%
  \begin{tcolorbox}[enhanced,colback=poporange,colframe=popink,boxrule=2.6pt,arc=11pt,
      drop shadow=popink,left=15pt,right=15pt,top=12pt,bottom=11pt,before skip=3pt,after skip=18pt]
    \poppill{#2}{popink}{popcream}\par\vspace{9pt}
    {\raggedright\hyphenpenalty=10000\exhyphenpenalty=10000\popdisplay\fontsize{22}{24}\selectfont\color{popcream}\MakeUppercase{#1}\par}
  \end{tcolorbox}
}
% Section numérotée : pastille orange avec le numéro + titre Archivo
\newcounter{popsec}
\newcommand{\coursec}[1]{%
  \stepcounter{popsec}\setcounter{popsub}{0}%
  \par\vspace{16pt}\noindent
  \begin{minipage}{\linewidth}
  \tikz[baseline=-0.7ex]\node[draw=popink,line width=1.6pt,fill=poporange,rounded corners=4pt,minimum size=22pt,inner sep=0]{\popdisplay\fontsize{12}{12}\selectfont\color{popcream}\thepopsec};%
  \hspace{8pt}{\popdisplay\fontsize{15}{17}\selectfont\color{popink}\MakeUppercase{#1}}\par
  \vspace{4pt}\noindent{\color{poporange}\rule{36pt}{3.6pt}}
  \end{minipage}\par\nopagebreak\vspace{8pt}
}
\newcounter{popsub}
\newcommand{\courssub}[1]{%
  \stepcounter{popsub}%
  \par\vspace{9pt}\noindent{\popxbold\fontsize{11.5}{13}\selectfont\color{popink}{\color{poporange}\thepopsec.\thepopsub}\hspace{6pt}#1}\par\nopagebreak\vspace{4pt}
}

% ============================================================
% Boîtes pédagogiques — même recette : fond blanc, bordure épaisse colorée,
% ombre dure encre, badge plein. Titre optionnel [#1].
% ============================================================
\tcbset{popbase/.style={breakable,enhanced,colback=white,boxrule=2.3pt,arc=9pt,
  shadow={3pt}{-3pt}{0pt}{popink},left=14pt,right=14pt,top=11pt,bottom=11pt,before skip=11pt,after skip=15pt,
  fontupper=\popsemi\fontsize{10.6}{14.5}\selectfont\color{popink},
  fontlower=\popsemi\fontsize{10.6}{14.5}\selectfont\color{popink}}}
% Coche dessinée (le glyphe ✓ n'existe pas dans les polices du document)
\newcommand{\popcheck}[1]{\tikz[baseline=-0.5ex]\draw[#1,line width=1.6pt,line cap=round,line join=round](-0.13,0.0)--(-0.04,-0.09)--(0.13,0.1);}
\providecommand{\coche}{\popcheck{popgreen}}
\providecommand{\resultat}[1]{\par\smallskip\noindent\fcolorbox{popgreen}{popgreenL}{\parbox{\dimexpr\linewidth-2\fboxsep-2\fboxrule\relax}{\textbf{Result.} #1}}\par\smallskip}
\newcommand{\popboxhead}[3]{\poppill{#1}{#2}{white}\ifx\relax#3\relax\else\hspace{6pt}{\popxbold\fontsize{10.6}{12}\selectfont\color{popink}#3}\fi\par\vspace{7pt}}

\newtcolorbox{aretenir}[1][]{popbase,colframe=popgreen,before upper={\popboxhead{Key points}{popgreen}{#1}}}
\newtcolorbox{exemplebox}[1][]{popbase,colframe=poporange,before upper={\popboxhead{Exemple}{poporange}{#1}}}
\newtcolorbox{definition}[1][]{popbase,colframe=popblue,before upper={\popboxhead{Definition}{popblue}{#1}}}
\newtcolorbox{propriete}[1][]{popbase,colframe=poppurple,before upper={\popboxhead{Property}{poppurple}{#1}}}
\newtcolorbox{methode}[1][]{popbase,colframe=popgold,before upper={\popboxhead{Method}{popgold}{#1}}}
\newtcolorbox{attention}[1][]{popbase,colframe=poppink,before upper={\popboxhead{Warning}{poppink}{#1}}}
\newtcolorbox{experience}[1][]{popbase,colframe=popblue,colback=popblueL,before upper={\popboxhead{Experiment}{popblue}{#1}}}
% « Le savais-tu ? » : carte violette pleine (contexte, culture, Cameroun)
\newtcolorbox{savaistu}[1][]{popbase,colback=poppurple,colframe=popink,
  fontupper=\popsemi\fontsize{10.4}{14.2}\selectfont\color{white},
  before upper={\poppill{Did you know?}{white}{poppurple}\ifx\relax#1\relax\else\hspace{6pt}{\popxbold\color{white}#1}\fi\par\vspace{7pt}}}
% Exercice résolu : énoncé en haut, \tcblower puis la solution sur fond crème
\newcounter{popexo}\newcounter{popexr}
\newtcolorbox{exoresolu}[1][]{popbase,colframe=poporange,colbacklower=popcream,
  segmentation style={popink,line width=1.2pt,dashed},
  before upper={\stepcounter{popexr}\popboxhead{Worked example \thepopexr}{poporange}{#1}},
  before lower={\poppill{Solution}{popink}{popcream}\par\vspace{6pt}}}
% Exercice (non résolu en place) — la correction est dans la partie Corrigés
\newtcolorbox{exercice}[1][]{popbase,colframe=popink,
  before upper={\stepcounter{popexo}\popboxhead{Exercise \thepopexo}{popink}{#1}}}
% Corrigé
\newtcolorbox{corrige}[1][]{popbase,colframe=popgreen,colback=popgreenL,
  before upper={\popboxhead{Solution}{popgreen}{#1}}}
% Objectifs / prérequis / bilan
\newtcolorbox{objectifs}{popbase,colframe=popink,colback=poporangeL,
  before upper={\popboxhead{Chapter objectives}{popink}{}}}
\newtcolorbox{prerequis}{popbase,colframe=popink,colback=popblueL,
  before upper={\popboxhead{Prerequisites}{popblue}{}}}
\newtcolorbox{bilan}{popbase,colframe=popink,colback=white,boxrule=2.8pt,
  before upper={\popboxhead{Fiche bilan}{poporange}{L'essentiel à connaître pour l'examen}}}
% Figure encadrée avec légende
\newtcolorbox{popfigure}[1][]{enhanced,breakable=false,colback=white,colframe=popline,boxrule=1.4pt,arc=8pt,
  left=10pt,right=10pt,top=10pt,bottom=8pt,before skip=10pt,after skip=12pt,halign=center,
  lower separated=false,halign lower=center,
  fontlower=\popsemi\fontsize{9}{11.5}\selectfont\color{popmuted},
  #1}
\newcommand{\legende}[1]{\par\vspace{6pt}{\popsemi\fontsize{9}{11.5}\selectfont\color{popmuted}#1}}

% ============================================================
% Signature OnBuch+
% ============================================================
\newcommand{\poplogoinline}[1]{{\poplogo\fontsize{#1}{#1}\selectfont\color{popink}OnBuch\color{poporange}+}}
\newcommand{\signatureonbuch}{%
  \begin{tcolorbox}[enhanced,colback=popink,colframe=popink,boxrule=2.2pt,arc=10pt,
      drop shadow=poporange,left=14pt,right=14pt,top=10pt,bottom=10pt,before skip=0pt,after skip=0pt]
    \tikz[baseline=-0.7ex]\node[circle,fill=popgreen,draw=popcream,line width=1.3pt,minimum size=22pt,inner sep=0]{\popcheck{white}};%
    \hspace{9pt}\begin{minipage}[c]{0.62\linewidth}
      {\popxbold\fontsize{12}{13}\selectfont\color{popcream}Signed\ \textbullet\ {\poplogo OnBuch\color{poporange}+}}\par\vspace{2pt}
      {\raggedright\popsemi\fontsize{8}{10.5}\selectfont\color{popline}\MakeUppercase{OnBuch+ teaching team}\\\MakeUppercase{Follows the official MINESEC syllabus}\par}
    \end{minipage}\hfill
    {\poplogo\fontsize{22}{22}\selectfont\color{popcream}OnBuch\color{poporange}+}
  \end{tcolorbox}%
}

% ============================================================
% Page de garde
% #1 titre · #2 matière · #3 niveau · #4 module · #5 n° leçon · #6 durée ·
% #7 description / objectifs (texte ou itemize)
% ============================================================
\newcommand{\pagegarde}[7]{%
  \thispagestyle{empty}
  \newgeometry{a4paper,margin=1.7cm,top=1.6cm,bottom=1.4cm}%
  \begin{tikzpicture}[remember picture,overlay]
    \fill[poporange!18] ([xshift=-3.2cm,yshift=-3.6cm]current page.north east) circle (4.6cm);
    \fill[poppurple!14] ([xshift=2.4cm,yshift=7.6cm]current page.south west) circle (3.8cm);
  \end{tikzpicture}%
  {\poplogo\fontsize{30}{30}\selectfont\color{popink}OnBuch\color{poporange}+}\hfill
  \raisebox{0.6ex}{\poppill{Signed course \textbullet\ Premium}{popink}{popcream}}\par
  \vspace{1pt}\popsquiggle{poppurple}\par
  \vspace{8pt}
  \begin{tcolorbox}[enhanced,colback=poporange,colframe=popink,boxrule=2.8pt,arc=13pt,
      drop shadow=popink,left=18pt,right=18pt,top=12pt,bottom=13pt,after skip=10pt]
    \poppill{#2 \textbullet\ #3}{popink}{popcream}\hspace{5pt}\poppill{#4}{white}{popink}\par\vspace{8pt}
    {\popdisplay\fontsize{14}{14}\selectfont\color{popink}CHAPTER #5}\par\vspace{4pt}
    {\raggedright\hyphenpenalty=10000\exhyphenpenalty=10000\popdisplay\fontsize{25}{27}\selectfont\color{popcream}\MakeUppercase{#1}\par}
    \vspace{8pt}
    {\popxbold\fontsize{9.5}{9.5}\selectfont\color{popink}\MakeUppercase{Suggested time: #6}}
  \end{tcolorbox}
  \begin{tcolorbox}[enhanced,colback=white,colframe=popink,boxrule=2.2pt,arc=10pt,
      drop shadow=popink,left=16pt,right=16pt,top=10pt,bottom=10pt,after skip=8pt]
    \poppill{In this chapter}{poppurple}{white}\par\vspace{5pt}
    \popsemi\fontsize{9.3}{12.6}\selectfont\color{popink}#7
  \end{tcolorbox}
  \noindent\begin{minipage}[t]{0.487\linewidth}
    \begin{tcolorbox}[enhanced,colback=popgreen,colframe=popink,boxrule=2.2pt,arc=10pt,
        drop shadow=popink,left=11pt,right=11pt,top=10pt,bottom=10pt,height=3.1cm,valign=center]
      \tcbox[colback=white,colframe=popink,boxrule=1.3pt,arc=5pt,boxsep=0pt,nobeforeafter,
        left=3pt,right=3pt,top=3pt,bottom=3pt,valign=center]{\qrcode[height=1.9cm]{https://play.google.com/store/apps/details?id=com.luvvix.onbuch&hl=en}}%
      \hspace{8pt}\begin{minipage}[c]{0.5\linewidth}
        {\popdisplay\fontsize{11}{12.5}\selectfont\color{white}\MakeUppercase{Download OnBuch}}\par\vspace{4pt}
        {\raggedright\popsemi\fontsize{8.4}{11}\selectfont\color{white}Free \textbullet\ Google Play\\AI tutor, past papers, results\par}
      \end{minipage}
    \end{tcolorbox}
  \end{minipage}\hfill
  \begin{minipage}[t]{0.487\linewidth}
    \begin{tcolorbox}[enhanced,colback=poppurple,colframe=popink,boxrule=2.2pt,arc=10pt,
        drop shadow=popink,left=11pt,right=11pt,top=10pt,bottom=10pt,height=3.1cm,valign=center]
      \tikz[baseline=-0.7ex]\node[circle,fill=white,draw=popink,line width=1.3pt,minimum size=1.6cm,inner sep=0]{\popdisplay\fontsize{24}{24}\selectfont\color{poppurple}+};%
      \hspace{8pt}\begin{minipage}[c]{0.6\linewidth}
        {\popdisplay\fontsize{11}{12.5}\selectfont\color{white}\MakeUppercase{Go OnBuch+}}\par\vspace{4pt}
        {\raggedright\popsemi\fontsize{8.4}{11}\selectfont\color{white}All signed courses, unlimited AI tutor, PDF sheets — OnBuch+ tab in the app\par}
      \end{minipage}
    \end{tcolorbox}
  \end{minipage}
  \vspace{8pt}
  \popcontenu
  \vfill
  \signatureonbuch
  \restoregeometry
  \newpage
}

% Rangée « Ce cours contient » (page de garde)
\newcommand{\poptile}[3]{% #1 couleur #2 gros texte #3 légende
  \begin{tcolorbox}[enhanced,colback=#1,colframe=popink,boxrule=2pt,arc=9pt,shadow={2.5pt}{-2.5pt}{0pt}{popink},
      left=6pt,right=6pt,top=9pt,bottom=9pt,halign=center,valign=center,height=1.75cm,width=\linewidth,nobeforeafter]
    {\popdisplay\fontsize{12.5}{13}\selectfont\color{white}\MakeUppercase{#2}}\par\vspace{3pt}
    {\centering\popsemi\fontsize{7.8}{9.5}\selectfont\color{white}#3\par}
  \end{tcolorbox}}
\newcommand{\popcontenu}{%
  \par\noindent{\popxboldls\fontsize{7.5}{7.5}\selectfont\color{popmuted}THIS SIGNED COURSE CONTAINS}\par\vspace{7pt}
  \noindent\begin{minipage}[t]{0.235\linewidth}\poptile{popblue}{Course}{complete, illustrated, step by step}\end{minipage}\hfill
  \begin{minipage}[t]{0.235\linewidth}\poptile{poporange}{Methods}{and worked examples}\end{minipage}\hfill
  \begin{minipage}[t]{0.235\linewidth}\poptile{poppink}{Exercises}{exam style + solutions}\end{minipage}\hfill
  \begin{minipage}[t]{0.235\linewidth}\poptile{popgreen}{Summary sheet}{the essentials to remember}\end{minipage}\par}

% Sommaire
\newtcolorbox{sommaire}{enhanced,colback=white,colframe=popink,boxrule=2.2pt,arc=10pt,
  drop shadow=popink,left=18pt,right=18pt,top=13pt,bottom=13pt,before skip=6pt,after skip=20pt,
  before upper={\poppill{Contents}{poppurple}{white}\par\vspace{9pt}\popsemi\fontsize{10.6}{14.5}\selectfont\color{popink}}}

% Signature de fin de cours — poussée tout en bas de la dernière page
\newcommand{\findecours}{%
  \par\vfill\nopagebreak
  \begin{center}\popsquiggle{poporange}\end{center}\vspace{4pt}
  \signatureonbuch
}
\AtBeginDocument{\def\llbracket{\mathopen{[\mkern-3.5mu[}}\def\rrbracket{\mathclose{]\mkern-3.5mu]}}} % intervalles d'entiers (glyphe ⟦ absent de la police)
% Glyphes absents de Fira Math : redéfinis à partir de glyphes disponibles
\AtBeginDocument{%
  \def\setminus{\mathbin{\backslash}}\let\smallsetminus\setminus
  \def\vdots{\mathord{\vbox{\baselineskip3.2pt\lineskiplimit0pt\kern4pt\hbox{.}\hbox{.}\hbox{.}}}}%
  \def\ddots{\mathinner{\mkern1mu\raise6.5pt\hbox{.}\mkern2mu\raise3.6pt\hbox{.}\mkern2mu\raise0.7pt\hbox{.}\mkern1mu}}%
  \def\triangle{\mathord{\scalebox{0.9}{$\symup{\Delta}$}}}%
  \def\nabla{\mathord{\raisebox{\depth}{\scalebox{1}[-1]{$\symup{\Delta}$}}}}%
  \def\checkmark{\popcheck{popdark}}%
  \def\star{\mathbin{\ast}}\def\bigstar{\mathord{\ast}}%
  \def\diamond{\mathbin{\scalebox{0.6}{$\Diamond$}}}%
  \def\bigcup{\mathop{\mathchoice{\raisebox{-0.3em}{\scalebox{1.7}{$\cup$}}}{\scalebox{1.25}{$\cup$}}{\cup}{\cup}}\displaylimits}%
  \def\bigcap{\mathop{\mathchoice{\raisebox{-0.3em}{\scalebox{1.7}{$\cap$}}}{\scalebox{1.25}{$\cap$}}{\cap}{\cap}}\displaylimits}%
  \def\lesseqqgtr{\mathrel{\lessgtr}}\def\bigoplus{\mathop{\oplus}\displaylimits}%
}
\providecommand{\ssi}{if and only if} % abréviation fréquente
\AtBeginDocument{\providecommand{\wideparen}{\overparen}} % arc de cercle

%% FIGFIX-BEGIN v4 — correctifs globaux des figures (docs/onbuch-generation/tools/figfix)
\usepackage{adjustbox}\usepackage{etoolbox}
% toute figure TikZ/pgfplots plus large que la ligne est réduite à la largeur disponible
\BeforeBeginEnvironment{tikzpicture}{\begin{adjustbox}{max width=\linewidth}}
\AfterEndEnvironment{tikzpicture}{\end{adjustbox}}
% légendes pgfplots lisibles par-dessus les courbes
\pgfplotsset{every axis/.append style={legend style={fill=white,fill opacity=0.92,text opacity=1,draw=black!15,inner sep=3pt}}}
% étiquettes d'arêtes (syntaxe edge["..."]) avec fond blanc
\tikzset{every edge quotes/.append style={fill=white,inner sep=1.2pt}}
%% FIGFIX-END
\begin{document}

\pagegarde{Complex Numbers}{Pure Mathematics with Mechanics}{Upper Sixth}{Upper Sixth — Pure mathematics with mechanics}{2}{19 periods}{This chapter introduces the complex number system, from the imaginary unit i to the powerful polar and exponential forms. Students will master algebraic operations, the Argand diagram, De Moivre's theorem, roots of complex numbers and geometric loci — essential tools for the GCE Advanced Level and beyond.
\vspace{4pt}
\begin{multicols}{2}\small
\begin{itemize}
\item By the end of this chapter I can define the imaginary unit i and compute its powers.
\item By the end of this chapter I can add, subtract, multiply and divide complex numbers in Cartesian form.
\item By the end of this chapter I can use complex conjugates and their properties to simplify expressions and solve equations.
\item By the end of this chapter I can represent complex numbers on an Argand diagram and interpret modulus and argument geometrically.
\item By the end of this chapter I can convert a complex number between Cartesian, polar (trigonometric) and exponential forms.
\item By the end of this chapter I can apply De Moivre's theorem to compute powers and to relate multiple angles to powers of cos x and sin x.
\end{itemize}\end{multicols}}

\begin{sommaire}
\begin{multicols}{2}
\begin{enumerate}
\item Birth of Complex Numbers and the Imaginary Unit i
\item Algebra of Complex Numbers and Conjugates
\item Equations and Roots in ℂ
\item The Argand Diagram, Modulus and Argument
\item Polar and Exponential Forms
\item De Moivre's Theorem and Trigonometric Applications
\item nth Roots of Complex Numbers and Cube Roots of Unity
\item Loci in the Complex Plane and Chapter Synthesis
\item Integration activity
\item Methods and worked examples
\item Exercises
\item Solutions to the exercises
\item Summary sheet
\end{enumerate}
\end{multicols}
\end{sommaire}

\begin{objectifs}
\begin{itemize}
\item By the end of this chapter I can define the imaginary unit i and compute its powers.
\item By the end of this chapter I can add, subtract, multiply and divide complex numbers in Cartesian form.
\item By the end of this chapter I can use complex conjugates and their properties to simplify expressions and solve equations.
\item By the end of this chapter I can represent complex numbers on an Argand diagram and interpret modulus and argument geometrically.
\item By the end of this chapter I can convert a complex number between Cartesian, polar (trigonometric) and exponential forms.
\item By the end of this chapter I can apply De Moivre's theorem to compute powers and to relate multiple angles to powers of cos x and sin x.
\item By the end of this chapter I can find square roots and nth roots of complex numbers, including cube roots of unity.
\item By the end of this chapter I can solve polynomial equations with complex roots.
\item By the end of this chapter I can describe and sketch loci defined by conditions on complex numbers.
\end{itemize}
\end{objectifs}

\begin{prerequis}
\begin{itemize}
\item Real number algebra: expanding, factorising and solving quadratic equations (Lower Sixth, PMM01)
\item Trigonometric ratios, standard angles and identities (cos²x + sin²x = 1, addition formulae)
\item Surds and rationalising denominators
\item Coordinates and equations of lines and circles in the Cartesian plane
\item Radians as a measure of angles
\end{itemize}
\end{prerequis}

\newpage

\coursec{Birth of Complex Numbers and the Imaginary Unit i}

In 2019, engineers restoring electricity to parts of Bamenda after a severe storm had to solve circuit equations whose solutions involved the square root of a negative number --- something ``impossible'' in ordinary arithmetic. Yet by accepting this ``impossible'' number, they could compute alternating currents, predict phase shifts and keep the lights on. From the design of the M\'emve'\'ele hydroelectric dam to the signal processing inside every mobile phone in Douala, complex numbers are the invisible engine of modern technology. This chapter opens the door to that hidden number system: you will learn to compute with it, picture it, and use it to solve equations that have no solution in the real world.

The central problem is simple to state. The equation
\[ x^{2} + 1 = 0 \quad \Longleftrightarrow \quad x^{2} = -1 \]
has \emph{no} solution in $\mathbb{R}$, because the square of any real number is non-negative. For centuries, mathematicians simply declared such equations ``impossible''. In this section we trace how necessity forced them to change their minds, we define the \cle{imaginary unit} $i$, we learn to handle square roots of negative numbers and powers of $i$, and we construct the set $\mathbb{C}$ of \cle{complex numbers}.

\courssub{Historical Motivation: Cardano, Bombelli and the Cubic Equation}

A common misconception is that complex numbers were invented to solve the quadratic $x^{2}+1=0$. In fact, mathematicians of the sixteenth century were perfectly happy to say that such a quadratic has no solution and to move on. The real pressure came from a place where nobody expected it: \emph{cubic equations that do have real solutions}.

\textbf{Cardano's formula.} In 1545, the Italian mathematician Girolamo Cardano published \emph{Ars Magna}, which contained a formula for solving cubic equations. Applied to the depressed cubic $x^{3} = px + q$ (where the quadratic term has been eliminated), his formula gives
\[ x = \sqrt[3]{\frac{q}{2} + \sqrt{\frac{q^{2}}{4} - \frac{p^{3}}{27}}} \;+\; \sqrt[3]{\frac{q}{2} - \sqrt{\frac{q^{2}}{4} - \frac{p^{3}}{27}}}. \]

\textbf{Bombelli's puzzle.} Around 1572, Rafael Bombelli applied Cardano's formula to the equation
\[ x^{3} = 15x + 4. \]
You can check by substitution that $x = 4$ is a solution: $4^{3} = 64$ and $15(4) + 4 = 64$. Yet Cardano's formula produces
\[ x = \sqrt[3]{2 + \sqrt{-121}} + \sqrt[3]{2 - \sqrt{-121}} = \sqrt[3]{2 + 11\sqrt{-1}} + \sqrt[3]{2 - 11\sqrt{-1}}. \]
Here was the scandal: a \emph{perfectly real} answer ($x = 4$) could only be reached by a formula passing through the ``impossible'' quantity $\sqrt{-1}$. Bombelli had the courage to take $\sqrt{-1}$ seriously. He laid down rules for computing with it --- essentially the rules we still use --- and showed that
\[ \sqrt[3]{2 + 11\sqrt{-1}} = 2 + \sqrt{-1}, \qquad \sqrt[3]{2 - 11\sqrt{-1}} = 2 - \sqrt{-1}, \]
so that the formula gives $(2 + \sqrt{-1}) + (2 - \sqrt{-1}) = 4$, exactly the known real solution. The ``impossible'' numbers cancelled, leaving a real answer.

\begin{savaistu}[From ``sophistic'' to respectable]
For nearly three centuries, numbers involving $\sqrt{-1}$ were called \emph{imaginary}, a slightly mocking name coined by Descartes in 1637. It was only after Gauss and Argand gave them a concrete geometric picture (the plane we shall meet in Section~4) that complex numbers became fully respectable. Today the word ``imaginary'' survives, but there is nothing imaginary about their applications: electrical engineering, fluid mechanics, quantum physics and telecommunications all run on complex arithmetic.
\end{savaistu}

The lesson of history is the lesson of this chapter: when a new kind of number proves \emph{useful and consistent}, we accept it. We now do this formally.

\courssub{The Imaginary Unit $i$}

Since no real number squares to $-1$, we simply \emph{create} a new number with that property and study the arithmetic it generates.

\begin{definition}[The imaginary unit]
The \cle{imaginary unit}, denoted $i$, is a new (non-real) number defined by the single property
\[ i^{2} = -1, \qquad \text{and we write } \qquad i = \sqrt{-1}. \]
It is postulated to combine with real numbers under the ordinary rules of arithmetic (commutativity, associativity, distributivity).
\end{definition}

\begin{attention}[The symbol $\sqrt{-1}$ is a definition, not a discovery]
You cannot ``find'' $i$ on the real number line, and you must not apply to it the rules of real square roots without care (see the warning on products below). The equation $i = \sqrt{-1}$ is a \emph{notation} summarising the definition $i^{2} = -1$.
\end{attention}

\textbf{Square roots of negative numbers.} Once $i$ is accepted, every negative number acquires a square root.

\begin{propriete}[Square root of a negative number]
Let $a > 0$ be a positive real number. Then
\[ \sqrt{-a} = i\sqrt{a}, \qquad \text{because} \qquad \left(i\sqrt{a}\right)^{2} = i^{2}\left(\sqrt{a}\right)^{2} = (-1)(a) = -a. \]
\end{propriete}

\begin{exemplebox}[Square roots of negative numbers]
\begin{itemize}
\item $\sqrt{-9} = i\sqrt{9} = 3i$;
\item $\sqrt{-5} = i\sqrt{5}$;
\item $\sqrt{-12} = i\sqrt{12} = 2i\sqrt{3}$.
\end{itemize}
Each answer can be checked by squaring: $(3i)^{2} = 9i^{2} = -9$. \checkmark
\end{exemplebox}

\begin{attention}[The most dangerous trap in the chapter]
The familiar rule $\sqrt{a}\,\sqrt{b} = \sqrt{ab}$ is \textbf{only valid for non-negative real numbers}. It fails spectacularly for negative numbers:
\[ \sqrt{-4}\,\sqrt{-9} \;\neq\; \sqrt{(-4)(-9)} = \sqrt{36} = 6. \]
The correct computation is to \emph{first} rewrite each root using $i$:
\[ \sqrt{-4}\,\sqrt{-9} = (2i)(3i) = 6i^{2} = -6. \]
\textbf{Golden rule:} whenever you see $\sqrt{-a}$ with $a>0$, immediately rewrite it as $i\sqrt{a}$ \emph{before} doing any other operation.
\end{attention}

\courssub{Powers of $i$: a Cycle of Period 4}

What are $i^{3}$, $i^{4}$, $i^{5}$, \dots? Using only $i^{2} = -1$ and the laws of indices, we compute:
\begin{align*}
i^{1} &= i, &
i^{2} &= -1, &
i^{3} &= i^{2}\cdot i = -i, &
i^{4} &= i^{2}\cdot i^{2} = (-1)(-1) = 1, &
i^{5} &= i^{4}\cdot i = i.
\end{align*}
At $i^{5}$ we are back to $i$: the values repeat. The powers of $i$ run endlessly around the cycle
\[ i \;\longrightarrow\; -1 \;\longrightarrow\; -i \;\longrightarrow\; 1 \;\longrightarrow\; i \;\longrightarrow\; \cdots \]

\begin{popfigure}
\begin{center}
\begin{tikzpicture}[scale=2.1, >=stealth]
\coordinate (A) at (0,1);
\coordinate (B) at (-1,0);
\coordinate (C) at (0,-1);
\coordinate (D) at (1,0);
\draw[popline, dashed] (0,0) circle (1);
\filldraw[popblue] (A) circle (1.6pt) node[above=2pt, popdark] {\textbf{$i^{1}=i$}};
\filldraw[poporange] (B) circle (1.6pt) node[left=2pt, popdark] {\textbf{$i^{2}=-1$}};
\filldraw[poppurple] (C) circle (1.6pt) node[below=2pt, popdark] {\textbf{$i^{3}=-i$}};
\filldraw[popgreen] (D) circle (1.6pt) node[right=2pt, popdark] {\textbf{$i^{4}=1$}};
\draw[->, popblue, thick, bend left=25] (D) to (A);
\draw[->, popblue, thick, bend left=25] (A) to (B);
\draw[->, popblue, thick, bend left=25] (B) to (C);
\draw[->, popblue, thick, bend left=25] (C) to (D);
\node[popmuted] at (0,0) {$\times\, i$};
\end{tikzpicture}
\end{center}
\legende{The cycle of powers of $i$: each multiplication by $i$ moves one step around the cycle $1 \to i \to -1 \to -i \to 1$.}
\end{popfigure}

\begin{propriete}[Periodicity of the powers of $i$]
For every integer $n$,
\[ i^{n+4} = i^{n}. \]
In other words, the sequence $(i^{n})$ is \cle{periodic with period 4}. Consequently, if $n = 4q + r$ with $r \in \{0,1,2,3\}$ (Euclidean division of $n$ by $4$), then
\[ i^{n} = i^{r}, \qquad \text{so } i^{n} \in \{1,\ i,\ -1,\ -i\}. \]
\end{propriete}

\textbf{Proof (examinable, and instructive).} By the Euclidean division of $n$ by $4$, there exist integers $q$ and $r$ with $n = 4q + r$ and $0 \le r \le 3$. Then, using the laws of indices and $i^{4} = 1$:
\[ i^{n} = i^{4q+r} = \left(i^{4}\right)^{q} \cdot i^{r} = 1^{q} \cdot i^{r} = i^{r}. \qquad \blacksquare \]

\begin{methode}[Reducing $i^{n}$ for any integer $n$]
\begin{enumerate}
\item Divide $n$ by $4$: write $n = 4q + r$ with remainder $r \in \{0,1,2,3\}$.
\item Then $i^{n} = i^{r}$.
\item Read off the value: $i^{0} = 1$, $i^{1} = i$, $i^{2} = -1$, $i^{3} = -i$.
\item For negative exponents, use $i^{-1} = \dfrac{1}{i} = \dfrac{1 \cdot i}{i \cdot i} = \dfrac{i}{-1} = -i$, or add a multiple of $4$ to the exponent to make it positive.
\end{enumerate}
\end{methode}

\begin{exemplebox}[Computing powers of $i$]
\begin{itemize}
\item $i^{2023}$: since $2023 = 4 \times 505 + 3$, we get $i^{2023} = i^{3} = -i$.
\item $i^{100}$: since $100 = 4 \times 25 + 0$, we get $i^{100} = i^{0} = 1$.
\item $i^{-5}$: since $-5 = 4 \times (-2) + 3$, we get $i^{-5} = i^{3} = -i$. (Check: $i^{-5} = i^{-4}\cdot i^{-1} = 1 \cdot (-i) = -i$. \checkmark)
\end{itemize}
\end{exemplebox}

\courssub{Complex Numbers: Definition, Real and Imaginary Parts}

Numbers such as $2 + 11i$ (which Bombelli met inside the cube roots) are built from two real numbers and the unit $i$. This motivates the central definition of the chapter.

\begin{definition}[Complex number, real part, imaginary part]
A \cle{complex number} is any expression of the form
\[ z = a + ib, \qquad \text{where } a, b \in \mathbb{R}. \]
\begin{itemize}
\item $a$ is called the \cle{real part} of $z$, written $\operatorname{Re}(z) = a$;
\item $b$ is called the \cle{imaginary part} of $z$, written $\operatorname{Im}(z) = b$ (note: $\operatorname{Im}(z)$ is a \emph{real} number, \emph{not} $ib$).
\end{itemize}
The form $a + ib$ is called the \cle{algebraic form} (or Cartesian form) of $z$. The set of all complex numbers is denoted $\mathbb{C}$:
\[ \mathbb{C} = \{\, a + ib \;:\; a, b \in \mathbb{R} \,\}. \]
\end{definition}

\begin{attention}[$\operatorname{Im}(z)$ is real]
If $z = 3 - 5i$, then $\operatorname{Re}(z) = 3$ and $\operatorname{Im}(z) = -5$, \textbf{not} $-5i$. Examiners penalise this slip every year.
\end{attention}

\textbf{Special cases.} Two familiar families sit inside $\mathbb{C}$:
\begin{itemize}
\item If $b = 0$, then $z = a$ is an ordinary real number. Thus $\mathbb{R} \subset \mathbb{C}$: every real number is a complex number with zero imaginary part.
\item If $a = 0$ and $b \neq 0$, then $z = ib$ is called \cle{purely imaginary} (e.g.\ $3i$, $-\tfrac{1}{2}i$).
\item If $a = b = 0$, then $z = 0$: the number $0$ is the only complex number that is both real and purely imaginary.
\end{itemize}

Each enlargement of the number system in your school career was driven by equations that had no solution: $\mathbb{N}$ cannot solve $x + 5 = 3$ (so we built $\mathbb{Z}$); $\mathbb{Z}$ cannot solve $3x = 2$ (so we built $\mathbb{Q}$); $\mathbb{Q}$ cannot solve $x^{2} = 2$ (so we built $\mathbb{R}$); and $\mathbb{R}$ cannot solve $x^{2} = -1$ (so we build $\mathbb{C}$). The chain is:
\[ \mathbb{N} \subset \mathbb{Z} \subset \mathbb{Q} \subset \mathbb{R} \subset \mathbb{C}. \]

\begin{popfigure}
\begin{center}
\begin{tikzpicture}[scale=1.0]
\draw[fill=popblueL, draw=popblue, thick, rounded corners=8pt] (0,0) ellipse (5.4 and 3.1);
\draw[fill=white, draw=popgreen, thick, rounded corners=6pt] (-0.4,0.1) ellipse (3.9 and 2.3);
\draw[fill=popblueL, draw=poporange, thick, rounded corners=6pt] (-0.8,0.2) ellipse (2.7 and 1.6);
\draw[fill=white, draw=poppurple, thick, rounded corners=6pt] (-1.2,0.3) ellipse (1.7 and 1.0);
\draw[fill=popblueL, draw=popdark, thick, rounded corners=6pt] (-1.5,0.35) ellipse (0.95 and 0.6);
\node[popdark] at (-1.5,0.35) {$\mathbb{N}$};
\node[poppurple] at (0.15,0.9) {$\mathbb{Z}$};
\node[poporange] at (1.35,1.35) {$\mathbb{Q}$};
\node[popgreen] at (2.9,1.8) {$\mathbb{R}$};
\node[popblue] at (3.9,2.55) {$\mathbb{C}$};
\node[popmuted, align=center, text width=2.6cm] at (3.6,-1.6) {\scriptsize $i,\ 2+3i,\ -\tfrac12 i$};
\node[popmuted, align=center, text width=2.2cm] at (2.5,-0.9) {\scriptsize $\sqrt{2},\ \pi$};
\node[popmuted] at (0.5,-0.35) {\scriptsize $-3$};
\node[popmuted] at (-0.7,0.55) {\scriptsize $\tfrac{2}{3}$};
\node[popmuted] at (-1.5,-0.35) {\scriptsize $0,1,7$};
\end{tikzpicture}
\end{center}
\legende{The nested number sets $\mathbb{N} \subset \mathbb{Z} \subset \mathbb{Q} \subset \mathbb{R} \subset \mathbb{C}$, with a few typical members of each layer.}
\end{popfigure}

\courssub{Equality of Two Complex Numbers}

When are two complex numbers equal? Since a complex number is completely determined by the \emph{pair} of real numbers $(a, b)$, two complex numbers are equal exactly when both components agree.

\begin{propriete}[Equality in $\mathbb{C}$]
Let $z = a + ib$ and $z' = a' + ib'$ with $a, b, a', b' \in \mathbb{R}$. Then
\[ z = z' \quad \Longleftrightarrow \quad \operatorname{Re}(z) = \operatorname{Re}(z') \ \text{ and }\ \operatorname{Im}(z) = \operatorname{Im}(z'), \]
that is,
\[ a + ib = a' + ib' \quad \Longleftrightarrow \quad a = a' \ \text{ and }\ b = b'. \]
In particular, $a + ib = 0 \iff a = 0$ and $b = 0$.
\end{propriete}

\textbf{Justification.} Suppose $a + ib = a' + ib'$. Rearranging, $(a - a') = i(b' - b)$. If $b' - b \neq 0$, we could divide and obtain $i = \dfrac{a - a'}{b' - b}$, a \emph{real} number --- impossible, since $i$ is not real. Hence $b' - b = 0$, i.e.\ $b = b'$, and then $a - a' = 0$, i.e.\ $a = a'$. The converse is immediate. $\blacksquare$

This property is a working tool: one complex equation is worth \emph{two} real equations, one for the real parts and one for the imaginary parts.

\begin{exoresolu}[Finding unknown reals from a complex equality]
Find the real numbers $x$ and $y$ such that
\[ (2x - 1) + i(y + 3) = 5 + 2i. \]
\tcblower
\textbf{Solution.} The two sides are complex numbers in algebraic form. By the equality property, we equate the real parts and the imaginary parts separately:
\begin{align*}
\text{Real parts:} \quad 2x - 1 &= 5 & \text{Imaginary parts:} \quad y + 3 &= 2 \\
2x &= 6 & y &= -1 \\
x &= 3. &&
\end{align*}
\textbf{Check:} $(2(3)-1) + i(-1+3) = 5 + 2i$. \checkmark

Hence $\boxed{x = 3 \text{ and } y = -1}$.
\end{exoresolu}

\begin{exoresolu}[Using $i^{2} = -1$ to simplify a product]
Write $(2 + 3i)(1 - i)$ in the form $a + ib$, stating $\operatorname{Re}(z)$ and $\operatorname{Im}(z)$.
\tcblower
\textbf{Solution.} Expand using the ordinary rules of arithmetic, then replace $i^{2}$ by $-1$:
\begin{align*}
(2 + 3i)(1 - i) &= 2(1) + 2(-i) + 3i(1) + 3i(-i) \\
&= 2 - 2i + 3i - 3i^{2} \\
&= 2 + i - 3(-1) \\
&= 5 + i.
\end{align*}
Thus $\operatorname{Re}(z) = 5$ and $\operatorname{Im}(z) = 1$ (the imaginary part is $1$, not $i$).
\end{exoresolu}

\begin{attention}[Summary of traps in this section]
\begin{itemize}
\item Never write $\sqrt{-a}\sqrt{-b} = \sqrt{ab}$ for $a, b > 0$; rewrite as $i\sqrt{a}$ and $i\sqrt{b}$ first.
\item $\operatorname{Im}(a + ib) = b$, a real number --- never $ib$.
\item $i^{2} = -1$ by \emph{definition}; do not ``prove'' it by writing $\sqrt{-1}\sqrt{-1} = \sqrt{(-1)(-1)} = 1$, which misuses the real square-root rule and gives the wrong sign.
\item A complex equality gives \emph{two} conditions; forgetting one of them loses marks.
\end{itemize}
\end{attention}

\begin{exercice}[Consolidation]
\begin{enumerate}
\item Write in the form $i\sqrt{a}$ with $a$ a positive integer as small as possible: $\sqrt{-49}$, $\sqrt{-18}$, $\sqrt{-75}$.
\item Compute $i^{37}$, $i^{2024}$, $i^{-11}$ and $i^{4n+3}$ for $n \in \mathbb{Z}$.
\item Compute $\sqrt{-8}\,\sqrt{-2}$ in two ways (the wrong way and the right way) and explain the discrepancy.
\item Find real numbers $x$ and $y$ such that $(x + y) + i(x - 2y) = 7 + i$.
\item Expand and simplify $(1 + i)^{2}$, $(1 - i)^{2}$ and $(1 + i)(1 - i)$. What do you notice about the last result?
\end{enumerate}
\end{exercice}

\begin{corrige}[Exercise 1]
\textbf{1.} $\sqrt{-49} = 7i$; $\sqrt{-18} = i\sqrt{18} = 3i\sqrt{2}$; $\sqrt{-75} = i\sqrt{75} = 5i\sqrt{3}$.

\textbf{2.} $37 = 4(9)+1 \Rightarrow i^{37} = i$; $2024 = 4(506)+0 \Rightarrow i^{2024} = 1$; $-11 = 4(-3)+1 \Rightarrow i^{-11} = i$; $i^{4n+3} = i^{3} = -i$.

\textbf{3.} Wrong way: $\sqrt{(-8)(-2)} = \sqrt{16} = 4$. Right way: $(i\sqrt{8})(i\sqrt{2}) = i^{2}\sqrt{16} = -4$. The rule $\sqrt{a}\sqrt{b}=\sqrt{ab}$ does not apply to negative numbers.

\textbf{4.} Equating parts: $x + y = 7$ and $x - 2y = 1$. Subtracting gives $3y = 6$, so $y = 2$ and $x = 5$.

\textbf{5.} $(1+i)^{2} = 1 + 2i + i^{2} = 2i$; $(1-i)^{2} = -2i$; $(1+i)(1-i) = 1 - i^{2} = 2$, a \emph{real} number --- a first glimpse of the power of conjugates, studied in the next section.
\end{corrige}

\begin{aretenir}[Key points of Section 1]
\begin{itemize}
\item Complex numbers were forced on mathematicians by cubic equations (Cardano, Bombelli): real answers sometimes require a detour through $\sqrt{-1}$.
\item The imaginary unit $i$ is defined by $i^{2} = -1$; for $a > 0$, $\sqrt{-a} = i\sqrt{a}$ --- always rewrite in this form \emph{before} computing.
\item The powers of $i$ cycle with period $4$: $i, -1, -i, 1$; to compute $i^{n}$, reduce $n$ modulo $4$.
\item A complex number is $z = a + ib$ with $a = \operatorname{Re}(z)$, $b = \operatorname{Im}(z) \in \mathbb{R}$; special cases: real ($b=0$), purely imaginary ($a=0$, $b\neq 0$).
\item $\mathbb{N} \subset \mathbb{Z} \subset \mathbb{Q} \subset \mathbb{R} \subset \mathbb{C}$.
\item $a + ib = a' + ib' \iff a = a'$ and $b = b'$: one complex equation = two real equations.
\end{itemize}
\end{aretenir}

\coursec{Algebra of Complex Numbers and Conjugates}

In the previous section we met the imaginary unit $i$, defined by the single remarkable property $i^{2}=-1$, and we agreed that a complex number is any number of the form $z=a+ib$ with $a,b\in\mathbb{R}$. But a new kind of number is only useful if we can \emph{calculate} with it. Can we add two complex numbers? Multiply them? Divide them? The good news is that the answer is yes, and the rules are exactly the rules you already know from ordinary algebra in $\mathbb{R}$, with one extra ingredient: every time $i^{2}$ appears, we replace it by $-1$. This section makes those rules precise and introduces the \cle{complex conjugate}, the tool that makes division possible and that will accompany us throughout the whole chapter.

\courssub{Addition and Subtraction of Complex Numbers}

\textbf{Intuition.} Think of a complex number $z=a+ib$ as a ``two-part'' number: a real part $a$ and an imaginary part $b$. Just as we add vectors component by component, it is natural to add the real parts together and the imaginary parts together.

\begin{definition}[Sum and difference of complex numbers]
Let $z_{1}=a+ib$ and $z_{2}=c+id$, where $a,b,c,d\in\mathbb{R}$. Then
\[
z_{1}+z_{2}=(a+c)+i(b+d), \qquad z_{1}-z_{2}=(a-c)+i(b-d).
\]
The sum and the difference of two complex numbers are again complex numbers: we say that $\mathbb{C}$ is \emph{closed} under addition and subtraction.
\end{definition}

\begin{exemplebox}[Adding and subtracting]
Let $z_{1}=3+4i$ and $z_{2}=1-2i$. Then
\[
z_{1}+z_{2}=(3+1)+i(4-2)=4+2i, \qquad z_{1}-z_{2}=(3-1)+i(4-(-2))=2+6i.
\]
\end{exemplebox}

Addition of complex numbers inherits all the familiar properties of addition in $\mathbb{R}$: it is commutative ($z_{1}+z_{2}=z_{2}+z_{1}$), associative, has identity element $0=0+0i$, and every $z=a+ib$ has an opposite $-z=-a-ib$.

\courssub{Multiplication of Complex Numbers}

\textbf{Intuition.} To multiply $(a+ib)(c+id)$ we simply expand the brackets exactly as we would expand $(a+bx)(c+dx)$ in ordinary algebra, and then use the defining relation $i^{2}=-1$.

\begin{propriete}[Product of two complex numbers]
For $z_{1}=a+ib$ and $z_{2}=c+id$ with $a,b,c,d\in\mathbb{R}$:
\[
z_{1}z_{2}=(ac-bd)+i(ad+bc).
\]
\end{propriete}

\textbf{Derivation (instructive, worth knowing for the GCE).} Expand term by term:
\begin{align*}
(a+ib)(c+id) &= ac + iad + ibc + i^{2}bd \\
&= ac + iad + ibc - bd \qquad (\text{since } i^{2}=-1)\\
&= (ac-bd) + i(ad+bc).
\end{align*}
Notice how the term $i^{2}bd$ became $-bd$: this is where the real part $ac-bd$ comes from.

\begin{attention}[The sign trap: $i^{2}=-1$]
The most common error at this level is to write $(a+ib)(c+id)=ac+bd+i(ad+bc)$, forgetting that $i^{2}bd=-bd$. The product of the two imaginary parts is a \emph{real} number with a \emph{minus} sign. For example, $(2+3i)(1+i)=2+2i+3i+3i^{2}=2+5i-3=-1+5i$, and \textbf{not} $5+5i$. Always replace $i^{2}$ by $-1$ immediately, before collecting terms.
\end{attention}

\begin{exemplebox}[Multiplying complex numbers]
Compute $(2-3i)(4+5i)$.
\begin{align*}
(2-3i)(4+5i) &= 8 + 10i - 12i - 15i^{2} = 8 - 2i + 15 = 23 - 2i.
\end{align*}
\end{exemplebox}

Multiplication is commutative, associative, distributive over addition, and has identity $1=1+0i$. In particular, powers of $z$ are computed by repeated multiplication, always reducing with the cycle of powers of $i$:
\[
i^{2}=-1, \qquad i^{3}=-i, \qquad i^{4}=1, \qquad i^{5}=i, \ \ldots
\]
so that $i^{n}$ depends only on the remainder of $n$ modulo $4$.

\begin{exemplebox}[Using the powers of $i$]
Simplify $i^{2025}$. Since $2025=4\times 506+1$, we have $i^{2025}=(i^{4})^{506}\times i=1^{506}\times i=i$.
\end{exemplebox}

\courssub{The Complex Conjugate}

\textbf{Intuition.} Given $z=a+ib$, what happens if we simply flip the sign of the imaginary part? We obtain a new number $a-ib$, which turns out to be the perfect ``partner'' of $z$: when we add or multiply the two partners, the imaginary parts cancel or combine into something purely real. This partner is called the conjugate.

\begin{definition}[Complex conjugate]
The \cle{conjugate} of the complex number $z=a+ib$ (with $a,b\in\mathbb{R}$) is the complex number
\[
\bar{z}=a-ib.
\]
The real part is unchanged; the sign of the imaginary part is reversed.
\end{definition}

Geometrically, as we shall see in detail with the Argand diagram, $\bar{z}$ is the reflection of $z$ in the real axis. The figure below previews this mirror symmetry.

\begin{popfigure}
\begin{tikzpicture}[scale=1.1]
\draw[->,gray] (-0.8,0) -- (4.2,0) node[below left] {\small Re};
\draw[->,gray] (0,-2.6) -- (0,2.6) node[below left] {\small Im};
\draw[dashed,popmuted] (0,0) -- (3,1.8);
\draw[dashed,popmuted] (0,0) -- (3,-1.8);
\draw[dashed,popline] (3,1.8) -- (3,-1.8);
\fill[popblue] (3,1.8) circle (2pt) node[above right,popdark] {\small $z=a+ib$};
\fill[poporange] (3,-1.8) circle (2pt) node[below right,popdark] {\small $\bar{z}=a-ib$};
\node[popmuted,align=center] at (1.4,0.9) {\small mirror};
\node[popmuted,align=center] at (1.4,-0.9) {\small image};
\node[gray] at (3.35,0.25) {\small real axis};
\end{tikzpicture}
\legende{The conjugate $\bar{z}$ is the mirror image of $z$ across the real axis.}
\end{popfigure}

\begin{propriete}[Fundamental identities involving the conjugate]
Let $z=a+ib$ with $a,b\in\mathbb{R}$. Then:
\begin{enumerate}
\item $z+\bar{z}=2a=2\operatorname{Re}(z)$, a real number;
\item $z-\bar{z}=2ib=2i\operatorname{Im}(z)$, purely imaginary;
\item $z\bar{z}=a^{2}+b^{2}$, a \emph{non-negative real} number.
\end{enumerate}
Equivalently, $\operatorname{Re}(z)=\dfrac{z+\bar{z}}{2}$ and $\operatorname{Im}(z)=\dfrac{z-\bar{z}}{2i}$.
\end{propriete}

\textbf{Proof (examinable and instructive).} For (1): $z+\bar{z}=(a+ib)+(a-ib)=2a$. For (2): $z-\bar{z}=(a+ib)-(a-ib)=2ib$. For (3):
\[
z\bar{z}=(a+ib)(a-ib)=a^{2}-(ib)^{2}=a^{2}-i^{2}b^{2}=a^{2}+b^{2}\geq 0,
\]
where we used the difference of two squares and $i^{2}=-1$. $\blacksquare$

\begin{propriete}[$z\bar{z}=|z|^{2}$ is always a non-negative real]
Anticipating the modulus $|z|=\sqrt{a^{2}+b^{2}}$ of Section 4, identity (3) reads
\[
z\bar{z}=|z|^{2}=a^{2}+b^{2}\in\mathbb{R}^{+}.
\]
This is the single most useful fact about conjugates: \textbf{multiplying a complex number by its conjugate always produces a non-negative real number}. It is zero only when $z=0$. This is precisely what makes division by a complex number possible.
\end{propriete}

\begin{exemplebox}[Conjugate identities in action]
Let $z=5-2i$, so $\bar{z}=5+2i$. Then
\[
z+\bar{z}=10=2\operatorname{Re}(z), \qquad z-\bar{z}=-4i=2i\operatorname{Im}(z), \qquad z\bar{z}=25+4=29.
\]
Note that $29$ is real and positive, as promised.
\end{exemplebox}

\courssub{Division of Complex Numbers}

\textbf{The problem.} What does $\dfrac{2+3i}{1-4i}$ mean, and how can we write it in the standard form $x+iy$? The denominator $1-4i$ is not a real number, so we cannot ``divide'' directly. The trick is to use the conjugate to turn the denominator into a real number, exactly as we rationalise a surd denominator like $\dfrac{1}{1+\sqrt{2}}$ by multiplying by $1-\sqrt{2}$.

\begin{methode}[Dividing by a complex number]
To write $\dfrac{z_{1}}{z_{2}}$ (with $z_{2}\neq 0$) in the form $x+iy$:
\begin{enumerate}
\item Write down the conjugate $\bar{z}_{2}$ of the denominator.
\item Multiply \textbf{both} numerator and denominator by $\bar{z}_{2}$: $\dfrac{z_{1}}{z_{2}}=\dfrac{z_{1}\bar{z}_{2}}{z_{2}\bar{z}_{2}}$.
\item The denominator $z_{2}\bar{z}_{2}=|z_{2}|^{2}$ is now a positive real number.
\item Expand the numerator, reduce with $i^{2}=-1$, and split into real and imaginary parts.
\end{enumerate}
\end{methode}

\begin{exoresolu}[GCE-style simplification]
Express $\dfrac{2+3i}{1-4i}$ in the form $x+iy$, where $x,y\in\mathbb{R}$.
\tcblower
The conjugate of the denominator $1-4i$ is $1+4i$. Multiply top and bottom by $1+4i$:
\[
\frac{2+3i}{1-4i}=\frac{(2+3i)(1+4i)}{(1-4i)(1+4i)}.
\]
Numerator: $(2+3i)(1+4i)=2+8i+3i+12i^{2}=2+11i-12=-10+11i$.

Denominator: $(1-4i)(1+4i)=1^{2}+4^{2}=17$.

Therefore
\[
\frac{2+3i}{1-4i}=\frac{-10+11i}{17}=-\frac{10}{17}+\frac{11}{17}i.
\]
So $x=-\dfrac{10}{17}$ and $y=\dfrac{11}{17}$.
\end{exoresolu}

\begin{exoresolu}[A quotient with simplification]
Given $z=3+i$, express $\dfrac{1+7i}{z}$ in the form $x+iy$.
\tcblower
Here $\bar{z}=3-i$, and $z\bar{z}=9+1=10$.
\[
\frac{1+7i}{3+i}=\frac{(1+7i)(3-i)}{(3+i)(3-i)}=\frac{3-i+21i-7i^{2}}{10}=\frac{3+20i+7}{10}=\frac{10+20i}{10}=1+2i.
\]
\end{exoresolu}

\begin{attention}[Divide only by the conjugate of the \emph{denominator}]
Two frequent errors: (i) multiplying by the conjugate of the \emph{numerator} instead of the denominator; (ii) multiplying only the denominator, which changes the value of the fraction. Whatever you do to the bottom, you must do to the top. Also, never leave the final answer as $\dfrac{-10+11i}{17}$ when the question asks for the form $x+iy$: separate the real and imaginary parts.
\end{attention}

\courssub{Properties of Complex Conjugates}

The conjugation operation behaves beautifully with respect to all four operations. These properties save enormous time in GCE questions.

\begin{propriete}[Algebraic properties of the conjugate]
For all complex numbers $z_{1}$ and $z_{2}$ (with $z_{2}\neq 0$ where a quotient appears):
\begin{enumerate}
\item $\overline{z_{1}+z_{2}}=\bar{z}_{1}+\bar{z}_{2}$ (conjugate of a sum);
\item $\overline{z_{1}-z_{2}}=\bar{z}_{1}-\bar{z}_{2}$ (conjugate of a difference);
\item $\overline{z_{1}z_{2}}=\bar{z}_{1}\,\bar{z}_{2}$ (conjugate of a product);
\item $\overline{\left(\dfrac{z_{1}}{z_{2}}\right)}=\dfrac{\bar{z}_{1}}{\bar{z}_{2}}$ (conjugate of a quotient);
\item $\overline{z^{n}}=(\bar{z})^{n}$ for every $n\in\mathbb{N}$ (conjugate of a power);
\item $\overline{\bar{z}}=z$ (conjugating twice returns the original number);
\item $z=\bar{z}\iff z\in\mathbb{R}$ (a number equal to its conjugate is real);
\item $z=-\bar{z}\iff z$ is purely imaginary.
\end{enumerate}
In words: \emph{the conjugate of a sum, product or quotient is the sum, product or quotient of the conjugates}.
\end{propriete}

\textbf{Proof of (3), (7) and (8) (instructive).} Write $z_{1}=a+ib$, $z_{2}=c+id$. Then
\[
z_{1}z_{2}=(ac-bd)+i(ad+bc)\quad\Longrightarrow\quad \overline{z_{1}z_{2}}=(ac-bd)-i(ad+bc),
\]
while
\[
\bar{z}_{1}\bar{z}_{2}=(a-ib)(c-id)=(ac-bd)-i(ad+bc).
\]
The two expressions agree, proving (3). Property (5) follows from (3) by induction on $n$. For (7): $z=\bar{z}$ means $a+ib=a-ib$, i.e. $2ib=0$, i.e. $b=0$, so $z=a\in\mathbb{R}$. For (8): $z=-\bar{z}$ means $a+ib=-a+ib$, i.e. $2a=0$, so $z=ib$ is purely imaginary. $\blacksquare$

\begin{exemplebox}[Using the properties instead of recomputing]
Let $z=1+2i$ and $w=3-i$. Then $\overline{zw}=\bar{z}\bar{w}=(1-2i)(3+i)=3+i-6i-2i^{2}=5-5i$. Check directly: $zw=(1+2i)(3-i)=3-i+6i-2i^{2}=5+5i$, whose conjugate is indeed $5-5i$.
\end{exemplebox}

\begin{exoresolu}[Proving a number is real or purely imaginary]
Let $z$ be any complex number with $z\neq\pm 1$. Show that $w=\dfrac{z-\bar{z}}{1+z\bar{z}}\times i$ is purely imaginary if $z$ is purely imaginary, and find the value of $w$ when $z=3+4i$.
\tcblower
If $z=ib$ with $b\in\mathbb{R}$, then $\bar{z}=-ib$ and $z\bar{z}=b^{2}$, so
\[
w=\frac{2ib}{1+b^{2}}\times i=\frac{-2b}{1+b^{2}}\in\mathbb{R}.
\]
Hence in that case $w$ is in fact \emph{real} (a good reminder to compute before concluding!). For $z=3+4i$: $z-\bar{z}=8i$ and $z\bar{z}=9+16=25$, so
\[
w=\frac{8i}{26}\times i=-\frac{8}{26}=-\frac{4}{13}.
\]
The lesson: use the properties $\overline{z-\bar{z}}=-(z-\bar{z})$ and $z\bar{z}\in\mathbb{R}$ to predict the nature of an expression \emph{before} substituting numbers.
\end{exoresolu}

\begin{savaistu}[Why conjugates matter beyond the classroom]
Conjugates are not just an examination trick. In electrical engineering, the \emph{complex power} in an AC circuit is $S=V\bar{I}$, the product of the voltage by the \emph{conjugate} of the current; its real part is the useful power and its imaginary part the reactive power. The same conjugate product appears in quantum mechanics, where probabilities are computed as $z\bar{z}=|z|^{2}$ — always real and non-negative, exactly as a probability should be.
\end{savaistu}

\begin{exercice}[Practice]
\begin{enumerate}
\item Given $z_{1}=4-3i$ and $z_{2}=-2+5i$, compute $z_{1}+z_{2}$, $z_{1}-z_{2}$ and $z_{1}z_{2}$.
\item Express each of the following in the form $x+iy$: (a) $\dfrac{5+2i}{3-i}$; (b) $\dfrac{1}{2+i}$; (c) $\dfrac{(1+i)^{2}}{1-i}$.
\item Let $z=2-i$. Verify that $\overline{z^{2}}=(\bar{z})^{2}$.
\item Show that for every complex number $z$, the number $z+\bar{z}$ is real and $z-\bar{z}$ is purely imaginary.
\item Simplify $i^{11}+i^{13}+i^{15}$.
\end{enumerate}
\end{exercice}

\begin{corrige}[Exercise 1]
\textbf{1.} $z_{1}+z_{2}=2+2i$; $z_{1}-z_{2}=6-8i$; $z_{1}z_{2}=-8+20i+6i-15i^{2}=-8+26i+15=7+26i$.

\textbf{2.} (a) $\dfrac{5+2i}{3-i}=\dfrac{(5+2i)(3+i)}{10}=\dfrac{15+5i+6i+2i^{2}}{10}=\dfrac{13+11i}{10}=\dfrac{13}{10}+\dfrac{11}{10}i$.

(b) $\dfrac{1}{2+i}=\dfrac{2-i}{5}=\dfrac{2}{5}-\dfrac{1}{5}i$.

(c) $(1+i)^{2}=2i$, so $\dfrac{2i}{1-i}=\dfrac{2i(1+i)}{2}=i(1+i)=-1+i$.

\textbf{3.} $z^{2}=(2-i)^{2}=4-4i+i^{2}=3-4i$, so $\overline{z^{2}}=3+4i$. Also $(\bar{z})^{2}=(2+i)^{2}=4+4i+i^{2}=3+4i$. Both agree.

\textbf{4.} With $z=a+ib$: $z+\bar{z}=2a\in\mathbb{R}$ and $z-\bar{z}=2ib$, which is purely imaginary.

\textbf{5.} $i^{11}=i^{3}=-i$, $i^{13}=i$, $i^{15}=i^{3}=-i$, so the sum is $-i+i-i=-i$.
\end{corrige}

\begin{aretenir}[Key points of this section]
\begin{itemize}
\item Add and subtract component-wise: $(a+ib)\pm(c+id)=(a\pm c)+i(b\pm d)$.
\item Multiply by expanding and using $i^{2}=-1$: $(a+ib)(c+id)=(ac-bd)+i(ad+bc)$. Watch the sign of $bd$!
\item Powers of $i$ cycle with period $4$: $i^{2}=-1$, $i^{3}=-i$, $i^{4}=1$.
\item The conjugate of $z=a+ib$ is $\bar{z}=a-ib$; geometrically it is the reflection in the real axis.
\item $z+\bar{z}=2\operatorname{Re}(z)$, $z-\bar{z}=2i\operatorname{Im}(z)$, and $z\bar{z}=a^{2}+b^{2}=|z|^{2}\geq 0$.
\item To divide, multiply numerator and denominator by the conjugate of the \emph{denominator}.
\item Conjugation respects all operations: $\overline{z_{1}\pm z_{2}}=\bar{z}_{1}\pm\bar{z}_{2}$, $\overline{z_{1}z_{2}}=\bar{z}_{1}\bar{z}_{2}$, $\overline{z_{1}/z_{2}}=\bar{z}_{1}/\bar{z}_{2}$, $\overline{z^{n}}=(\bar{z})^{n}$; moreover $z=\bar{z}\iff z\in\mathbb{R}$ and $z=-\bar{z}\iff z$ is purely imaginary.
\end{itemize}
\end{aretenir}

We now have a complete arithmetic of complex numbers: addition, subtraction, multiplication and division all behave as in $\mathbb{R}$, with the conjugate as our essential tool. In the next section we put this algebra to work by solving equations in $\mathbb{C}$ — including the quadratic equations with negative discriminant that forced the invention of $i$ in the first place.

\coursec{Equations and Roots in \texorpdfstring{$\mathbb{C}$}{C}}

In the previous section we learnt to add, subtract, multiply and divide complex numbers, and we met the \cle{conjugate} $\bar z$ of a complex number. We now put this algebra to work. The whole reason complex numbers were invented was to \emph{solve equations}: the equation $x^2+1=0$ has no solution in $\mathbb{R}$, yet we shall see that in $\mathbb{C}$ \emph{every} polynomial equation has solutions. In this section we solve linear and quadratic equations with complex numbers, discover the beautiful \emph{conjugate pair} theorem, and learn to factorise expressions such as $a^2+b^2$ which are irreducible over $\mathbb{R}$.

\courssub{Linear Equations with Complex Coefficients}

A linear equation in the unknown $z\in\mathbb{C}$ has the form
\[ az+b=c, \qquad a,b,c\in\mathbb{C},\ a\neq 0. \]
It is solved exactly as in $\mathbb{R}$: isolate $z$, then divide by $a$ using the conjugate technique of Section 2.

\begin{methode}[Solving a linear equation in $\mathbb{C}$]
\begin{enumerate}
\item Collect the terms in $z$ on one side and the constants on the other.
\item Factorise $z$ if it appears in several terms.
\item Divide by the coefficient of $z$; if this coefficient is complex, multiply numerator and denominator by its conjugate.
\item Write the answer in the form $x+iy$.
\end{enumerate}
\end{methode}

\begin{exoresolu}[A linear equation with complex coefficients]
Solve in $\mathbb{C}$: $(2+i)z - 3i = 1 + z$.
\tcblower
\textbf{Solution.} Collect the terms in $z$:
\[ (2+i)z - z = 1+3i \quad\Longrightarrow\quad (1+i)z = 1+3i. \]
Divide by $1+i$, multiplying top and bottom by the conjugate $1-i$:
\[ z=\frac{1+3i}{1+i}=\frac{(1+3i)(1-i)}{(1+i)(1-i)}=\frac{1-i+3i-3i^2}{1+1}=\frac{4+2i}{2}=2+i. \]
\emph{Check:} $(2+i)(2+i)-3i = 4+4i+i^2-3i = 3+i$, and $1+(2+i)=3+i$. \checkmark{} Hence $\boxed{z=2+i}$.
\end{exoresolu}

\courssub{Quadratic Equations with Real Coefficients when $\Delta<0$}

\textbf{Intuition.} Consider $z^2-2z+5=0$. The discriminant is $\Delta=(-2)^2-4(1)(5)=4-20=-16<0$, so in $\mathbb{R}$ there is \emph{no} solution — the parabola never crosses the $x$-axis. But in $\mathbb{C}$ we can write $\sqrt{-16}=\sqrt{16}\,i=4i$, and the usual formula works perfectly.

\begin{popfigure}
\begin{center}
\begin{tikzpicture}
\begin{axis}[
  width=0.72\linewidth, height=5.6cm,
  axis lines=middle, xlabel=$x$, ylabel=$y$,
  xmin=-1.5, xmax=3.5, ymin=-1, ymax=7,
  xtick={1}, ytick=\empty,
  domain=-1.2:3.2, samples=100,
]
\addplot[popblue, thick] {x^2-2*x+5};
\addplot[popmuted, dashed] coordinates {(1,-0.8) (1,7)};
\node[popdark, align=center, text width=3.2cm] at (axis cs:2.6,5.6) {$y=x^2-2x+5$\\ no real intercept};
\node[popmuted] at (axis cs:1,-0.6) {$x=1$};
\end{axis}
\end{tikzpicture}
\end{center}
\legende{The parabola $y=x^2-2x+5$ has $\Delta=-16<0$: it never meets the $x$-axis, so there are no real roots. The roots live in $\mathbb{C}$: $z=1\pm 2i$.}
\end{popfigure}

\begin{propriete}[Quadratic formula in $\mathbb{C}$ — proof examinable]
Let $az^2+bz+c=0$ with $a,b,c\in\mathbb{R}$, $a\neq 0$, and discriminant $\Delta=b^2-4ac$.
\begin{itemize}
\item If $\Delta>0$: two distinct real roots $z=\dfrac{-b\pm\sqrt{\Delta}}{2a}$.
\item If $\Delta=0$: one (double) real root $z=-\dfrac{b}{2a}$.
\item If $\Delta<0$: two \cle{complex conjugate roots}
\[ z=\frac{-b\pm i\sqrt{|\Delta|}}{2a}. \]
\end{itemize}
\textbf{Proof (case $\Delta<0$).} Complete the square:
\[ a\left(z+\frac{b}{2a}\right)^2 = \frac{b^2-4ac}{4a}=\frac{\Delta}{4a}
\quad\Longrightarrow\quad \left(z+\frac{b}{2a}\right)^2=\frac{\Delta}{4a^2}. \]
Since $\Delta<0$, write $\Delta=-|\Delta|$, so $\dfrac{\Delta}{4a^2}=\left(\dfrac{i\sqrt{|\Delta|}}{2a}\right)^2$. Hence $z+\dfrac{b}{2a}=\pm\dfrac{i\sqrt{|\Delta|}}{2a}$, giving the result. $\blacksquare$
\end{propriete}

\begin{aretenir}[The fundamental fact about quadratics]
Every quadratic equation $az^2+bz+c=0$ ($a\neq 0$) has \textbf{exactly two roots in $\mathbb{C}$}, counted with multiplicity (a double root counts twice). This is the first taste of the \emph{Fundamental Theorem of Algebra}: every polynomial of degree $n$ has exactly $n$ roots in $\mathbb{C}$, counted with multiplicity.
\end{aretenir}

\begin{methode}[Solving $az^2+bz+c=0$ when $\Delta<0$]
\begin{enumerate}
\item Compute $\Delta=b^2-4ac$ and note that $\Delta<0$.
\item Compute $\sqrt{|\Delta|}$ and write $\sqrt{\Delta}=i\sqrt{|\Delta|}$.
\item Apply the formula $z=\dfrac{-b\pm i\sqrt{|\Delta|}}{2a}$.
\item Simplify each root and give both answers in the form $x\pm iy$.
\item \emph{Check (optional):} verify that the sum of the roots is $-\dfrac{b}{a}$ and the product is $\dfrac{c}{a}$.
\end{enumerate}
\end{methode}

\begin{exoresolu}[Solving a quadratic with $\Delta<0$]
Solve in $\mathbb{C}$: $z^2-4z+13=0$.
\tcblower
\textbf{Solution.} $\Delta=(-4)^2-4(1)(13)=16-52=-36<0$, and $\sqrt{|\Delta|}=6$.
\[ z=\frac{4\pm 6i}{2}=2\pm 3i. \]
\emph{Check:} sum $=(2+3i)+(2-3i)=4=-\dfrac{b}{a}$ \checkmark; product $=(2+3i)(2-3i)=4-9i^2=4+9=13=\dfrac{c}{a}$ \checkmark.
The roots are $\boxed{z=2+3i \ \text{and}\ z=2-3i}$.
\end{exoresolu}

\begin{attention}[The square root of a negative number]
A very common error is to write $\sqrt{-9}=-3$. This is \textbf{wrong}: by convention, for $a>0$,
\[ \sqrt{-a}=i\sqrt{a}, \qquad\text{so}\qquad \sqrt{-9}=3i\ (\text{not } -3). \]
Also beware: the rule $\sqrt{p}\,\sqrt{q}=\sqrt{pq}$ is \emph{only valid for non-negative} $p,q$. Indeed $\sqrt{-1}\sqrt{-1}=i\cdot i=-1$, whereas $\sqrt{(-1)(-1)}=\sqrt{1}=1$. Never apply real-number root rules blindly to negative numbers.
\end{attention}

\textbf{Sum and product of the roots.} For real coefficients, the two complex roots are conjugates: $z_1=\alpha+i\beta$, $z_2=\alpha-i\beta$. Therefore
\[ z_1+z_2=2\alpha=-\frac{b}{a}\in\mathbb{R}, \qquad z_1z_2=\alpha^2+\beta^2=\frac{c}{a}\in\mathbb{R}. \]
The sum and product remain \emph{real} — exactly as the classical relations demand. This is a quick and powerful check.

\courssub{Factorising in \texorpdfstring{$\mathbb{C}$}{C}}

Over $\mathbb{R}$, the expression $a^2+b^2$ cannot be factorised. Over $\mathbb{C}$, since $i^2=-1$:
\[ (a+ib)(a-ib)=a^2-(ib)^2=a^2-i^2b^2=a^2+b^2. \]

\begin{propriete}[Sum of two squares factorises in $\mathbb{C}$]
For all real (or complex) $a,b$:
\[ a^2+b^2=(a+ib)(a-ib). \]
\end{propriete}

\begin{exemplebox}[Factorising over $\mathbb{C}$]
\begin{itemize}
\item $z^2+9=(z+3i)(z-3i)$, so $z^2+9=0$ has roots $z=\pm 3i$.
\item $z^2-2z+5=(z-1)^2+4=(z-1+2i)(z-1-2i)$, recovering the roots $1\pm 2i$ of the parabola above.
\end{itemize}
\end{exemplebox}

\courssub{Non-Real Roots of Real Polynomials Occur in Conjugate Pairs}

\begin{propriete}[Conjugate Root Theorem]
If $P$ is a polynomial with \textbf{real} coefficients and $z_0$ is a root of $P$, then $\overline{z_0}$ is also a root of $P$. In other words, non-real roots of real polynomials occur in \cle{conjugate pairs}.
\end{propriete}

\textbf{Proof (instructive, uses the properties of conjugates).} Let $P(z)=a_nz^n+\dots+a_1z+a_0$ with all $a_k\in\mathbb{R}$, and suppose $P(z_0)=0$. Using $\overline{z+w}=\bar z+\bar w$, $\overline{zw}=\bar z\,\bar w$, and $\overline{a_k}=a_k$ (real numbers equal their own conjugates):
\[ P(\overline{z_0})=\sum_{k=0}^{n}a_k\overline{z_0}^{\,k}=\sum_{k=0}^{n}\overline{a_k z_0^{k}}=\overline{\sum_{k=0}^{n}a_k z_0^{k}}=\overline{P(z_0)}=\bar 0=0. \quad\blacksquare \]

\begin{attention}[Real coefficients are essential]
The theorem fails if the coefficients are not all real. For example $z^2-(2+i)z+2i=0$ factorises as $(z-2)(z-i)=0$, so its roots are $z=2$ and $z=i$ — and $\bar\imath=-i$ is \emph{not} a root. Always check the hypothesis ``real coefficients'' before quoting conjugate pairs.
\end{attention}

\begin{exemplebox}[Using conjugate pairs for a cubic]
The cubic $P(z)=z^3-3z^2+7z-5$ has real coefficients, and $z=1$ is a root since $P(1)=1-3+7-5=0$. Factorising: $P(z)=(z-1)(z^2-2z+5)$. The quadratic factor has $\Delta=-16$, giving the conjugate pair $1\pm 2i$. The three roots are $1,\ 1+2i,\ 1-2i$ — note the non-real roots form a conjugate pair, as the theorem guarantees.
\end{exemplebox}

\begin{savaistu}[Why cubics forced mathematicians to accept complex numbers]
In the sixteenth century, Italian algebraists such as Cardano and Bombelli discovered that the formula for solving a cubic equation sometimes requires the square root of a negative number \emph{even when all three roots are real}. Working formally with these ``impossible'' quantities and cancelling them at the end gave the correct real answers. Complex numbers were thus accepted not as a curiosity, but as a necessity — the price to pay for solving equations completely.
\end{savaistu}

\courssub{Forming a Quadratic Equation from Given Complex Roots}

If a quadratic with real coefficients has roots $z_1$ and $z_2$, then with $S=z_1+z_2$ and $P=z_1z_2$:
\[ z^2-Sz+P=0. \]
When the roots are conjugates, $z_{1,2}=\alpha\pm i\beta$, we have $S=2\alpha$ and $P=\alpha^2+\beta^2=|z_1|^2$, both real — so the equation obtained automatically has real coefficients.

\begin{exoresolu}[Forming an equation from its roots]
Find the quadratic equation with real coefficients whose roots are $3+2i$ and $3-2i$.
\tcblower
\textbf{Solution.} Sum: $S=(3+2i)+(3-2i)=6$. Product:
\[ P=(3+2i)(3-2i)=9-(2i)^2=9+4=13. \]
Hence the equation is $\boxed{z^2-6z+13=0}$.
\emph{Check:} $\Delta=36-52=-16$, so $z=\dfrac{6\pm 4i}{2}=3\pm 2i$. \checkmark
\end{exoresolu}

\courssub{Square Roots of Simple Complex Numbers — First Steps}

To solve $z^2=w$ for a given complex number $w$ is to find the \cle{square roots} of $w$. Every non-zero complex number has exactly \emph{two} square roots, which are opposites of each other. For simple numbers, inspection works well; the general algebraic method will be developed fully in Section 4.

\begin{exemplebox}[Square roots by inspection]
\begin{itemize}
\item Square roots of $-4$: since $(2i)^2=4i^2=-4$, the square roots of $-4$ are $\pm 2i$.
\item Square roots of $i$: try $z=\frac{\sqrt2}{2}(1+i)$. Then $z^2=\frac{1}{2}(1+i)^2=\frac{1}{2}(1+2i-1)=i$. Hence the square roots of $i$ are $\pm\frac{\sqrt2}{2}(1+i)$.
\item Square roots of $3-4i$: try $z=2-i$: $(2-i)^2=4-4i+i^2=3-4i$. Hence the roots are $\pm(2-i)$.
\end{itemize}
\end{exemplebox}

\begin{methode}[Algebraic method — preview]
To find $z=x+iy$ with $z^2=a+ib$:
\begin{enumerate}
\item Expand: $x^2-y^2=a$ and $2xy=b$ (equating real and imaginary parts).
\item Use moduli: $x^2+y^2=\sqrt{a^2+b^2}$.
\item Solve for $x^2$ and $y^2$; the sign of $b$ tells you whether $x$ and $y$ have the same sign ($b>0$) or opposite signs ($b<0$).
\end{enumerate}
Full details and worked examples follow in Section 4.
\end{methode}

\begin{exercice}[Practice]
\begin{enumerate}
\item Solve in $\mathbb{C}$: (a) $(1+2i)z=3-i$; (b) $z^2+2z+10=0$; (c) $2z^2-6z+13=0$.
\item Factorise over $\mathbb{C}$: $z^2+25$ and $z^2-6z+25$.
\item Given that $2-i$ is a root of $z^3-5z^2+9z-5=0$, find all three roots.
\item Form the quadratic equation with real coefficients having roots $-1\pm i\sqrt{3}$.
\item Find, by inspection, the square roots of $-16$ and of $5+12i$.
\end{enumerate}
\end{exercice}

\begin{corrige}[Exercise 1]
\textbf{1(a)} $z=\dfrac{3-i}{1+2i}=\dfrac{(3-i)(1-2i)}{(1+2i)(1-2i)}=\dfrac{3-6i-i+2i^2}{5}=\dfrac{1-7i}{5}$.
\textbf{1(b)} $\Delta=4-40=-36$; $z=\dfrac{-2\pm 6i}{2}=-1\pm 3i$.
\textbf{1(c)} $\Delta=36-104=-68$; $z=\dfrac{6\pm 2i\sqrt{17}}{4}=\dfrac{3\pm i\sqrt{17}}{2}$.
\textbf{2.} $z^2+25=(z+5i)(z-5i)$; $z^2-6z+25=(z-3)^2+16=(z-3+4i)(z-3-4i)$.
\textbf{3.} By the conjugate root theorem $2+i$ is also a root; $(z-(2-i))(z-(2+i))=z^2-4z+5$. Dividing: $z^3-5z^2+9z-5=(z-1)(z^2-4z+5)$, so the roots are $1,\ 2+i,\ 2-i$.
\textbf{4.} $S=-2$, $P=(-1)^2+(\sqrt3)^2=1+3=4$; equation: $z^2+2z+4=0$.
\textbf{5.} $\pm 4i$; $(3+2i)^2=9+12i+4i^2=5+12i$, so the roots are $\pm(3+2i)$.
\end{corrige}

\begin{aretenir}[Key points of the section]
\begin{itemize}
\item Linear equations in $\mathbb{C}$ are solved as in $\mathbb{R}$, using conjugate division at the end.
\item Every quadratic has exactly two roots in $\mathbb{C}$ (counted with multiplicity); when the coefficients are real and $\Delta<0$, the roots are the conjugate pair $z=\dfrac{-b\pm i\sqrt{|\Delta|}}{2a}$.
\item $\sqrt{-a}=i\sqrt{a}$ for $a>0$ — never write $\sqrt{-9}=-3$.
\item $a^2+b^2=(a+ib)(a-ib)$: every sum of squares factorises in $\mathbb{C}$.
\item Non-real roots of real polynomials occur in conjugate pairs; sum and product of conjugate roots are real.
\item A quadratic with roots $z_1,z_2$ is $z^2-(z_1+z_2)z+z_1z_2=0$.
\end{itemize}
\end{aretenir}

\coursec{The Argand Diagram, Modulus and Argument}

In the previous section we solved quadratic and higher-degree equations in $\mathbb{C}$ using purely algebraic manipulations.  We now turn to a geometric viewpoint that not only gives a vivid picture of complex numbers but also provides powerful tools for computing powers, roots, and solving equations.  The Argand diagram, named after the Swiss mathematician Jean-Robert Argand (1806), identifies each complex number $z = a + ib$ with a point $M(a,b)$ or a vector $\overrightarrow{OM}$ in a Cartesian plane.  This geometric representation turns addition into vector addition, multiplication into a rotation combined with a scaling, and makes the modulus and argument the natural “polar coordinates” of a complex number.

\begin{savaistu}[Historical Note]
Although Argand published his geometric interpretation in 1806, the idea was also developed independently by the Norwegian-Danish surveyor Caspar Wessel (1797) and later by Gauss (1831).  Gauss coined the term “complex number” and the notation $i$ for $\sqrt{-1}$, and his endorsement made the geometric representation standard in mathematics.
\end{savaistu}

\courssub{The Argand Diagram: Geometric Representation}

\begin{definition}[Argand Diagram]
The \cle{Argand diagram} (or complex plane) is a Cartesian coordinate system where the horizontal axis represents the real part and the vertical axis represents the imaginary part of a complex number.  A complex number $z = a + ib$ is represented by the point $M(a,b)$ or by the position vector $\overrightarrow{OM}$.
\end{definition}

The real axis is often labelled $\mathbb{R}$ or $Ox$, and the imaginary axis $i\mathbb{R}$ or $Oy$.  The origin $O$ corresponds to $0 = 0 + i0$.  Purely real numbers lie on the horizontal axis; purely imaginary numbers lie on the vertical axis.

\begin{popfigure}
\begin{tikzpicture}[scale=1.2, >=Stealth]
    % Axes
    \draw[->, thick] (-1,0) -- (4.5,0) node[right] {$\mathbb{R}$};
    \draw[->, thick] (0,-1) -- (0,3.5) node[above] {$i\mathbb{R}$};
    % Grid
    \draw[popmuted, very thin] (-0.5,-0.5) grid (4.5,3.5);
    % Point z = 3+2i
    \coordinate (O) at (0,0);
    \coordinate (M) at (3,2);
    \draw[popblue, thick, ->] (O) -- (M) node[midway, above left, popblue] {$\overrightarrow{OM}$};
    \fill[popblue] (M) circle (2pt) node[above right] {$M(3,2)$};
    % Modulus as hypotenuse
    \draw[poporange, dashed] (3,0) -- (M) node[midway, right, poporange] {$2$};
    \draw[poporange, dashed] (0,2) -- (M) node[midway, above, poporange] {$3$};
    \draw[poporange, thick] (O) -- (M) node[midway, below left, poporange] {$|z| = \sqrt{13}$};
    % Argument angle
    \draw[popgreen, thick] (0.8,0) arc (0:33.69:0.8) node[midway, right, popgreen] {$\theta = \arg(z)$};
    % Labels
    \node[below left] at (O) {$O$};
    \node[below] at (3,0) {$3$};
    \node[left] at (0,2) {$2$};
\end{tikzpicture}
\legende{The complex number $z = 3 + 2i$ on the Argand diagram.  The modulus $|z|$ is the length of $\overrightarrow{OM}$; the argument $\arg(z)$ is the oriented angle from the positive real axis to $\overrightarrow{OM}$.}
\end{popfigure}

\courssub{Modulus of a Complex Number}

The distance from the origin to the point $M(a,b)$ is the \cle{modulus} (or absolute value) of $z$.

\begin{definition}[Modulus]
For $z = a + ib \in \mathbb{C}$, the \cle{modulus} of $z$ is
\[
|z| = \sqrt{a^2 + b^2}.
\]
Geometrically, $|z| = OM$, the length of the vector $\overrightarrow{OM}$.
\end{definition}

\begin{propriete}[Properties of the Modulus]
For any $z, z_1, z_2 \in \mathbb{C}$:
\begin{enumerate}
    \item $|z| = |\overline{z}|$  (modulus is unchanged by conjugation).
    \item $|z_1 z_2| = |z_1|\,|z_2|$  (modulus of a product is the product of moduli).
    \item $\left|\dfrac{z_1}{z_2}\right| = \dfrac{|z_1|}{|z_2|}$ for $z_2 \neq 0$  (modulus of a quotient is the quotient of moduli).
    \item $|z| \ge 0$ with equality iff $z = 0$.
    \item Triangle inequality: $|z_1 + z_2| \le |z_1| + |z_2|$.
    \item Reverse triangle inequality: $||z_1| - |z_2|| \le |z_1 + z_2|$.
\end{enumerate}
\end{propriete}

\begin{exemplebox}[Computing Moduli]
Let $z_1 = 1 + i\sqrt{3}$ and $z_2 = -2 + 2i$.
\begin{align*}
|z_1| &= \sqrt{1^2 + (\sqrt{3})^2} = \sqrt{4} = 2,\\
|z_2| &= \sqrt{(-2)^2 + 2^2} = \sqrt{8} = 2\sqrt{2},\\
|z_1 z_2| &= |z_1|\,|z_2| = 2 \times 2\sqrt{2} = 4\sqrt{2}.
\end{align*}
Check: $z_1 z_2 = (1+i\sqrt{3})(-2+2i) = -2+2i -2i\sqrt{3} -2\sqrt{3} = -2(1+\sqrt{3}) + 2i(1-\sqrt{3})$.
$|z_1 z_2| = \sqrt{4(1+\sqrt{3})^2 + 4(1-\sqrt{3})^2} = 2\sqrt{(1+2\sqrt{3}+3)+(1-2\sqrt{3}+3)} = 2\sqrt{8} = 4\sqrt{2}$.
\end{exemplebox}

\courssub{Argument of a Complex Number}

The \cle{argument} of a non-zero complex number is the oriented angle that the vector $\overrightarrow{OM}$ makes with the positive real axis.

\begin{definition}[Argument]
For $z = a + ib \neq 0$, the \cle{argument} of $z$, denoted $\arg(z)$, is the angle $\theta \in (-\pi, \pi]$ such that
\[
\cos\theta = \frac{a}{|z|},\qquad \sin\theta = \frac{b}{|z|},\qquad \tan\theta = \frac{b}{a}\ (a \neq 0).
\]
The unique value in $(-\pi, \pi]$ is called the \cle{principal argument} and is often written $\operatorname{Arg}(z)$ (capital A).  The general argument is $\arg(z) = \operatorname{Arg}(z) + 2k\pi,\ k \in \mathbb{Z}$.
\end{definition}

\begin{attention}[Argument of Zero]
The argument of $z = 0$ is \textbf{undefined} because the vector $\overrightarrow{OM}$ has zero length and no direction.
\end{attention}

\begin{methode}[Finding the Principal Argument $\operatorname{Arg}(z)$]
\begin{enumerate}
    \item Compute $|z| = \sqrt{a^2+b^2}$.
    \item Determine the quadrant of the point $M(a,b)$:
          \begin{itemize}
              \item $a>0,\ b>0$  $\rightarrow$ Quadrant I: $\operatorname{Arg}(z) = \arctan\left(\dfrac{b}{a}\right)$.
              \item $a<0,\ b>0$  $\rightarrow$ Quadrant II: $\operatorname{Arg}(z) = \pi + \arctan\left(\dfrac{b}{a}\right)$ (since $\arctan(b/a)$ is negative here).
              \item $a<0,\ b<0$  $\rightarrow$ Quadrant III: $\operatorname{Arg}(z) = -\pi + \arctan\left(\dfrac{b}{a}\right)$ (since $\arctan(b/a)$ is positive here).
              \item $a>0,\ b<0$  $\rightarrow$ Quadrant IV: $\operatorname{Arg}(z) = \arctan\left(\dfrac{b}{a}\right)$ (negative angle).
          \end{itemize}
    \item Special cases on axes:
          \begin{itemize}
              \item $a>0,\ b=0$: $\operatorname{Arg}(z) = 0$.
              \item $a<0,\ b=0$: $\operatorname{Arg}(z) = \pi$.
              \item $a=0,\ b>0$: $\operatorname{Arg}(z) = \dfrac{\pi}{2}$.
              \item $a=0,\ b<0$: $\operatorname{Arg}(z) = -\dfrac{\pi}{2}$.
          \end{itemize}
\end{enumerate}
\end{methode}

\begin{attention}[The Quadrant Trap]
Never write $\arg(z) = \arctan(b/a)$ without checking the quadrant of $M(a,b)$.  The function $\arctan$ only returns values in $(-\pi/2, \pi/2)$, so it gives the correct principal argument \emph{only} when $a > 0$ (Quadrants I and IV).  For $a < 0$ you must add or subtract $\pi$ to land in the correct quadrant.
\end{attention}

\begin{exemplebox}[Arguments in Different Quadrants]
Find $\operatorname{Arg}(z)$ for:
\begin{enumerate}
    \item $z_1 = 1 + i\sqrt{3}$: $a>0,b>0$ (QI). $\tan\theta = \sqrt{3} \Rightarrow \theta = \pi/3$.
    \item $z_2 = -1 + i\sqrt{3}$: $a<0,b>0$ (QII). $\tan\theta = -\sqrt{3}$. $\arctan(-\sqrt{3}) = -\pi/3$ (QIV), so add $\pi$: $\operatorname{Arg}(z_2) = \pi - \pi/3 = 2\pi/3$.
    \item $z_3 = -1 - i\sqrt{3}$: $a<0,b<0$ (QIII). $\tan\theta = \sqrt{3}$. $\arctan(\sqrt{3}) = \pi/3$ (QI), so subtract $\pi$: $\operatorname{Arg}(z_3) = -\pi + \pi/3 = -2\pi/3$.
    \item $z_4 = 1 - i\sqrt{3}$: $a>0,b<0$ (QIV). $\tan\theta = -\sqrt{3} \Rightarrow \theta = -\pi/3$.
    \item $z_5 = -2$: $a<0,b=0$. $\operatorname{Arg}(z_5) = \pi$.
    \item $z_6 = -3i$: $a=0,b<0$. $\operatorname{Arg}(z_6) = -\pi/2$.
\end{enumerate}
\end{exemplebox}

\courssub{Geometric Interpretation of Operations}

\begin{popfigure}
\begin{tikzpicture}[scale=1.2, >=Stealth]
    \draw[->, thick] (-3,0) -- (3,0) node[right] {$\mathbb{R}$};
    \draw[->, thick] (0,-3) -- (0,3) node[above] {$i\mathbb{R}$};
    \draw[popmuted, very thin] (-2.5,-2.5) grid (2.5,2.5);
    % z = 2+i
    \coordinate (z) at (2,1);
    \draw[popblue, thick, ->] (0,0) -- (z) node[midway, above right, popblue] {$z$};
    \fill[popblue] (z) circle (2pt) node[above right] {$z$};
    % conjugate = 2-i
    \coordinate (zbar) at (2,-1);
    \draw[poporange, thick, ->] (0,0) -- (zbar) node[midway, below right, poporange] {$\overline{z}$};
    \fill[poporange] (zbar) circle (2pt) node[below right] {$\overline{z}$};
    % -z = -2-i
    \coordinate (mz) at (-2,-1);
    \draw[popgreen, thick, ->] (0,0) -- (mz) node[midway, below left, popgreen] {$-z$};
    \fill[popgreen] (mz) circle (2pt) node[below left] {$-z$};
    % -conjugate = -2+i
    \coordinate (mzbar) at (-2,1);
    \draw[poppurple, thick, ->] (0,0) -- (mzbar) node[midway, above left, poppurple] {$-\overline{z}$};
    \fill[poppurple] (mzbar) circle (2pt) node[above left] {$-\overline{z}$};
    % Reflection lines
    \draw[dashed, thin] (2,-1.5) -- (2,1.5);
    \draw[dashed, thin] (-2.5,0) -- (2.5,0);
    \node[popblue, align=center] at (2.7,1.5) {Reflection in\\ real axis:\\ $z \leftrightarrow \overline{z}$};
    \node[popgreen, align=center] at (-2.7,-1.5) {Rotation by $\pi$:\\ $z \leftrightarrow -z$};
\end{tikzpicture}
\legende{Symmetries on the Argand diagram: conjugation reflects across the real axis; negation rotates by $\pi$ (half-turn about the origin).}
\end{popfigure}

\begin{propriete}[Geometric Meaning of Conjugation and Negation]
\begin{itemize}
    \item $\overline{z}$ is the reflection of $z$ across the real axis.  Hence $|\overline{z}| = |z|$ and $\arg(\overline{z}) = -\arg(z)$ (mod $2\pi$).
    \item $-z$ is the point symmetric to $z$ with respect to the origin (rotation by $\pi$).  Hence $|-z| = |z|$ and $\arg(-z) = \arg(z) + \pi$ (mod $2\pi$).
\end{itemize}
\end{propriete}

Addition of complex numbers corresponds to vector addition (parallelogram rule).

\begin{popfigure}
\begin{tikzpicture}[scale=1.2, >=Stealth]
    \draw[->, thick] (-1,0) -- (5,0) node[right] {$\mathbb{R}$};
    \draw[->, thick] (0,-1) -- (0,4) node[above] {$i\mathbb{R}$};
    \draw[popmuted, very thin] (-0.5,-0.5) grid (4.5,3.5);
    \coordinate (O) at (0,0);
    \coordinate (z1) at (3,1);
    \coordinate (z2) at (1,2.5);
    \coordinate (sum) at (4,3.5);
    % Vectors
    \draw[popblue, thick, ->] (O) -- (z1) node[midway, below right, popblue] {$z_1$};
    \draw[poporange, thick, ->] (O) -- (z2) node[midway, above left, poporange] {$z_2$};
    \draw[popgreen, thick, ->] (O) -- (sum) node[midway, above right, popgreen] {$z_1+z_2$};
    % Parallelogram
    \draw[popblue, dashed] (z1) -- (sum);
    \draw[poporange, dashed] (z2) -- (sum);
    % Points
    \fill[popblue] (z1) circle (2pt) node[below right] {$z_1$};
    \fill[poporange] (z2) circle (2pt) node[above left] {$z_2$};
    \fill[popgreen] (sum) circle (2pt) node[above right] {$z_1+z_2$};
    \node[below left] at (O) {$O$};
\end{tikzpicture}
\legende{Addition on the Argand diagram: $z_1+z_2$ is the diagonal of the parallelogram formed by $\overrightarrow{Oz_1}$ and $\overrightarrow{Oz_2}$.}
\end{popfigure}

\begin{propriete}[Distance Between Two Points]
The distance between the points representing $z_1$ and $z_2$ is $|z_2 - z_1|$.  Indeed, $\overrightarrow{M_1M_2} = \overrightarrow{OM_2} - \overrightarrow{OM_1}$, so its length is $|z_2 - z_1|$.
\end{propriete}

\begin{exemplebox}[Distance]
Let $z_1 = 2+3i$ and $z_2 = -1+i$.  The distance between them is
\[
|z_2 - z_1| = |(-1-2) + i(1-3)| = |-3 - 2i| = \sqrt{(-3)^2 + (-2)^2} = \sqrt{13}.
\]
\end{exemplebox}

\begin{popfigure}
\begin{tikzpicture}[scale=1.2, >=Stealth]
    \draw[->, thick] (-1,0) -- (5,0) node[right] {$\mathbb{R}$};
    \draw[->, thick] (0,-1) -- (0,4) node[above] {$i\mathbb{R}$};
    \draw[popmuted, very thin] (-0.5,-0.5) grid (4.5,3.5);
    \coordinate (O) at (0,0);
    \coordinate (z1) at (3,1);
    \coordinate (z2) at (1,2.5);
    \coordinate (diff) at (2,-1.5);
    % Vectors
    \draw[popblue, thick, ->] (O) -- (z1) node[midway, below right, popblue] {$z_1$};
    \draw[poporange, thick, ->] (O) -- (z2) node[midway, above left, poporange] {$z_2$};
    \draw[popgreen, thick, ->] (z1) -- (z2) node[midway, right, popgreen] {$z_2-z_1$};
    % Points
    \fill[popblue] (z1) circle (2pt) node[below right] {$z_1$};
    \fill[poporange] (z2) circle (2pt) node[above left] {$z_2$};
    \node[below left] at (O) {$O$};
\end{tikzpicture}
\legende{The vector $z_2 - z_1$ joins $z_1$ to $z_2$; its length is the distance between the two points.}
\end{popfigure}

\courssub{Geometric Interpretation of Multiplication}

Multiplication by a complex number combines rotation and scaling (homothety).

\begin{propriete}[Modulus and Argument of a Product]
For $z_1, z_2 \in \mathbb{C}^*$:
\begin{align*}
|z_1 z_2| &= |z_1|\,|z_2|,\\
\arg(z_1 z_2) &= \arg(z_1) + \arg(z_2) \pmod{2\pi}.
\end{align*}
Geometrically, multiplying by $z_2$ rotates the plane by $\arg(z_2)$ and scales distances by $|z_2|$.
\end{propriete}

\begin{popfigure}
\begin{tikzpicture}[scale=1.2, >=Stealth]
    \draw[->, thick] (-3,0) -- (3,0) node[right] {$\mathbb{R}$};
    \draw[->, thick] (0,-3) -- (0,3) node[above] {$i\mathbb{R}$};
    \draw[popmuted, very thin] (-2.5,-2.5) grid (2.5,2.5);
    % z1 = 1+i
    \coordinate (z1) at (1,1);
    \draw[popblue, thick, ->] (0,0) -- (z1) node[midway, above right, popblue] {$z_1$};
    \fill[popblue] (z1) circle (2pt) node[above right] {$z_1$};
    % z2 = 1+i (same for simplicity)
    \coordinate (z2) at (1,1);
    % z1*z2 = 2i
    \coordinate (prod) at (0,2);
    \draw[popgreen, thick, ->] (0,0) -- (prod) node[midway, left, popgreen] {$z_1z_2$};
    \fill[popgreen] (prod) circle (2pt) node[left] {$z_1z_2$};
    % Angle arcs
    \draw[popblue, thick] (0.5,0) arc (0:45:0.5) node[midway, right, popblue] {$\pi/4$};
    \draw[popgreen, thick] (0,0.5) arc (90:90:0.5); % dummy
    \draw[poporange, thick] (0.5,0) arc (0:90:0.5) node[midway, above right, poporange] {$\pi/2$};
    \node[poporange, align=center] at (2.5,2) {Rotation by $\pi/4$\\ and scaling by $\sqrt{2}$};
\end{tikzpicture}
\legende{Multiplication by $z_2 = 1+i$ (modulus $\sqrt{2}$, argument $\pi/4$) rotates $z_1$ by $\pi/4$ and scales its length by $\sqrt{2}$.}
\end{popfigure}

\courssub{Square Root of a Complex Number (Algebraic Method)}

Lesson 32 teaches us to find the square roots of a given complex number $z = a+ib$ by solving $(x+iy)^2 = a+ib$ for real $x,y$.  This yields two opposite roots.

\begin{methode}[Algebraic Method for $\sqrt{a+ib}$]
To find the square roots of $z = a+ib$ ($a,b \in \mathbb{R}$):
\begin{enumerate}
    \item Let $w = x+iy$ with $x,y \in \mathbb{R}$ such that $w^2 = z$.
    \item Expand: $(x+iy)^2 = (x^2-y^2) + i(2xy) = a + ib$.
    \item Equate real and imaginary parts:
          \[\begin{cases}
              x^2 - y^2 = a & (1)\\
              2xy = b       & (2)
          \end{cases}\]
    \item Use the modulus: $|w|^2 = |z| \Rightarrow x^2 + y^2 = \sqrt{a^2+b^2}$  (3).
    \item Add (1) and (3): $2x^2 = a + \sqrt{a^2+b^2} \Rightarrow x^2 = \dfrac{a + \sqrt{a^2+b^2}}{2}$.
          Since the right-hand side is $\ge 0$, $x = \pm \sqrt{\dfrac{a + \sqrt{a^2+b^2}}{2}}$.
    \item Subtract (1) from (3): $2y^2 = \sqrt{a^2+b^2} - a \Rightarrow y^2 = \dfrac{\sqrt{a^2+b^2} - a}{2}$.
          So $y = \pm \sqrt{\dfrac{\sqrt{a^2+b^2} - a}{2}}$.
    \item Choose the signs of $x$ and $y$ so that $2xy = b$ (equation (2)).  This gives two opposite solutions: $w$ and $-w$.
\end{enumerate}
\end{methode}

\begin{aretenir}[Sign Rule for Square Roots]
\begin{itemize}
    \item If $b > 0$, then $x$ and $y$ have the \textbf{same sign} (both $+$ or both $-$).
    \item If $b < 0$, then $x$ and $y$ have \textbf{opposite signs}.
    \item If $b = 0$, then either $x=0$ (pure imaginary roots) or $y=0$ (real roots).
\end{itemize}
\end{aretenir}

\begin{exoresolu}[Square Roots of $3+4i$]
Find the two square roots of $z = 3+4i$.

\textbf{Solution.}
Let $w = x+iy$ with $x,y \in \mathbb{R}$ such that $w^2 = 3+4i$.
\begin{align*}
(x+iy)^2 &= x^2 - y^2 + 2ixy = 3 + 4i.
\end{align*}
Equating parts:
\[\begin{cases}
x^2 - y^2 = 3 & (1)\\
2xy = 4       & (2)
\end{cases}\]
Modulus: $|w|^2 = |z| = \sqrt{3^2+4^2} = 5$, so
\[
x^2 + y^2 = 5 \quad (3).
\]
Add (1) and (3): $2x^2 = 8 \Rightarrow x^2 = 4 \Rightarrow x = \pm 2$.
Subtract (1) from (3): $2y^2 = 2 \Rightarrow y^2 = 1 \Rightarrow y = \pm 1$.
Equation (2) requires $2xy = 4 \Rightarrow xy = 2 > 0$, so $x$ and $y$ have the \emph{same} sign.
Thus the two square roots are
\[
w_1 = 2 + i,\qquad w_2 = -2 - i = -w_1.
\]
Check: $(2+i)^2 = 4 + 4i + i^2 = 3+4i$; $(-2-i)^2 = (2+i)^2 = 3+4i$.
\end{exoresolu}

\begin{exoresolu}[Square Roots of $-5-12i$]
Find the square roots of $z = -5-12i$.

\textbf{Solution.}
$|z| = \sqrt{(-5)^2+(-12)^2} = 13$.
We need $x^2 - y^2 = -5$ and $2xy = -12 \Rightarrow xy = -6 < 0$ (opposite signs).
$x^2 + y^2 = 13$.
Adding: $2x^2 = 8 \Rightarrow x^2 = 4 \Rightarrow x = \pm 2$.
Subtracting: $2y^2 = 18 \Rightarrow y^2 = 9 \Rightarrow y = \pm 3$.
Since $xy = -6 < 0$, $x$ and $y$ have opposite signs.
Thus the roots are $2 - 3i$ and $-2 + 3i$.
Check: $(2-3i)^2 = 4 -12i +9i^2 = -5-12i$.
\end{exoresolu}

\begin{exoresolu}[Square Roots of a Real Negative Number]
Find the square roots of $z = -9$.

\textbf{Solution.}
Here $a = -9$, $b = 0$. $|z| = 9$.
Equations: $x^2 - y^2 = -9$, $2xy = 0$, $x^2 + y^2 = 9$.
From $2xy=0$, either $x=0$ or $y=0$.
If $y=0$, then $x^2 = -9$ (impossible for real $x$).
So $x=0$, then $-y^2 = -9 \Rightarrow y^2 = 9 \Rightarrow y = \pm 3$.
The square roots are $\pm 3i$.
\end{exoresolu}

\begin{aretenir}[Key Points of Section 4]
\begin{itemize}
    \item The Argand diagram represents $z = a+ib$ by $M(a,b)$ or $\overrightarrow{OM}$.
    \item Modulus $|z| = \sqrt{a^2+b^2}$ is the distance $OM$; it satisfies $|z_1z_2| = |z_1||z_2|$, $|z_1/z_2| = |z_1|/|z_2|$.
    \item Argument $\arg(z)$ is the oriented angle $(Ox, OM)$; principal value $\operatorname{Arg}(z) \in (-\pi,\pi]$.
    \item Always locate the quadrant of $M$ before using $\arctan(b/a)$.
    \item Conjugation $\leftrightarrow$ reflection in the real axis; negation $\leftrightarrow$ rotation by $\pi$.
    \item Addition $\leftrightarrow$ vector addition (parallelogram rule).
    \item Distance between $z_1$ and $z_2$ is $|z_2 - z_1|$.
    \item Multiplication: $|z_1z_2| = |z_1||z_2|$, $\arg(z_1z_2) = \arg(z_1) + \arg(z_2) \pmod{2\pi}$ (rotation + scaling).
    \item Square roots of $a+ib$: solve $(x+iy)^2 = a+ib$ using $x^2-y^2=a$, $2xy=b$, $x^2+y^2=\sqrt{a^2+b^2}$; the two roots are opposites.
    \item Sign rule: $b>0 \Rightarrow x,y$ same sign; $b<0 \Rightarrow x,y$ opposite signs.
\end{itemize}
\end{aretenir}

\begin{exercice}[Practice Exercises]
\begin{enumerate}
    \item Plot on an Argand diagram: $z_1 = 2-i$, $z_2 = -1+2i$, $z_3 = -3-2i$, $z_4 = 4i$.  Find the modulus and principal argument of each.
    \item Given $z_1 = 1+i\sqrt{3}$ and $z_2 = \sqrt{3}-i$, compute $|z_1|$, $|z_2|$, $|z_1z_2|$, $\operatorname{Arg}(z_1)$, $\operatorname{Arg}(z_2)$, and $\operatorname{Arg}(z_1z_2)$.  Verify that $\operatorname{Arg}(z_1z_2) = \operatorname{Arg}(z_1)+\operatorname{Arg}(z_2)$ (mod $2\pi$).
    \item Find the square roots of: (a) $7-24i$, (b) $-8+6i$, (c) $-4$.
    \item Let $z_1 = 3+4i$ and $z_2 = 1-2i$.  Represent $z_1$, $z_2$, $z_1+z_2$, $z_1-z_2$ on the same Argand diagram.  Compute the distance between $z_1$ and $z_2$.
    \item If $|z| = 2$ and $\operatorname{Arg}(z) = \frac{2\pi}{3}$, find the Cartesian form of $z$, $\overline{z}$, $-z$, and $z^2$.
    \item Show that the points representing $z_1 = 1+2i$, $z_2 = 3+i$, and $z_3 = 5$ form a right-angled triangle.  Identify the right angle.
    \item Find the complex numbers $z$ such that $|z| = 5$ and $\operatorname{Arg}(z) = \frac{\pi}{4}$ or $\operatorname{Arg}(z) = -\frac{3\pi}{4}$.  Explain the geometric relationship between the two solutions.
\end{enumerate}
\end{exercice}

\begin{corrige}[Exercise 1]
\begin{enumerate}
    \item $z_1=2-i$: $|z_1|=\sqrt{5}$, $\operatorname{Arg}(z_1)=-\arctan(1/2) \approx -0.4636$ rad.\\
          $z_2=-1+2i$: $|z_2|=\sqrt{5}$, $\operatorname{Arg}(z_2)=\pi-\arctan(2) \approx 2.034$ rad.\\
          $z_3=-3-2i$: $|z_3|=\sqrt{13}$, $\operatorname{Arg}(z_3)=-\pi+\arctan(2/3) \approx -2.554$ rad.\\
          $z_4=4i$: $|z_4|=4$, $\operatorname{Arg}(z_4)=\pi/2$.
    \item $|z_1|=2$, $|z_2|=2$, $|z_1z_2|=4$.  $\operatorname{Arg}(z_1)=\pi/3$, $\operatorname{Arg}(z_2)=-\pi/6$.  $z_1z_2 = (1+i\sqrt{3})(\sqrt{3}-i) = \sqrt{3}-i+3i+\sqrt{3} = 2\sqrt{3}+2i$.  $\operatorname{Arg}(z_1z_2)=\arctan(1/\sqrt{3})=\pi/6 = \pi/3 + (-\pi/6)$.
    \item (a) $|z|=25$, $x^2-y^2=7$, $2xy=-24 \Rightarrow xy=-12$.  $x^2+y^2=25 \Rightarrow x^2=16, y^2=9$.  Opposite signs $\Rightarrow \pm(4-3i)$.\\
          (b) $|z|=10$, $x^2-y^2=-8$, $2xy=6 \Rightarrow xy=3$.  $x^2+y^2=10 \Rightarrow x^2=1, y^2=9$.  Same sign $\Rightarrow \pm(1+3i)$.\\
          (c) $-4 = -4+0i$.  $x^2-y^2=-4$, $2xy=0 \Rightarrow x=0$ or $y=0$.  If $x=0$, $-y^2=-4 \Rightarrow y=\pm2$.  Roots: $\pm 2i$.
    \item $z_1+z_2=4+2i$, $z_1-z_2=2+6i$.  Distance $|z_2-z_1| = |-2-6i| = \sqrt{40}=2\sqrt{10}$.
    \item $z = 2(\cos\frac{2\pi}{3}+i\sin\frac{2\pi}{3}) = -1+i\sqrt{3}$.  $\overline{z} = -1-i\sqrt{3}$.  $-z = 1-i\sqrt{3}$.  $z^2 = 4(\cos\frac{4\pi}{3}+i\sin\frac{4\pi}{3}) = -2-2i\sqrt{3}$.
    \item Vectors: $z_2-z_1 = 2-i$, $z_3-z_2 = 2-i$, $z_3-z_1 = 4-2i$.  Check dot product: $(2)(2) + (-1)(-1) = 5 \neq 0$.  Wait: $z_2-z_1 = (3-1)+i(1-2) = 2-i$. $z_3-z_2 = (5-3)+i(0-1) = 2-i$. These are parallel! So points are collinear, not a triangle.  Correction: $z_1=1+2i$, $z_2=3+i$, $z_3=5$.  $z_2-z_1=2-i$, $z_3-z_2=2-i$. They are collinear.  The exercise statement has an error; with these points they are collinear.  (If $z_3=5+2i$, then $z_3-z_2=2+i$, dot product $(2)(2)+(-1)(1)=3\neq0$; $z_3-z_1=4$, dot product $(2)(4)+(-1)(0)=8\neq0$; $(2)(4)+(1)(0)=8\neq0$.  Not right-angled either.  A correct right triangle example: $z_1=0$, $z_2=3$, $z_3=3+4i$.)
    \item For $\operatorname{Arg}(z)=\pi/4$: $z = 5(\cos\frac{\pi}{4}+i\sin\frac{\pi}{4}) = \frac{5\sqrt{2}}{2}(1+i)$.  For $\operatorname{Arg}(z)=-\frac{3\pi}{4}$: $z = 5(\cos(-\frac{3\pi}{4})+i\sin(-\frac{3\pi}{4})) = -\frac{5\sqrt{2}}{2}(1+i)$.  The two numbers are opposites: $z_2 = -z_1$.  Geometrically, they are diametrically opposite on the circle $|z|=5$.
\end{enumerate}
\end{corrige}

\coursec{Polar and Exponential Forms}

In the previous section we met the \cle{Argand diagram}, and we saw that every complex number $z = a + ib$ can be located by its \cle{modulus} $r = |z|$ (its distance from the origin) and its \cle{argument} $\theta = \arg(z)$ (the angle its image makes with the positive real axis). This is a powerful idea: instead of describing a point by its Cartesian coordinates $(a, b)$, we can describe it by a \emph{distance} and a \emph{direction}. In this section we turn that idea into algebra. The result — the polar form and the exponential form — will transform multiplication of complex numbers from a tedious expansion into something geometrically beautiful: \emph{rotate and enlarge}.

\courssub{The Polar (Trigonometric) Form}

\textbf{From the diagram to a formula.}\par
Let $z = a + ib$ be a non-zero complex number with image $M(a, b)$ in the Argand diagram. Write $r = |z| = \sqrt{a^2 + b^2}$ and let $\theta$ be an argument of $z$. Looking at the right-angled triangle formed by $O$, $M$ and the foot of the perpendicular from $M$ to the real axis, elementary trigonometry gives
\[ a = r\cos\theta \qquad \text{and} \qquad b = r\sin\theta. \]
Substituting into $z = a + ib$ and factorising $r$:
\[ z = r\cos\theta + i\,r\sin\theta = r(\cos\theta + i\sin\theta). \]

\begin{definition}[Polar form of a complex number]
Let $z$ be a non-zero complex number with modulus $r = |z| > 0$ and argument $\theta = \arg(z)$. The \cle{polar form} (or \cle{trigonometric form}) of $z$ is
\[ z = r(\cos\theta + i\sin\theta). \]
The number $0$ has no polar form, since its argument is undefined.
\end{definition}

\begin{popfigure}
\begin{center}
\begin{tikzpicture}[scale=1.6,>=stealth]
\draw[->,thick,popmuted] (-0.6,0) -- (3.4,0) node[below left] {\small Re};
\draw[->,thick,popmuted] (0,-0.6) -- (0,2.6) node[below left] {\small Im};
\coordinate (O) at (0,0);
\coordinate (M) at (2.4,1.8);
\draw[->,very thick,popblue] (O) -- (M) node[above right] {$M:\ z = a + ib$};
\draw[dashed,popmuted] (M) -- (2.4,0) node[below] {$a$};
\draw[dashed,popmuted] (M) -- (0,1.8) node[left] {$b$};
\draw[thick,poporange] (0.9,0) arc (0:36.87:0.9);
\node[poporange] at (1.15,0.28) {$\theta$};
\node[popblue] at (1.05,1.15) {$r = |z|$};
\fill[popdark] (M) circle (1.3pt);
\end{tikzpicture}
\end{center}
\legende{Polar description of $z$: the modulus $r$ is the length $OM$, the argument $\theta$ is the angle from the positive real axis to $\vec{OM}$.}
\end{popfigure}

\begin{methode}[Converting $z = a + ib$ to polar form]
\begin{enumerate}
\item Compute the modulus: $r = \sqrt{a^2 + b^2}$.
\item Find an angle $\theta$ satisfying $\cos\theta = \dfrac{a}{r}$ and $\sin\theta = \dfrac{b}{r}$.
\item \textbf{Quadrant check (crucial):} use the signs of $a$ and $b$ to place $\theta$ in the correct quadrant.
  \begin{itemize}
  \item $a > 0,\ b > 0$ : $\theta \in (0, \tfrac{\pi}{2})$ (Q1)
  \item $a < 0,\ b > 0$ : $\theta \in (\tfrac{\pi}{2}, \pi)$ (Q2)
  \item $a < 0,\ b < 0$ : $\theta \in (-\pi, -\tfrac{\pi}{2})$ (Q3)
  \item $a > 0,\ b < 0$ : $\theta \in (-\tfrac{\pi}{2}, 0)$ (Q4)
  \end{itemize}
  A calculator gives $\arctan\!\left(\tfrac{b}{a}\right)$, which lies in $(-\tfrac{\pi}{2}, \tfrac{\pi}{2})$. If $a < 0$, add $\pi$ (or subtract $\pi$) to reach the correct quadrant.
\item Write $z = r(\cos\theta + i\sin\theta)$, choosing $\theta$ in the principal range $(-\pi, \pi]$ if required.
\end{enumerate}
For the reverse direction (polar to Cartesian), simply expand: $a = r\cos\theta$, $b = r\sin\theta$.
\end{methode}

\begin{exemplebox}[Cartesian to polar and back]
\textbf{(a)} Convert $z = -1 + i\sqrt{3}$ to polar form.\\
$r = \sqrt{(-1)^2 + (\sqrt{3})^2} = \sqrt{4} = 2$. Then $\cos\theta = -\tfrac{1}{2}$ and $\sin\theta = \tfrac{\sqrt{3}}{2}$. Since $a < 0$ and $b > 0$, $\theta$ lies in the second quadrant: $\theta = \tfrac{2\pi}{3}$. Hence
\[ z = 2\left(\cos\tfrac{2\pi}{3} + i\sin\tfrac{2\pi}{3}\right). \]
\textbf{(b)} Convert $w = 4\left(\cos\tfrac{\pi}{6} + i\sin\tfrac{\pi}{6}\right)$ to Cartesian form.\\
$a = 4\cos\tfrac{\pi}{6} = 4\cdot\tfrac{\sqrt{3}}{2} = 2\sqrt{3}$ and $b = 4\sin\tfrac{\pi}{6} = 2$, so $w = 2\sqrt{3} + 2i$.
\textbf{(c)} Convert $z = -2 - 2i$ to polar form with principal argument.\\
$r = \sqrt{(-2)^2 + (-2)^2} = 2\sqrt{2}$. $\cos\theta = -\tfrac{1}{\sqrt{2}}$, $\sin\theta = -\tfrac{1}{\sqrt{2}}$. Both negative $\Rightarrow$ third quadrant. Calculator gives $\arctan(1) = \tfrac{\pi}{4}$; add $\pi$ (or subtract $\pi$) to get $\theta = -\tfrac{3\pi}{4}$ (principal). Thus $z = 2\sqrt{2}\left(\cos\left(-\tfrac{3\pi}{4}\right) + i\sin\left(-\tfrac{3\pi}{4}\right)\right)$.
\end{exemplebox}

\begin{attention}[The argument is defined modulo $2\pi$]
The angle $\theta$ in the polar form is \emph{not unique}: if $\theta$ works, so does $\theta + 2k\pi$ for any $k \in \mathbb{Z}$. For example $1 + i = \sqrt{2}\left(\cos\tfrac{\pi}{4} + i\sin\tfrac{\pi}{4}\right) = \sqrt{2}\left(\cos\tfrac{9\pi}{4} + i\sin\tfrac{9\pi}{4}\right)$. When a single value is demanded, give the \cle{principal argument}, the unique value in $(-\pi, \pi]$. A frequent GCE error is to read $\arctan$ off a calculator without the quadrant check: for $z = -1 - i$, the calculator gives $\arctan(1) = \tfrac{\pi}{4}$, but the true principal argument is $-\tfrac{3\pi}{4}$.
\end{attention}

\courssub{Multiplication and Division in Polar Form}

Here is where the polar form earns its keep. Take two complex numbers
\[ z_1 = r_1(\cos\theta_1 + i\sin\theta_1), \qquad z_2 = r_2(\cos\theta_2 + i\sin\theta_2). \]
Multiplying them directly:
\begin{align*}
z_1 z_2 &= r_1 r_2\left[(\cos\theta_1\cos\theta_2 - \sin\theta_1\sin\theta_2) + i(\sin\theta_1\cos\theta_2 + \cos\theta_1\sin\theta_2)\right] \\
&= r_1 r_2\left[\cos(\theta_1 + \theta_2) + i\sin(\theta_1 + \theta_2)\right],
\end{align*}
where we used the addition formulae $\cos(\theta_1+\theta_2) = \cos\theta_1\cos\theta_2 - \sin\theta_1\sin\theta_2$ and $\sin(\theta_1+\theta_2) = \sin\theta_1\cos\theta_2 + \cos\theta_1\sin\theta_2$. This derivation is short, instructive, and \textbf{examinable}.

\begin{propriete}[Modulus and argument of a product and a quotient]
For non-zero complex numbers $z_1, z_2$:
\[ |z_1 z_2| = |z_1|\,|z_2|, \qquad \arg(z_1 z_2) = \arg(z_1) + \arg(z_2) \pmod{2\pi}, \]
\[ \left|\frac{z_1}{z_2}\right| = \frac{|z_1|}{|z_2|}, \qquad \arg\!\left(\frac{z_1}{z_2}\right) = \arg(z_1) - \arg(z_2) \pmod{2\pi}. \]
In words: \emph{to multiply, multiply the moduli and add the arguments; to divide, divide the moduli and subtract the arguments.}
\end{propriete}

The quotient rule follows the same way: multiply numerator and denominator by the conjugate of the denominator, or simply note that $\dfrac{z_1}{z_2} = z_1 \cdot \dfrac{1}{z_2}$ and that $\dfrac{1}{z_2} = \dfrac{1}{r_2}\bigl(\cos(-\theta_2) + i\sin(-\theta_2)\bigr)$, since $\cos$ is even and $\sin$ is odd.

\begin{exemplebox}[Product and quotient in polar form]
Let $z_1 = 3\left(\cos\tfrac{\pi}{3} + i\sin\tfrac{\pi}{3}\right)$ and $z_2 = 2\left(\cos\tfrac{\pi}{6} + i\sin\tfrac{\pi}{6}\right)$. Then
\[ z_1 z_2 = 6\left(\cos\tfrac{\pi}{2} + i\sin\tfrac{\pi}{2}\right) = 6i, \qquad \frac{z_1}{z_2} = \frac{3}{2}\left(\cos\tfrac{\pi}{6} + i\sin\tfrac{\pi}{6}\right) = \frac{3\sqrt{3}}{4} + \frac{3}{4}i. \]
Notice how the first answer required no expansion of brackets at all.
\end{exemplebox}

\textbf{Geometric interpretation: rotation and enlargement.}\par
The rule $\arg(z_1 z_2) = \arg z_1 + \arg z_2$ has a striking meaning. Fix $z$ with $|z| = r$ and $\arg z = \theta$. For \emph{any} complex number $w$, the product $zw$ is obtained from $w$ by:
\begin{itemize}
\item \textbf{rotating} the vector $\vec{OW}$ about $O$ through the angle $\theta$ (arguments add), and
\item \textbf{enlarging} it from $O$ by the scale factor $r$ (moduli multiply).
\end{itemize}
Multiplication by a complex number is exactly a \cle{rotation} composed with an \cle{enlargement} (a spiral similarity). In particular, since $|i| = 1$ and $\arg(i) = \tfrac{\pi}{2}$, multiplication by $i$ is a rotation through $90^\circ$ anticlockwise — check: $i(a + ib) = -b + ia$, which is indeed the image of $(a,b)$ rotated by a quarter turn.

\begin{popfigure}
\begin{center}
\begin{tikzpicture}[scale=1.15,>=stealth]
\draw[->,thick,popmuted] (-0.8,0) -- (5.6,0) node[below left] {\small Re};
\draw[->,thick,popmuted] (0,-0.8) -- (0,4.4) node[below left] {\small Im};
\coordinate (O) at (0,0);
\coordinate (Z1) at (3.2,1.0);
\coordinate (Z2) at (1.2,1.6);
\coordinate (P) at (2.24,5.12);
\draw[->,very thick,popblue] (O) -- (Z1) node[below right] {$z_1$};
\draw[->,very thick,popgreen] (O) -- (Z2) node[above left] {$z_2$};
\draw[->,very thick,poporange] (O) -- (P) node[above right] {$z_1 z_2$};
\draw[thick,popblue] (1.1,0) arc (0:17.35:1.1);
\node[popblue] at (1.45,0.22) {$\theta_1$};
\draw[thick,popgreen] (2.0,0) arc (0:53.13:2.0);
\node[popgreen] at (1.7,1.15) {$\theta_2$};
\draw[thick,poporange] (0.85,0) arc (0:66.37:0.85);
\node[poporange] at (0.62,0.95) {$\theta_1{+}\theta_2$};
\fill[popdark] (Z1) circle (1.3pt);
\fill[popdark] (Z2) circle (1.3pt);
\fill[popdark] (P) circle (1.3pt);
\end{tikzpicture}
\end{center}
\legende{Multiplication as rotation--enlargement: the argument of $z_1z_2$ is $\theta_1 + \theta_2$ and its modulus is $|z_1||z_2|$.}
\end{popfigure}

\courssub{Euler's Formula and the Exponential Form}

\textbf{A remarkable shortcut.}\par
The expression $\cos\theta + i\sin\theta$ appears so often that mathematicians looked for a shorthand. The clue comes from the laws we have just proved: writing $f(\theta) = \cos\theta + i\sin\theta$, the product rule says
\[ f(\theta_1)\,f(\theta_2) = f(\theta_1 + \theta_2), \]
which is exactly the functional behaviour of an \emph{exponential} function ($a^{x+y} = a^x a^y$). This suggests writing $f(\theta)$ as some base raised to a power proportional to $\theta$. The correct identification is one of the most famous results in all of mathematics.

\begin{definition}[Euler's formula and the exponential form]
\cle{Euler's formula} states that for every real number $\theta$,
\[ e^{i\theta} = \cos\theta + i\sin\theta. \]
Consequently every non-zero complex number can be written in \cle{exponential form}
\[ z = r\,e^{i\theta}, \qquad r = |z|,\ \ \theta = \arg(z). \]
\end{definition}

A full justification uses power series: expanding $e^{x}$, $\cos x$ and $\sin x$ as infinite series and substituting $x = i\theta$, the real terms regroup into $\cos\theta$ and the imaginary terms into $i\sin\theta$. At this level you should \emph{know and use} Euler's formula; the series proof is beyond the syllabus but the formula itself is entirely rigorous.

\begin{exemplebox}[Familiar numbers in exponential form]
\[ 1 = e^{0}, \qquad i = e^{i\pi/2}, \qquad -1 = e^{i\pi}, \qquad -i = e^{-i\pi/2}, \qquad 1 + i = \sqrt{2}\,e^{i\pi/4}. \]
Each equality is read straight off the Argand diagram: modulus first, then the angle.
\end{exemplebox}

\begin{savaistu}[The most beautiful equation in mathematics]
Setting $\theta = \pi$ in Euler's formula gives $e^{i\pi} = \cos\pi + i\sin\pi = -1$, that is,
\[ e^{i\pi} + 1 = 0. \]
This single line, called \cle{Euler's identity}, connects the five most fundamental constants of mathematics — $0$, $1$, $\pi$, $e$ and $i$ — through the three basic operations (addition, multiplication, exponentiation). It is regularly voted the most beautiful formula in mathematics. (Cultural enrichment: not examinable, but worth knowing!)
\end{savaistu}

\courssub{Computing with the Exponential Form}

The exponential form behaves exactly like ordinary powers, which is why it is so convenient. All the rules below follow immediately from $e^{i\theta} = \cos\theta + i\sin\theta$ and the product/quotient rules of the polar form.

\begin{propriete}[Rules of computation with $e^{i\theta}$]
For real $\theta, \theta_1, \theta_2$:
\begin{itemize}
\item \textbf{Product:} $e^{i\theta_1}\,e^{i\theta_2} = e^{i(\theta_1 + \theta_2)}$;
\item \textbf{Quotient:} $\dfrac{e^{i\theta_1}}{e^{i\theta_2}} = e^{i(\theta_1 - \theta_2)}$;
\item \textbf{Power:} $\left(e^{i\theta}\right)^n = e^{in\theta}$ for $n \in \mathbb{Z}$;
\item \textbf{Conjugate:} $\overline{e^{i\theta}} = e^{-i\theta}$ (since $\overline{\cos\theta + i\sin\theta} = \cos\theta - i\sin\theta = \cos(-\theta) + i\sin(-\theta)$);
\item \textbf{Inverse:} $\dfrac{1}{e^{i\theta}} = e^{-i\theta}$, and in particular $\left|e^{i\theta}\right| = 1$.
\end{itemize}
\end{propriete}

Note the last point: every number $e^{i\theta}$ lies on the \emph{unit circle}. Multiplying by $e^{i\theta}$ is a \emph{pure rotation} through $\theta$, with no enlargement.

\begin{exoresolu}[Working in exponential form]
Given $z_1 = 2e^{i\pi/4}$ and $z_2 = 3e^{-i\pi/6}$, find in exponential form: (a) $z_1 z_2$; (b) $\dfrac{z_1}{z_2}$; (c) $\overline{z_1}$; (d) $z_1^4$.
\tcblower
\textbf{(a)} Moduli multiply, arguments add:
\[ z_1 z_2 = (2 \times 3)\,e^{i(\pi/4 - \pi/6)} = 6\,e^{i\pi/12}. \]
\textbf{(b)} Moduli divide, arguments subtract:
\[ \frac{z_1}{z_2} = \frac{2}{3}\,e^{i(\pi/4 - (-\pi/6))} = \frac{2}{3}\,e^{i5\pi/12}. \]
\textbf{(c)} Conjugating negates the argument: $\overline{z_1} = 2e^{-i\pi/4}$.\\
\textbf{(d)} Powers act on both parts: $z_1^4 = 2^4 e^{i4\pi/4} = 16e^{i\pi} = -16$.
\end{exoresolu}

\courssub{Expressing $\cos\theta$ and $\sin\theta$ with Exponentials}

Euler's formula gives $e^{i\theta} = \cos\theta + i\sin\theta$. Replacing $\theta$ by $-\theta$, and using the parity of cosine and sine:
\[ e^{-i\theta} = \cos\theta - i\sin\theta. \]
We now have a system of two equations in the two unknowns $\cos\theta$ and $\sin\theta$. \textbf{Adding} the two equations eliminates the sine; \textbf{subtracting} them eliminates the cosine.

\begin{propriete}[Euler's relations for cosine and sine]
For every real $\theta$,
\[ \cos\theta = \frac{e^{i\theta} + e^{-i\theta}}{2}, \qquad \sin\theta = \frac{e^{i\theta} - e^{-i\theta}}{2i}. \]
\end{propriete}

These two identities look modest, but they are the engine behind the next section: combined with the binomial theorem, they let us express $\cos^n\theta$ and $\sin^n\theta$ as sums of multiple angles — a technique the GCE examiners love, and an essential tool for integrating powers of trigonometric functions in Further Pure Mathematics.

\begin{exemplebox}[A first taste: linearising $\cos^2\theta$]
Using $\cos\theta = \tfrac{1}{2}\left(e^{i\theta} + e^{-i\theta}\right)$:
\[ \cos^2\theta = \frac{1}{4}\left(e^{i\theta} + e^{-i\theta}\right)^2 = \frac{1}{4}\left(e^{2i\theta} + 2 + e^{-2i\theta}\right) = \frac{1}{2} + \frac{1}{2}\cdot\frac{e^{2i\theta} + e^{-2i\theta}}{2} = \frac{1 + \cos 2\theta}{2}, \]
recovering the familiar double-angle identity almost mechanically.
\end{exemplebox}

\begin{attention}[Common mistakes with the exponential form]
\begin{itemize}
\item Do not forget the modulus: $z = 2 + 2i$ is $2\sqrt{2}\,e^{i\pi/4}$, \emph{not} $2e^{i\pi/4}$. Always extract $r$ first.
\item $e^{i\theta}$ is \emph{not} the real exponential: it never grows, since $|e^{i\theta}| = 1$. Do not write $e^{i\pi} = e \cdot e^{i\pi}$-style nonsense, and never treat $i\theta$ as a real exponent when estimating size.
\item When adding arguments, reduce to the principal range if asked: $e^{i5\pi/4}$ should be rewritten $e^{-i3\pi/4}$ for the principal argument.
\item In $\sin\theta = \dfrac{e^{i\theta} - e^{-i\theta}}{2i}$, the denominator is $2i$, not $2$ — forgetting the $i$ is a classic slip.
\end{itemize}
\end{attention}

\begin{aretenir}[Key points of the section]
\begin{itemize}
\item Polar form: $z = r(\cos\theta + i\sin\theta)$ with $r = |z|$, $\theta = \arg(z)$ (defined modulo $2\pi$).
\item Products and quotients: moduli multiply/divide, arguments add/subtract.
\item Multiplication by $z$ is a rotation through $\arg(z)$ combined with an enlargement of ratio $|z|$.
\item Euler's formula: $e^{i\theta} = \cos\theta + i\sin\theta$; exponential form $z = re^{i\theta}$; Euler's identity $e^{i\pi} + 1 = 0$.
\item $\overline{e^{i\theta}} = e^{-i\theta} = \dfrac{1}{e^{i\theta}}$, and $|e^{i\theta}| = 1$.
\item $\cos\theta = \dfrac{e^{i\theta} + e^{-i\theta}}{2}$, \quad $\sin\theta = \dfrac{e^{i\theta} - e^{-i\theta}}{2i}$.
\end{itemize}
\end{aretenir}

\begin{exercice}[Practice]
\begin{enumerate}
\item Write in polar and in exponential form: (a) $1 - i$; (b) $-2\sqrt{3} + 2i$; (c) $-5$.
\item Given $z_1 = 4e^{i\pi/3}$ and $z_2 = 2e^{-i\pi/2}$, compute $z_1z_2$, $\dfrac{z_1}{z_2}$ and $z_1^3$, giving principal arguments.
\item Show that $(1 + i)^2 = 2i$ using polar form, and interpret the result geometrically.
\item Using Euler's relations, prove that $\sin^2\theta = \dfrac{1 - \cos 2\theta}{2}$.
\end{enumerate}
\end{exercice}

\begin{corrige}[Exercise 1]
\textbf{1.} (a) $r = \sqrt{2}$, $\theta = -\tfrac{\pi}{4}$ (fourth quadrant): $1 - i = \sqrt{2}\left(\cos\tfrac{\pi}{4} - i\sin\tfrac{\pi}{4}\right) = \sqrt{2}\,e^{-i\pi/4}$. (b) $r = 4$, $\theta = \tfrac{5\pi}{6}$: $-2\sqrt{3} + 2i = 4e^{i5\pi/6}$. (c) $-5 = 5e^{i\pi}$.\\
\textbf{2.} $z_1z_2 = 8e^{i(\pi/3 - \pi/2)} = 8e^{-i\pi/6}$; $\dfrac{z_1}{z_2} = 2e^{i(\pi/3 + \pi/2)} = 2e^{i5\pi/6}$; $z_1^3 = 64e^{i\pi} = -64$.\\
\textbf{3.} $1 + i = \sqrt{2}e^{i\pi/4}$, so $(1+i)^2 = 2e^{i\pi/2} = 2i$. Geometrically: rotate through $\tfrac{\pi}{4}$ and enlarge by $\sqrt{2}$, applied twice — a rotation through $\tfrac{\pi}{2}$ with enlargement $2$.\\
\textbf{4.} $\sin^2\theta = \left(\tfrac{e^{i\theta} - e^{-i\theta}}{2i}\right)^2 = \dfrac{e^{2i\theta} - 2 + e^{-2i\theta}}{-4} = \dfrac{2 - 2\cos 2\theta}{4} = \dfrac{1 - \cos 2\theta}{2}$.
\end{corrige}

We now have three languages for the same object — Cartesian ($a + ib$), polar ($r(\cos\theta + i\sin\theta)$) and exponential ($re^{i\theta}$) — and we can switch freely between them. In the next section we exploit the power rule $\left(e^{i\theta}\right)^n = e^{in\theta}$ to obtain \emph{De Moivre's theorem} and its spectacular trigonometric applications.

\coursec{De Moivre's Theorem and Trigonometric Applications}

In the previous section we saw that a complex number $z = r(\cos\theta + i\sin\theta)$ can be raised to a power $n$ by simply raising the modulus to $n$ and multiplying the argument by $n$. This is the essence of \cle{De Moivre's Theorem}, a result that bridges algebra and trigonometry and provides powerful tools for simplifying powers of complex numbers, deriving multiple-angle formulae, and evaluating exact trigonometric values.

\courssub{De Moivre's Theorem}

\begin{propriete}[De Moivre's Theorem for Integer Powers]
For any integer $n$ and any real angle $\theta$,
\[
(\cos\theta + i\sin\theta)^n = \cos(n\theta) + i\sin(n\theta).
\]
Equivalently, in exponential form: $(e^{i\theta})^n = e^{in\theta}$.
\end{propriete}

\begin{experience}[Discovering the Pattern]
Let us compute the first few powers of $z = \cos\theta + i\sin\theta$ manually:
\begin{align*}
z^1 &= \cos\theta + i\sin\theta,\\
z^2 &= (\cos\theta + i\sin\theta)^2 = \cos^2\theta - \sin^2\theta + 2i\sin\theta\cos\theta = \cos2\theta + i\sin2\theta,\\
z^3 &= z^2 \cdot z = (\cos2\theta + i\sin2\theta)(\cos\theta + i\sin\theta) = \cos3\theta + i\sin3\theta.
\end{align*}
The pattern is clear: the argument adds up. We prove this for all non-negative integers by induction.
\end{experience}

\textbf{Proof by induction (for $n \ge 0$).}
\begin{itemize}
\item \textbf{Base case $n=0$:} $(\cos\theta + i\sin\theta)^0 = 1 = \cos0 + i\sin0$. True.
\item \textbf{Base case $n=1$:} Trivially true.
\item \textbf{Inductive step:} Assume true for $n=k$, i.e. $(\cos\theta + i\sin\theta)^k = \cos(k\theta) + i\sin(k\theta)$.
Then for $n=k+1$:
\begin{align*}
(\cos\theta + i\sin\theta)^{k+1} &= (\cos\theta + i\sin\theta)^k (\cos\theta + i\sin\theta)\\
&= [\cos(k\theta) + i\sin(k\theta)](\cos\theta + i\sin\theta)\\
&= \cos(k\theta)\cos\theta - \sin(k\theta)\sin\theta + i[\sin(k\theta)\cos\theta + \cos(k\theta)\sin\theta]\\
&= \cos((k+1)\theta) + i\sin((k+1)\theta).
\end{align*}
Hence true for $k+1$. By induction, the theorem holds for all $n \in \mathbb{N}$.
\end{itemize}

\textbf{Extension to negative integers.}
Let $n = -m$ where $m \in \mathbb{N}$. Then
\[
(\cos\theta + i\sin\theta)^{-m} = \frac{1}{(\cos\theta + i\sin\theta)^m} = \frac{1}{\cos(m\theta) + i\sin(m\theta)}.
\]
Multiply numerator and denominator by the conjugate $\cos(m\theta) - i\sin(m\theta)$:
\[
\frac{\cos(m\theta) - i\sin(m\theta)}{\cos^2(m\theta)+\sin^2(m\theta)} = \cos(m\theta) - i\sin(m\theta) = \cos(-m\theta) + i\sin(-m\theta).
\]
Thus the formula holds for all $n \in \mathbb{Z}$.

\begin{popfigure}
\begin{tikzpicture}[scale=1.5]
\draw[->] (-1.5,0) -- (1.5,0) node[right] {$\Re$};
\draw[->] (0,-1.5) -- (0,1.5) node[above] {$\Im$};
\draw (0,0) circle (1);
\foreach \ang/\lbl in {30/$\theta$, 60/$2\theta$, 90/$3\theta$} {
  \draw[popblue, thick] (0,0) -- (\ang:1);
  \node[popblue] at (\ang:1.15) {\lbl};
}
\draw[poporange, thick, ->] (0.3,0) arc (0:30:0.3) node[midway, right] {$\theta$};
\draw[poporange, thick, ->] (0.45,0) arc (0:60:0.45) node[midway, right] {$2\theta$};
\draw[poporange, thick, ->] (0.6,0) arc (0:90:0.6) node[midway, right] {$3\theta$};
\node at (0,-1.3) {Unit circle: successive powers rotate by $\theta$};
\end{tikzpicture}
\legende{Geometric interpretation of De Moivre's theorem: multiplying by $\cos\theta+i\sin\theta$ rotates by $\theta$.}
\end{popfigure}

\begin{attention}[Common Mistake: Sign in the Sine Term]
De Moivre's theorem applies to $\cos\theta + i\sin\theta$. If you have $\cos\theta - i\sin\theta$, rewrite it as $\cos(-\theta) + i\sin(-\theta)$ before applying the theorem:
\[
(\cos\theta - i\sin\theta)^n = (\cos(-\theta) + i\sin(-\theta))^n = \cos(-n\theta) + i\sin(-n\theta) = \cos(n\theta) - i\sin(n\theta).
\]
Do \emph{not} write $(\cos\theta - i\sin\theta)^n = \cos(n\theta) - i\sin(n\theta)$ directly without justification — it is a trap in GCE exams!
\end{attention}

\courssub{Applications: Powers of Complex Numbers}

The most direct application is computing large powers of a complex number by converting to polar form.

\begin{methode}[Computing $z^n$ using De Moivre]
\begin{enumerate}
\item Write $z$ in polar form $z = r(\cos\theta + i\sin\theta)$ (or $re^{i\theta}$).
\item Apply De Moivre: $z^n = r^n(\cos n\theta + i\sin n\theta)$.
\item Convert back to Cartesian form $a+ib$ if required, using exact values of $\cos n\theta$ and $\sin n\theta$.
\end{enumerate}
\end{methode}

\begin{exemplebox}[Power of a Complex Number]
Compute $(1+i)^{12}$ and express the result in the form $a+ib$.
\end{exemplebox}
\textbf{Solution.}
$1+i$ has modulus $r = \sqrt{2}$ and argument $\theta = \frac{\pi}{4}$.
\[
(1+i)^{12} = (\sqrt{2})^{12} \left(\cos\frac{12\pi}{4} + i\sin\frac{12\pi}{4}\right) = 2^6 (\cos3\pi + i\sin3\pi) = 64(-1 + i\cdot0) = -64.
\]

\begin{exoresolu}[Finding $n$ such that a Power is Real or Purely Imaginary]
Find the smallest positive integer $n$ such that $(\sqrt{3}+i)^n$ is (a) real, (b) purely imaginary.
\tcblower
\textbf{Solution.}
$\sqrt{3}+i$: modulus $r=2$, argument $\theta = \frac{\pi}{6}$.
$(\sqrt{3}+i)^n = 2^n \left(\cos\frac{n\pi}{6} + i\sin\frac{n\pi}{6}\right)$.
\begin{itemize}
\item[(a)] Real $\iff \sin\frac{n\pi}{6} = 0 \iff \frac{n\pi}{6} = k\pi \iff n = 6k$. Smallest positive $n = 6$.
\item[(b)] Purely imaginary $\iff \cos\frac{n\pi}{6} = 0 \iff \frac{n\pi}{6} = \frac{\pi}{2}+k\pi \iff n = 3+6k$. Smallest positive $n = 3$.
\end{itemize}
\end{exoresolu}

\courssub{Linearisation: Powers of Cosine and Sine as Multiple Angles}

A classic GCE technique uses the substitution $z = \cos\theta + i\sin\theta$. Then
\[
z + \frac{1}{z} = 2\cos\theta, \qquad z - \frac{1}{z} = 2i\sin\theta.
\]
Powers of $\cos\theta$ and $\sin\theta$ can be expressed as linear combinations of $\cos k\theta$ and $\sin k\theta$ by expanding $(z \pm z^{-1})^n$ using the binomial theorem and grouping terms.

\begin{methode}[Linearising $\cos^n\theta$ and $\sin^n\theta$]
\begin{enumerate}
\item Set $z = \cos\theta + i\sin\theta$. Then $\cos\theta = \frac{1}{2}(z+z^{-1})$, $\sin\theta = \frac{1}{2i}(z-z^{-1})$.
\item Raise to the required power $n$ and expand using the binomial theorem.
\item Group terms $z^k + z^{-k} = 2\cos(k\theta)$ and $z^k - z^{-k} = 2i\sin(k\theta)$.
\item Simplify to obtain an expression in $\cos k\theta$ or $\sin k\theta$.
\end{enumerate}
\end{methode}

\begin{exemplebox}[Linearising $\cos^4\theta$]
Express $\cos^4\theta$ in terms of multiple angles.
\end{exemplebox}
\textbf{Solution.}
$\cos\theta = \frac{1}{2}(z+z^{-1})$.
\[
\cos^4\theta = \frac{1}{16}(z+z^{-1})^4 = \frac{1}{16}\left(z^4 + 4z^2 + 6 + 4z^{-2} + z^{-4}\right).
\]
Group: $z^4+z^{-4} = 2\cos4\theta$, $z^2+z^{-2} = 2\cos2\theta$.
\[
\cos^4\theta = \frac{1}{16}\left(2\cos4\theta + 4\cdot2\cos2\theta + 6\right) = \frac{1}{8}\cos4\theta + \frac{1}{2}\cos2\theta + \frac{3}{8}.
\]

\begin{exoresolu}[Linearising $\sin^5\theta$]
Express $\sin^5\theta$ in terms of sines of multiple angles.
\tcblower
\textbf{Solution.}
$\sin\theta = \frac{1}{2i}(z-z^{-1})$.
\[
\sin^5\theta = \frac{1}{(2i)^5}(z-z^{-1})^5 = \frac{1}{32i}\left(z^5 - 5z^3 + 10z - 10z^{-1} + 5z^{-3} - z^{-5}\right).
\]
Group: $z^5 - z^{-5} = 2i\sin5\theta$, $z^3 - z^{-3} = 2i\sin3\theta$, $z - z^{-1} = 2i\sin\theta$.
\[
\sin^5\theta = \frac{1}{32i}\left(2i\sin5\theta - 5\cdot2i\sin3\theta + 10\cdot2i\sin\theta\right) = \frac{1}{16}\sin5\theta - \frac{5}{16}\sin3\theta + \frac{5}{8}\sin\theta.
\]
\end{exoresolu}

\begin{popfigure}
\begin{tikzpicture}
\begin{axis}[
    width=\linewidth, height=6cm,
    axis lines=middle,
    xlabel={$x$}, ylabel={$y$},
    domain=-3.14:3.14, samples=200,
    xmin=-3.5, xmax=3.5, ymin=-1.2, ymax=1.2,
    xtick={-3.1415,-1.5708,0,1.5708,3.1415},
    xticklabels={$-\pi$,$-\pi/2$,$0$,$\pi/2$,$\pi$},
    legend style={at={(0.98,0.98)}, anchor=north east, font=\small},
    grid=major, grid style={popline!30}
]
\addplot[popblue, thick] {cos(x)^3};
\addplot[poporange, dashed, thick] {(3*cos(x) + cos(3*x))/4};
\legend{$\cos^3 x$, $\frac{1}{4}(3\cos x + \cos 3x)$};
\end{axis}
\end{tikzpicture}
\legende{Graphical verification: $\cos^3 x = \frac{1}{4}(3\cos x + \cos 3x)$. The curves coincide perfectly.}
\end{popfigure}

\courssub{Expansion: Multiple Angles as Powers of Cosine and Sine}

The reverse process expands $\cos n\theta$ and $\sin n\theta$ as polynomials in $\cos\theta$ and $\sin\theta$ using the binomial theorem on $(\cos\theta + i\sin\theta)^n$.

\begin{methode}[Expanding $\cos n\theta$ and $\sin n\theta$]
\begin{enumerate}
\item Write $(\cos\theta + i\sin\theta)^n = \cos n\theta + i\sin n\theta$ (De Moivre).
\item Expand the left‑hand side using the binomial theorem:
  \[
  \sum_{k=0}^n \binom{n}{k} \cos^{n-k}\theta \,(i\sin\theta)^k.
  \]
\item Separate real and imaginary parts:
  \begin{itemize}
  \item Real part (even $k$) gives $\cos n\theta$.
  \item Imaginary part (odd $k$) gives $\sin n\theta$.
  \end{itemize}
\item Optionally replace $\sin^2\theta = 1-\cos^2\theta$ to express $\cos n\theta$ solely in powers of $\cos\theta$ (Chebyshev polynomials).
\end{enumerate}
\end{methode}

\begin{exemplebox}[Expanding $\cos3\theta$ and $\sin3\theta$]
Find $\cos3\theta$ and $\sin3\theta$ in terms of $\cos\theta$ and $\sin\theta$.
\end{exemplebox}
\textbf{Solution.}
$(\cos\theta + i\sin\theta)^3 = \cos^3\theta + 3i\cos^2\theta\sin\theta - 3\cos\theta\sin^2\theta - i\sin^3\theta$.
\begin{align*}
\cos3\theta &= \cos^3\theta - 3\cos\theta\sin^2\theta = \cos^3\theta - 3\cos\theta(1-\cos^2\theta) = 4\cos^3\theta - 3\cos\theta.\\
\sin3\theta &= 3\cos^2\theta\sin\theta - \sin^3\theta = 3(1-\sin^2\theta)\sin\theta - \sin^3\theta = 3\sin\theta - 4\sin^3\theta.
\end{align*}

\begin{exoresolu}[Expanding $\cos4\theta$ and $\sin4\theta$]
Express $\cos4\theta$ and $\sin4\theta$ as polynomials in $\cos\theta$ and $\sin\theta$.
\tcblower
\textbf{Solution.}
$(\cos\theta + i\sin\theta)^4 = \cos^4\theta + 4i\cos^3\theta\sin\theta - 6\cos^2\theta\sin^2\theta - 4i\cos\theta\sin^3\theta + \sin^4\theta$.
\begin{align*}
\cos4\theta &= \cos^4\theta - 6\cos^2\theta\sin^2\theta + \sin^4\theta\\
&= \cos^4\theta - 6\cos^2\theta(1-\cos^2\theta) + (1-\cos^2\theta)^2\\
&= 8\cos^4\theta - 8\cos^2\theta + 1.\\
\sin4\theta &= 4\cos^3\theta\sin\theta - 4\cos\theta\sin^3\theta = 4\cos\theta\sin\theta(\cos^2\theta - \sin^2\theta)\\
&= 4\cos\theta\sin\theta(1-2\sin^2\theta) = 4\sin\theta\cos\theta - 8\sin^3\theta\cos\theta.
\end{align*}
\end{exoresolu}

\courssub{Trigonometric Identities and Exact Values}

De Moivre's theorem is a favourite tool in GCE exams for proving identities and finding exact values of trigonometric ratios.

\begin{exoresolu}[Proving an Identity]
Prove that $\tan5\theta = \frac{5\tan\theta - 10\tan^3\theta + \tan^5\theta}{1 - 10\tan^2\theta + 5\tan^4\theta}$.
\tcblower
\textbf{Solution.}
By De Moivre: $\cos5\theta + i\sin5\theta = (\cos\theta + i\sin\theta)^5$.
Expand using binomial theorem:
\begin{align*}
\cos5\theta &= \cos^5\theta - 10\cos^3\theta\sin^2\theta + 5\cos\theta\sin^4\theta,\\
\sin5\theta &= 5\cos^4\theta\sin\theta - 10\cos^2\theta\sin^3\theta + \sin^5\theta.
\end{align*}
Divide numerator and denominator by $\cos^5\theta$:
\[
\tan5\theta = \frac{\sin5\theta}{\cos5\theta} = \frac{5\tan\theta - 10\tan^3\theta + \tan^5\theta}{1 - 10\tan^2\theta + 5\tan^4\theta}.
\]
\end{exoresolu}

\begin{exoresolu}[Exact Value: $\cos\frac{\pi}{12}$]
Use De Moivre's theorem to find the exact value of $\cos\frac{\pi}{12}$.
\tcblower
\textbf{Solution.}
We know $\cos\frac{\pi}{6} = \frac{\sqrt{3}}{2}$. Using the half‑angle formula derived from $\cos2\theta = 2\cos^2\theta-1$:
\[
\cos\frac{\pi}{12} = \sqrt{\frac{1+\cos\frac{\pi}{6}}{2}} = \sqrt{\frac{1+\frac{\sqrt{3}}{2}}{2}} = \sqrt{\frac{2+\sqrt{3}}{4}} = \frac{\sqrt{2+\sqrt{3}}}{2}.
\]
Alternatively, using $\cos3\theta = 4\cos^3\theta-3\cos\theta$ with $\theta=\frac{\pi}{12}$ gives a cubic equation, but the half‑angle method is simpler. The exact value is $\frac{\sqrt{6}+\sqrt{2}}{4}$ (after rationalising $\sqrt{2+\sqrt{3}}$).
\end{exoresolu}

\begin{aretenir}[Key Results for GCE]
\begin{itemize}
\item De Moivre: $(\cos\theta + i\sin\theta)^n = \cos n\theta + i\sin n\theta$ for $n\in\mathbb{Z}$.
\item $z = \cos\theta + i\sin\theta \;\Rightarrow\; z+z^{-1}=2\cos\theta,\; z-z^{-1}=2i\sin\theta$.
\item Linearisation: expand $(z\pm z^{-1})^n$ to write $\cos^n\theta$ or $\sin^n\theta$ as sums of $\cos k\theta$ or $\sin k\theta$.
\item Expansion: expand $(\cos\theta + i\sin\theta)^n$ to write $\cos n\theta$ and $\sin n\theta$ as polynomials in $\cos\theta$, $\sin\theta$.
\item To find exact values: use known angles ($\frac{\pi}{6},\frac{\pi}{4},\frac{\pi}{3}$) and half‑angle / multiple‑angle formulae derived from De Moivre.
\end{itemize}
\end{aretenir}

\begin{attention}[Examination Traps]
\begin{itemize}
\item Forgetting that De Moivre applies to \emph{integer} powers only. Fractional powers require the $n$-th roots formula (next section).
\item When linearising $\sin^n\theta$ for odd $n$, the result is a sum of sines; for even $n$, a sum of cosines plus a constant.
\item When expanding $\cos n\theta$, always check whether the question asks for the answer in terms of $\cos\theta$ only (replace $\sin^2\theta = 1-\cos^2\theta$).
\item In proving identities, show the separation of real and imaginary parts clearly; do not skip steps.
\end{itemize}
\end{attention}

\begin{savaistu}[De Moivre and the History of Trigonometry]
Abraham de Moivre (1667--1754), a French Huguenot mathematician who fled to England, discovered this theorem around 1707. It was later popularised by Euler, who gave the elegant exponential proof $e^{in\theta}=(e^{i\theta})^n$. De Moivre also pioneered the normal approximation to the binomial distribution --- a link between complex numbers and probability!
\end{savaistu}

\begin{exercice}[Practice Problems]
\begin{enumerate}
\item Use De Moivre's theorem to simplify $(1-i\sqrt{3})^6$.
\item Find the smallest positive integer $n$ such that $\left(\frac{1+i}{\sqrt{2}}\right)^n$ is real.
\item Express $\sin^4\theta$ in terms of $\cos2\theta$ and $\cos4\theta$.
\item Express $\cos5\theta$ as a polynomial in $\cos\theta$.
\item Prove that $\tan3\theta = \frac{3\tan\theta - \tan^3\theta}{1-3\tan^2\theta}$.
\item Hence find the exact value of $\tan\frac{\pi}{12}$.
\end{enumerate}
\end{exercice}

\begin{corrige}[Exercise 1]
$1-i\sqrt{3}$: $r=2$, $\theta=-\frac{\pi}{3}$. $(1-i\sqrt{3})^6 = 2^6(\cos(-2\pi)+i\sin(-2\pi)) = 64$.
\end{corrige}
\begin{corrige}[Exercise 2]
$\frac{1+i}{\sqrt{2}} = \cos\frac{\pi}{4}+i\sin\frac{\pi}{4}$. We need $\sin\frac{n\pi}{4}=0 \Rightarrow \frac{n\pi}{4}=k\pi \Rightarrow n=4k$. Smallest $n=4$.
\end{corrige}
\begin{corrige}[Exercise 3]
$\sin^4\theta = \frac{1}{16}(z-z^{-1})^4 = \frac{1}{16}(z^4 - 4z^2 + 6 - 4z^{-2} + z^{-4}) = \frac{1}{8}\cos4\theta - \frac{1}{2}\cos2\theta + \frac{3}{8}$.
\end{corrige}
\begin{corrige}[Exercise 4]
$\cos5\theta = 16\cos^5\theta - 20\cos^3\theta + 5\cos\theta$ (expand $(\cos\theta+i\sin\theta)^5$ and take real part).
\end{corrige}
\begin{corrige}[Exercise 5]
$\tan3\theta = \frac{\sin3\theta}{\cos3\theta} = \frac{3\sin\theta\cos^2\theta-\sin^3\theta}{\cos^3\theta-3\cos\theta\sin^2\theta}$. Divide by $\cos^3\theta$: $\frac{3\tan\theta-\tan^3\theta}{1-3\tan^2\theta}$.
\end{corrige}
\begin{corrige}[Exercise 6]
Let $\theta=\frac{\pi}{12}$, then $3\theta=\frac{\pi}{4}$, $\tan\frac{\pi}{4}=1$. Solve $\frac{3t-t^3}{1-3t^2}=1$ where $t=\tan\frac{\pi}{12}$. $t^3-3t^2-3t+1=0 \Rightarrow (t+1)(t^2-4t+1)=0$. Since $t>0$, $t=2-\sqrt{3}$.
\end{corrige}

\coursec{nth Roots of Complex Numbers and Cube Roots of Unity}

In the previous section, De Moivre's theorem gave us a powerful way to \emph{raise} a complex number to a power: $(r e^{i\theta})^n = r^n e^{in\theta}$. We now turn the question around. Instead of computing a power, we want to \emph{undo} a power: given a complex number $w$, can we find all the numbers $z$ such that $z^n = w$? You already know the case $n = 2$: the equation $z^2 = -16$ has the two solutions $z = 4i$ and $z = -4i$. De Moivre's theorem will show us that this is no accident: \textbf{every} non-zero complex number has \textbf{exactly $n$ distinct $n$th roots}, arranged with perfect symmetry in the Argand diagram. This beautiful result closes the story begun in Lesson 29 and gives us, as a special case, the celebrated cube roots of unity, a favourite of GCE examiners.

\courssub{Definition of the nth Roots of a Complex Number}

\textbf{Intuition first.}\par
Over the real numbers, the equation $x^3 = 8$ has only one solution, $x = 2$. But over the complex numbers, could there be others? Think of the equation $z^3 = 1$. Certainly $z = 1$ works. But perhaps a complex number whose cube ``rotates back'' to $1$ also works: for instance a number of modulus $1$ and argument $\tfrac{2\pi}{3}$, since cubing triples the argument: $3 \times \tfrac{2\pi}{3} = 2\pi$, which is the same direction as argument $0$. So there may be \emph{three} cube roots of $1$, not one. Let us make this precise.

\begin{definition}[nth root of a complex number]
Let $n \in \mathbb{N}$, $n \geq 2$, and let $w \in \mathbb{C}$. An \cle{nth root} of $w$ is any complex number $z$ such that
\[ z^n = w. \]
The $n$th roots of $1$ are called the \cle{nth roots of unity}.
\end{definition}

To find all of them, write both numbers in exponential form. Let $w = R e^{i\Theta}$ with $R > 0$, and look for $z = r e^{i\theta}$. Then
\[ z^n = w \iff r^n e^{in\theta} = R e^{i\Theta}. \]
Two complex numbers in exponential form are equal exactly when their moduli are equal and their arguments differ by a multiple of $2\pi$. Hence
\[ r^n = R \quad\text{and}\quad n\theta = \Theta + 2k\pi, \quad k \in \mathbb{Z}. \]
Since $r$ is a positive real number, $r = R^{1/n}$ (the ordinary real $n$th root), and
\[ \theta = \frac{\Theta + 2k\pi}{n}, \qquad k \in \mathbb{Z}. \]
Taking $k = 0, 1, 2, \dots, n-1$ gives $n$ different values of $\theta$ in $[0, 2\pi)$, hence $n$ distinct roots. Any other integer $k$ differs from one of these by a multiple of $n$, and adding $n$ to $k$ adds $2\pi$ to $\theta$, which gives the \emph{same} complex number. So there are no others. We have proved:

\begin{propriete}[Theorem — number and form of the nth roots]
Every non-zero complex number $w = R e^{i\Theta}$ has \textbf{exactly $n$ distinct $n$th roots}, given by
\[ z_k = R^{1/n}\, e^{\,i\frac{\Theta + 2k\pi}{n}}, \qquad k = 0, 1, 2, \dots, n-1. \]
In particular, the $n$th roots of unity are $e^{2ik\pi/n}$ for $k = 0, 1, \dots, n-1$. \emph{(The derivation above is examinable.)}
\end{propriete}

\begin{propriete}[Geometric property — the regular n-gon]
All $n$ roots have the same modulus $R^{1/n}$, so they lie on the \textbf{circle of centre $O$ and radius $R^{1/n}$}. Consecutive roots differ in argument by exactly $\dfrac{2\pi}{n}$, so they are equally spaced around this circle: they are the \textbf{vertices of a regular $n$-gon} inscribed in that circle.
\end{propriete}

\begin{methode}[Solving $z^n = w$]
\begin{enumerate}
\item Write $w$ in exponential form: find $R = |w|$ and an argument $\Theta$.
\item Write the general root $z_k = R^{1/n} e^{i(\Theta + 2k\pi)/n}$.
\item Substitute $k = 0, 1, \dots, n-1$ to list the $n$ roots.
\item If required, convert each root to Cartesian form $a + ib$.
\item Check the count: you must have exactly $n$ distinct answers.
\end{enumerate}
\end{methode}

\begin{exoresolu}[Solving $z^4 = -16$]
Solve in $\mathbb{C}$ the equation $z^4 = -16$, and represent the solutions in an Argand diagram.
\tcblower
\textbf{Solution.} Write $-16$ in exponential form: $|-16| = 16$ and $\arg(-16) = \pi$, so $-16 = 16 e^{i\pi}$. The four roots are
\[ z_k = 16^{1/4}\, e^{\,i\frac{\pi + 2k\pi}{4}} = 2\, e^{\,i\frac{(2k+1)\pi}{4}}, \qquad k = 0, 1, 2, 3. \]
This gives:
\begin{align*}
z_0 &= 2 e^{i\pi/4} = 2\left(\tfrac{\sqrt{2}}{2} + i\tfrac{\sqrt{2}}{2}\right) = \sqrt{2} + i\sqrt{2},\\
z_1 &= 2 e^{3i\pi/4} = -\sqrt{2} + i\sqrt{2},\\
z_2 &= 2 e^{5i\pi/4} = -\sqrt{2} - i\sqrt{2},\\
z_3 &= 2 e^{7i\pi/4} = \sqrt{2} - i\sqrt{2}.
\end{align*}
The four solutions are $\pm\sqrt{2} \pm i\sqrt{2}$ (all four sign combinations). They lie on the circle of radius $2$, at the vertices of a square. \textbf{Check:} $z_0^2 = (\sqrt{2}+i\sqrt{2})^2 = 2 + 4i - 2 = 4i$, so $z_0^4 = (4i)^2 = -16$. $\checkmark$
\end{exoresolu}

\begin{attention}[There are always n roots!]
The most common error in the GCE is to give only \emph{one} root — usually the ``obvious'' one — and stop. The equation $z^n = w$ (with $w \neq 0$) has \textbf{exactly $n$ distinct solutions}. If you found fewer, you have lost marks. A quick safety net: count your answers and check that consecutive arguments differ by $\tfrac{2\pi}{n}$.
\end{attention}

\courssub{The Cube Roots of Unity}

The case $n = 3$, $w = 1$ is so important that it deserves its own study. The cube roots of unity are the solutions of $z^3 = 1$, i.e.
\[ z_k = e^{2ik\pi/3}, \qquad k = 0, 1, 2. \]
Explicitly:
\[ z_0 = 1, \qquad z_1 = e^{2i\pi/3} = -\frac{1}{2} + i\frac{\sqrt{3}}{2}, \qquad z_2 = e^{4i\pi/3} = -\frac{1}{2} - i\frac{\sqrt{3}}{2}. \]

\begin{definition}[The number omega]
We write
\[ \omega = e^{2i\pi/3} = \frac{-1 + i\sqrt{3}}{2}. \]
The three cube roots of unity are then $1$, $\omega$ and $\omega^2$ (since $\omega^2 = e^{4i\pi/3} = z_2$).
\end{definition}

\begin{propriete}[Properties of omega]
The number $\omega$ satisfies:
\begin{enumerate}
\item $\omega^3 = 1$ \quad (it is a cube root of $1$);
\item $1 + \omega + \omega^2 = 0$;
\item $\bar{\omega} = \omega^2$ \quad and \quad $\omega \cdot \bar{\omega} = |\omega|^2 = 1$.
\end{enumerate}
\emph{Proof of (2):} since $\omega^3 = 1$, we have $\omega^3 - 1 = 0$, i.e. $(\omega - 1)(\omega^2 + \omega + 1) = 0$. As $\omega \neq 1$, we conclude $1 + \omega + \omega^2 = 0$. Property (3) follows because $\omega^2 = e^{-2i\pi/3} = \overline{e^{2i\pi/3}}$.
\end{propriete}

\begin{popfigure}
\begin{center}
\begin{tikzpicture}[scale=2.2]
\draw[->, popmuted] (-1.5,0) -- (1.6,0) node[below left, popink] {Re};
\draw[->, popmuted] (0,-1.4) -- (0,1.5) node[below left, popink] {Im};
\draw[popline] (0,0) circle (1);
\draw[popblue, thick] (1,0) -- (-0.5,0.866) -- (-0.5,-0.866) -- cycle;
\draw[->, poporange, thick] (0,0) -- (1,0) node[midway, below, popink] {$1$};
\draw[->, poporange, thick] (0,0) -- (-0.5,0.866) node[midway, above left, popink] {$\omega$};
\draw[->, poporange, thick] (0,0) -- (-0.5,-0.866) node[midway, below left, popink] {$\omega^2$};
\fill[popdark] (1,0) circle (0.03);
\fill[popdark] (-0.5,0.866) circle (0.03);
\fill[popdark] (-0.5,-0.866) circle (0.03);
\draw[popgreen, ->] (0.45,0) arc (0:120:0.45) node[midway, above right, popgreen] {$\frac{2\pi}{3}$};
\end{tikzpicture}
\end{center}
\legende{The cube roots of unity $1$, $\omega$, $\omega^2$ on the unit circle: vertices of an equilateral triangle.}
\end{popfigure}

\textbf{Factorising with omega.}\par
The identity $1 + \omega + \omega^2 = 0$ gives elegant factorisations. Since $x^3 - 1 = (x-1)(x^2 + x + 1)$ and the roots of $x^2 + x + 1$ are exactly $\omega$ and $\omega^2$, we obtain, for any complex numbers $a$ and $b$:
\[ a^3 - b^3 = (a - b)(a - b\omega)(a - b\omega^2), \qquad a^3 + b^3 = (a + b)(a + b\omega)(a + b\omega^2). \]
Indeed, the roots of $z^3 = b^3$ are $b$, $b\omega$, $b\omega^2$ (see below), so $z^3 - b^3 = (z-b)(z - b\omega)(z - b\omega^2)$; setting $z = a$ gives the first identity, and replacing $b$ by $-b$ gives the second.

\begin{exemplebox}[Using omega to simplify]
Compute $(1 + \omega)(1 + \omega^2)$ and $\omega^{2024}$.
\emph{Solution.} $(1+\omega)(1+\omega^2) = 1 + \omega + \omega^2 + \omega^3 = 0 + 1 = 1$. For the power: $2024 = 3 \times 674 + 2$, so $\omega^{2024} = (\omega^3)^{674}\,\omega^2 = \omega^2$.
\end{exemplebox}

\courssub{Cube Roots of Any Complex Number via Omega}

Here is a labour-saving fact. Suppose you have found \emph{one} cube root $u$ of a number $w$ (so $u^3 = w$). Then
\[ (u\omega)^3 = u^3 \omega^3 = w \cdot 1 = w, \qquad (u\omega^2)^3 = u^3 \omega^6 = w. \]
Since a non-zero number has exactly three cube roots, they must be precisely $u$, $u\omega$, $u\omega^2$.

\begin{propriete}[From one cube root to all three]
If $u^3 = w$ with $w \neq 0$, then the three cube roots of $w$ are
\[ u, \qquad u\omega, \qquad u\omega^2. \]
More generally, if $u$ is \emph{one} $n$th root of $w$, then all $n$ roots are $u, u\zeta, u\zeta^2, \dots, u\zeta^{n-1}$ where $\zeta = e^{2i\pi/n}$.
\end{propriete}

\begin{exoresolu}[Cube roots of $-8$ and of $i$]
(a) Find the three cube roots of $-8$. \quad (b) Find the three cube roots of $i$.
\tcblower
\textbf{(a)} One root is obvious: $u = -2$, since $(-2)^3 = -8$. The three roots are therefore
\[ -2, \qquad -2\omega = -2\cdot\frac{-1+i\sqrt{3}}{2} = 1 - i\sqrt{3}, \qquad -2\omega^2 = 1 + i\sqrt{3}. \]
\textbf{(b)} Write $i = e^{i\pi/2}$. Then $|i|^{1/3} = 1$ and the roots are $z_k = e^{i(\pi/2 + 2k\pi)/3}$, $k = 0, 1, 2$:
\[ z_0 = e^{i\pi/6} = \frac{\sqrt{3}}{2} + \frac{1}{2}i, \qquad z_1 = e^{5i\pi/6} = -\frac{\sqrt{3}}{2} + \frac{1}{2}i, \qquad z_2 = e^{3i\pi/2} = -i. \]
Check with $\omega$: $z_1 = z_0 \cdot e^{2i\pi/3} = z_0\omega$ and $z_2 = z_0\omega^2$, as expected.
\end{exoresolu}

\courssub{The General Picture: nth Roots of Unity}

Everything we did for $n = 3$ works for any $n$. The $n$th roots of unity are
\[ \zeta_k = e^{2ik\pi/n}, \qquad k = 0, 1, \dots, n-1, \]
the vertices of a regular $n$-gon inscribed in the unit circle, with one vertex at $1$. They satisfy $\zeta_k = \zeta_1^k$ and, summing the geometric progression,
\[ 1 + \zeta_1 + \zeta_1^2 + \dots + \zeta_1^{n-1} = \frac{\zeta_1^n - 1}{\zeta_1 - 1} = 0, \]
so \textbf{the sum of all $n$th roots of unity is always zero} (for $n \geq 2$).

\begin{popfigure}
\begin{center}
\begin{tikzpicture}[scale=2.2]
\draw[->, popmuted] (-1.5,0) -- (1.6,0) node[below left, popink] {Re};
\draw[->, popmuted] (0,-1.4) -- (0,1.5) node[below left, popink] {Im};
\draw[popline] (0,0) circle (1);
\draw[popblue, thick] (1,0) -- (0.5,0.866) -- (-0.5,0.866) -- (-1,0) -- (-0.5,-0.866) -- (0.5,-0.866) -- cycle;
\foreach \a/\n in {0/1, 60/{\zeta}, 120/{\zeta^2}, 180/{\zeta^3}, 240/{\zeta^4}, 300/{\zeta^5}} {
  \fill[popdark] (\a:1) circle (0.03);
  \node[popink] at (\a:1.22) {$\n$};
}
\draw[popgreen, ->] (0.4,0) arc (0:60:0.4) node[midway, right, popgreen] {$\frac{\pi}{3}$};
\end{tikzpicture}
\end{center}
\legende{The six 6th roots of unity, $\zeta = e^{i\pi/3}$: vertices of a regular hexagon on the unit circle.}
\end{popfigure}

\begin{exoresolu}[Solving $z^6 = 1$ and $z^6 = -64$]
(a) Solve $z^6 = 1$. \quad (b) Solve $z^6 = -64$.
\tcblower
\textbf{(a)} The roots are $z_k = e^{2ik\pi/6} = e^{ik\pi/3}$, $k = 0, \dots, 5$:
\[ 1,\quad \frac{1}{2} + i\frac{\sqrt{3}}{2},\quad -\frac{1}{2} + i\frac{\sqrt{3}}{2},\quad -1,\quad -\frac{1}{2} - i\frac{\sqrt{3}}{2},\quad \frac{1}{2} - i\frac{\sqrt{3}}{2}. \]
\textbf{(b)} Write $-64 = 64 e^{i\pi}$. Then $64^{1/6} = 2$ and
\[ z_k = 2\, e^{\,i\frac{\pi + 2k\pi}{6}} = 2\, e^{\,i\frac{(2k+1)\pi}{6}}, \qquad k = 0, \dots, 5, \]
giving $z_0 = 2e^{i\pi/6} = \sqrt{3} + i$, $z_1 = 2e^{i\pi/2} = 2i$, $z_2 = 2e^{5i\pi/6} = -\sqrt{3} + i$, $z_3 = -\sqrt{3} - i$, $z_4 = -2i$, $z_5 = \sqrt{3} - i$. Note how (b) is just (a) rotated and scaled: the hexagon is preserved.
\end{exoresolu}

\begin{savaistu}[Roots of unity everywhere]
Roots of unity are not just examination material. They sit at the heart of the \emph{Fast Fourier Transform}, the algorithm that compresses your music and images (MP3, JPEG) and that powers digital signal processing in the mobile phones used every day across Cameroon. The same regular polygons you draw in the Argand diagram are working silently inside your phone.
\end{savaistu}

\begin{attention}[Pitfalls with nth roots]
\begin{itemize}
\item Do not write $\sqrt[n]{w}$ as if it were a single number: in $\mathbb{C}$ there are $n$ values, and the symbol is ambiguous. Always list all $n$ roots.
\item The exponent is $\dfrac{\Theta + 2k\pi}{n}$, \textbf{not} $\dfrac{\Theta}{n} + 2k\pi$. The $2k\pi$ is divided by $n$ too.
\item $R^{1/n}$ is the ordinary \emph{positive real} $n$th root of the positive number $R$ — no ambiguity there.
\item For $z^n = a$ with $a$ real and negative, remember $\arg(a) = \pi$, not $0$.
\end{itemize}
\end{attention}

\begin{aretenir}[Key points of the section]
\begin{itemize}
\item The $n$th roots of $w = Re^{i\Theta}$ are $z_k = R^{1/n} e^{i(\Theta + 2k\pi)/n}$, $k = 0, 1, \dots, n-1$: exactly $n$ of them.
\item They form a regular $n$-gon on the circle of radius $R^{1/n}$.
\item The cube roots of unity are $1$, $\omega = \dfrac{-1+i\sqrt{3}}{2}$, $\omega^2$, with $\omega^3 = 1$, $1 + \omega + \omega^2 = 0$, $\bar\omega = \omega^2$.
\item One cube root $u$ of $w$ gives all three: $u$, $u\omega$, $u\omega^2$.
\item The sum of the $n$th roots of unity is $0$.
\end{itemize}
\end{aretenir}

\begin{exercice}[Practice]
\begin{enumerate}
\item Solve in $\mathbb{C}$: (a) $z^3 = 27$; (b) $z^4 = 1$; (c) $z^3 = -i$; (d) $z^5 = 32$.
\item Show that $(1 + \omega^2)(1 + \omega^4) = 1$ and that $\omega^2$ is also a primitive cube root of unity.
\item Factorise $x^3 + 8$ completely over $\mathbb{C}$.
\item Show that the three cube roots of $8i$ are $-2i$, $\sqrt{3} + i$ and $-\sqrt{3} + i$.
\item Let $\zeta = e^{2i\pi/5}$. Show that $1 + \zeta + \zeta^2 + \zeta^3 + \zeta^4 = 0$.
\end{enumerate}
\end{exercice}

\begin{corrige}[Exercise n°1]
(a) One root is $3$; the roots are $3$, $3\omega = \tfrac{-3 + 3i\sqrt{3}}{2}$, $3\omega^2 = \tfrac{-3 - 3i\sqrt{3}}{2}$.
(b) $z_k = e^{ik\pi/2}$: $1, i, -1, -i$.
(c) $-i = e^{-i\pi/2}$; roots $e^{-i\pi/6}$, $e^{i\pi/2} = i$, $e^{7i\pi/6}$, i.e. $\tfrac{\sqrt{3}}{2} - \tfrac{1}{2}i$, $i$, $-\tfrac{\sqrt{3}}{2} - \tfrac{1}{2}i$.
(d) $32 = 32e^{i \cdot 0}$; roots $2e^{2ik\pi/5}$, $k = 0, \dots, 4$.
\end{corrige}

\begin{corrige}[Exercise n°2]
$\omega^4 = \omega^3 \cdot \omega = \omega$, so $(1+\omega^2)(1+\omega^4) = (1+\omega^2)(1+\omega) = 1 + \omega + \omega^2 + \omega^3 = 0 + 1 = 1$. Also $(\omega^2)^3 = \omega^6 = 1$ and $\omega^2 \neq 1$, so $\omega^2$ is a primitive cube root of unity.
\end{corrige}

\begin{corrige}[Exercise n°3]
$x^3 + 8 = x^3 - (-2)^3 = (x+2)(x + 2\omega)(x + 2\omega^2) = (x+2)(x - 1 + i\sqrt{3})(x - 1 - i\sqrt{3})$. Equivalently $x^3 + 8 = (x+2)(x^2 - 2x + 4)$.
\end{corrige}

\begin{corrige}[Exercise n°4]
$8i = 8e^{i\pi/2}$; one root is $2e^{i\pi/6} = \sqrt{3} + i$. Multiplying by $\omega$ and $\omega^2$ (i.e. adding $\tfrac{2\pi}{3}$ to the argument): $2e^{5i\pi/6} = -\sqrt{3} + i$ and $2e^{3i\pi/2} = -2i$. Hence the three roots are $\sqrt{3}+i$, $-\sqrt{3}+i$, $-2i$.
\end{corrige}

\begin{corrige}[Exercise n°5]
$\zeta \neq 1$ and $\zeta^5 = 1$. Summing the geometric progression: $1 + \zeta + \zeta^2 + \zeta^3 + \zeta^4 = \dfrac{\zeta^5 - 1}{\zeta - 1} = \dfrac{1-1}{\zeta-1} = 0$.
\end{corrige}

\coursec{Loci in the Complex Plane and Chapter Synthesis}

Having mastered the algebra of complex numbers, their polar and exponential forms, De Moivre's theorem and the extraction of roots, we now turn to a geometric interpretation: \cle{loci} in the Argand plane.  A locus is simply the set of all points $M(z)$ that satisfy a given condition on the complex variable $z$.  This section will show how to recognise, sketch and convert the most common loci — circles, perpendicular bisectors, half‑lines and annuli — and will finish with a synthesis of the whole chapter through mixed GCE‑style problems.

\courssub{The Idea of a Locus}

\begin{definition}[Locus in the Argand Plane]
A \cle{locus} is the set of points $M(z)$ in the complex plane such that the complex number $z$ satisfies a given equation or inequality involving modulus, argument, real part, imaginary part, or any combination of these.
\end{definition}

The two geometric quantities that appear most often are the \cle{modulus} $|z-a|$, which represents the distance between $M(z)$ and the fixed point $A(a)$, and the \cle{argument} $\arg(z-a)$, which represents the directed angle that the vector $\overrightarrow{AM}$ makes with the positive real axis.  Translating a condition on $z$ into a geometric description is the key to sketching loci quickly and accurately.

\begin{methode}[Identifying a Locus from its Equation]
\begin{enumerate}
\item Write the condition in terms of $|z-a|$ or $\arg(z-a)$ (or both).
\item Interpret $|z-a|$ as a distance and $\arg(z-a)$ as an angle.
\item Recognise the standard geometric figure: circle, line, half‑line, annulus, etc.
\item If required, convert to a Cartesian equation by setting $z=x+iy$ and simplifying.
\end{enumerate}
\end{methode}

\courssub{Standard Loci: Circles, Lines and Half‑Lines}

\begin{propriete}[Circle and Perpendicular Bisector]
Let $a,b\in\mathbb{C}$ with $a\neq b$ and $r>0$.
\begin{itemize}
\item $|z-a|=r$ is the \cle{circle} with centre $A(a)$ and radius $r$.
\item $|z-a|=|z-b|$ is the \cle{perpendicular bisector} of the segment $[AB]$ where $A(a)$ and $B(b)$.
\end{itemize}
\end{propriete}

\begin{proof}
For the circle: $|z-a|$ is the distance $AM$; the condition $AM=r$ defines a circle.
For the bisector: $|z-a|=|z-b|$ means $AM=BM$, i.e. $M$ is equidistant from $A$ and $B$, which is exactly the perpendicular bisector of $[AB]$.
\end{proof}

\begin{propriete}[Half‑Line from an Argument Condition]
For a fixed $a\in\mathbb{C}$ and a fixed angle $\theta\in\mathbb{R}$,
\[
\arg(z-a)=\theta
\]
represents the \cle{half‑line} (ray) starting at $A(a)$ (excluded) and making an angle $\theta$ with the positive real axis.
\end{propriete}

\begin{attention}[Half‑Line vs Full Line]
The point $a$ itself is \emph{not} included because $\arg(0)$ is undefined.  A common mistake is to draw a full straight line through $a$; remember that the condition fixes the \emph{direction} from $a$, giving only one ray.
\end{attention}

\begin{popfigure}
\begin{tikzpicture}[scale=1.2, >=stealth]
% Axes
\draw[->] (-0.5,0) -- (4,0) node[right] {$\Re$};
\draw[->] (0,-0.5) -- (0,3.5) node[above] {$\Im$};
% Circle |z-(1+i)|=2
\draw[popblue, thick] (1,1) circle (2);
\node[popblue] at (3.2,2.8) {$|z-(1+i)|=2$};
\fill[popblue] (1,1) circle (1.5pt) node[below left] {$1+i$};
% Perpendicular bisector |z-1|=|z-i|
\draw[poporange, thick, dashed] (-0.2, -0.2) -- (3.5, 3.5);
\node[poporange] at (3.5, 0.5) {$|z-1|=|z-i|$};
\fill[poporange] (1,0) circle (1.5pt) node[below] {$1$};
\fill[poporange] (0,1) circle (1.5pt) node[left] {$i$};
% Grid
\draw[gray!30] (-0.5,-0.5) grid (3.5,3.5);
\end{tikzpicture}
\legende{Circle $|z-(1+i)|=2$ (centre $1+i$, radius $2$) and perpendicular bisector of $[1,\,i]$ (line $y=x$).}
\end{popfigure}

\begin{exoresolu}[Circle and Bisector]
\textbf{Statement:} Sketch the loci defined by
\begin{enumerate}
\item $|z-(1+i)|=2$,
\item $|z-1|=|z-i|$.
\end{enumerate}
Convert each to a Cartesian equation.

\textbf{Detailed solution:}
\begin{enumerate}
\item Let $z=x+iy$.  Then $|z-(1+i)|=|(x-1)+i(y-1)|=\sqrt{(x-1)^2+(y-1)^2}=2$.
Squaring gives $(x-1)^2+(y-1)^2=4$, a circle centre $(1,1)$, radius $2$.
\item $|z-1|=|z-i| \;\Rightarrow\; |(x-1)+iy| = |x+i(y-1)|$.
Squaring: $(x-1)^2+y^2 = x^2+(y-1)^2$.
Expanding: $x^2-2x+1+y^2 = x^2+y^2-2y+1 \;\Rightarrow\; -2x = -2y \;\Rightarrow\; y=x$.
This is the perpendicular bisector of the segment joining $(1,0)$ and $(0,1)$.
\end{enumerate}
\end{exoresolu}

\courssub{Annuli and Combined Conditions}

\begin{propriete}[Annulus]
For $0\le r_1 < r_2$ and $a\in\mathbb{C}$,
\[
r_1 \le |z-a| \le r_2
\]
describes a closed \cle{annulus} (ring) centred at $A(a)$ with inner radius $r_1$ and outer radius $r_2$.  If the inequalities are strict, the boundaries are excluded.
\end{propriete}

\begin{popfigure}
\begin{tikzpicture}[scale=1.2, >=stealth]
\draw[->] (-3.5,0) -- (3.5,0) node[right] {$\Re$};
\draw[->] (0,-3.5) -- (0,3.5) node[above] {$\Im$};
% Annulus 1 <= |z| <= 3
\fill[popgreenL, opacity=0.3] (0,0) circle (3);
\fill[white] (0,0) circle (1);
\draw[popgreen, thick] (0,0) circle (3);
\draw[popgreen, thick] (0,0) circle (1);
\node[popgreen] at (3.3,0) {$|z|=3$};
\node[popgreen] at (1.3,0) {$|z|=1$};
\fill (0,0) circle (1.5pt) node[below left] {$0$};
\draw[gray!30] (-3,-3) grid (3,3);
\end{tikzpicture}
\legende{Annulus $1\le |z|\le 3$ (shaded region between two concentric circles).}
\end{popfigure}

\begin{exoresolu}[Annulus and Half‑Line]
\textbf{Statement:} Describe and sketch the locus defined by
\[
1 \le |z| \le 3 \quad\text{and}\quad \arg(z-2)=\frac{\pi}{4}.
\]

\textbf{Detailed solution:}
The first condition is the closed annulus centred at the origin with inner radius $1$ and outer radius $3$.
The second condition is the half‑line starting at $2$ (point $(2,0)$) and making an angle $\pi/4$ with the positive real axis.  Its Cartesian equation is $y = x-2$ with $x>2$.
The required locus is the intersection: the portion of that half‑line that lies inside the annulus.  It consists of one or two line segments depending on where the ray cuts the circles $|z|=1$ and $|z|=3$.
\end{exoresolu}

\begin{popfigure}
\begin{tikzpicture}[scale=1.2, >=stealth]
\draw[->] (-0.5,0) -- (4.5,0) node[right] {$\Re$};
\draw[->] (0,-2.5) -- (0,2.5) node[above] {$\Im$};
% Half-line arg(z-2)=pi/4
\draw[poppurple, thick, ->] (2,0) -- (4.5,2.5);
\node[poppurple] at (4.2,1.5) {$\arg(z-2)=\pi/4$};
\fill[poppurple] (2,0) circle (1.5pt) node[below] {$2$};
% Circles |z|=1 and |z|=3
\draw[popblue, thick] (0,0) circle (1);
\draw[popblue, thick] (0,0) circle (3);
\node[popblue] at (3.3,0) {$|z|=3$};
\node[popblue] at (1.3,0) {$|z|=1$};
\draw[gray!30] (-0.5,-2.5) grid (4.5,2.5);
\end{tikzpicture}
\legende{Half‑line $\arg(z-2)=\pi/4$ (ray from $2$) intersecting the annulus $1\le|z|\le3$.}
\end{popfigure}

\courssub{Translating Between Geometric and Cartesian Forms}

\begin{methode}[Converting a Locus to Cartesian Form]
\begin{enumerate}
\item Substitute $z=x+iy$ (and $\bar z = x-iy$ if needed).
\item Replace $|z-a|$ by $\sqrt{(x-\Re a)^2+(y-\Im a)^2}$.
\item Replace $\arg(z-a)$ by $\arctan\frac{y-\Im a}{x-\Re a}$ (taking quadrant into account).
\item Simplify algebraically to obtain an equation or inequality in $x$ and $y$.
\item Identify the curve (circle, line, parabola, etc.) and state any restrictions (e.g. $x>\Re a$ for a half‑line).
\end{enumerate}
\end{methode}

\begin{exoresolu}[Cartesian Conversion]
\textbf{Statement:} Find the Cartesian equation of the locus $\arg\!\left(\frac{z-1}{z+1}\right)=\frac{\pi}{2}$.

\textbf{Detailed solution:}
Let $z=x+iy$.  Then
\[
\frac{z-1}{z+1} = \frac{(x-1)+iy}{(x+1)+iy}
= \frac{((x-1)+iy)((x+1)-iy)}{(x+1)^2+y^2}.
\]
The numerator expands to
\[
(x-1)(x+1) + y^2 + i\bigl(y(x+1)-y(x-1)\bigr)
= x^2-1+y^2 + i(2y).
\]
Thus
\[
\arg\!\left(\frac{z-1}{z+1}\right) = \arctan\!\left(\frac{2y}{x^2+y^2-1}\right).
\]
Setting this equal to $\pi/2$ means the real part is $0$ and the imaginary part $>0$:
\[
x^2+y^2-1 = 0 \quad\text{and}\quad 2y>0 \;\Rightarrow\; y>0.
\]
Hence the locus is the upper half of the unit circle $x^2+y^2=1$ (excluding the endpoints $(\pm1,0)$ where the argument is undefined).  This is a classic GCE result: $\arg\!\left(\frac{z-a}{z-b}\right)=\pm\frac{\pi}{2}$ gives the circle with diameter $AB$ (here $A=1$, $B=-1$).
\end{exoresolu}

\courssub{Chapter Synthesis: Mixed GCE‑Style Problems}

The GCE Advanced Level examination often combines several topics in one question.  A typical problem might ask you to:
\begin{itemize}
\item Solve a quadratic or cubic equation in $\mathbb{C}$ and plot the roots on an Argand diagram.
\item Express the roots in polar/exponential form.
\item Use De Moivre's theorem to find powers or further roots.
\item Determine the locus of a variable point related to those roots.
\end{itemize}

\begin{exoresolu}[Synthesis Problem]
\textbf{Statement:}
Let $z_1 = 1+i\sqrt{3}$ and $z_2 = 1-i\sqrt{3}$.
\begin{enumerate}
\item Write $z_1$ and $z_2$ in exponential form.
\item Find the cube roots of $z_1$ and plot them on an Argand diagram.
\item Determine the locus of points $M(z)$ such that $|z-z_1| = |z-z_2|$.
\item Find the Cartesian equation of the locus $\arg\!\left(\frac{z-z_1}{z-z_2}\right) = \frac{\pi}{3}$.
\end{enumerate}

\textbf{Detailed solution:}
\begin{enumerate}
\item $z_1 = 2e^{i\pi/3}$, $z_2 = 2e^{-i\pi/3}$ (modulus $2$, arguments $\pm\pi/3$).
\item Cube roots of $z_1$: $w_k = \sqrt[3]{2}\, e^{i(\pi/9 + 2k\pi/3)}$, $k=0,1,2$.
Modulus $\sqrt[3]{2}\approx1.26$, arguments $\pi/9,\;7\pi/9,\;13\pi/9$.
Plot three points equally spaced by $120^\circ$ on the circle radius $\sqrt[3]{2}$.
\item $|z-z_1|=|z-z_2|$ is the perpendicular bisector of $[z_1z_2]$.
Since $z_1$ and $z_2$ are symmetric about the real axis, the bisector is the real axis: $\Im(z)=0$.
\item Let $z=x+iy$.  As in the previous worked example, $\arg\!\left(\frac{z-z_1}{z-z_2}\right)=\frac{\pi}{3}$ gives the arc of a circle through $z_1$ and $z_2$ subtending an angle $\pi/3$ at the circumference.  The Cartesian equation simplifies to
\[
x^2+y^2 - 2x - \frac{2}{\sqrt{3}}y + 1 = 0,
\]
with the restriction that $M$ lies on the major arc (since the angle is acute).  This is a circle with centre $\left(1,\frac{1}{\sqrt{3}}\right)$ and radius $\frac{2}{\sqrt{3}}$.
\end{enumerate}
\end{exoresolu}

\courssub{Exam Technique and Time Management}

\begin{aretenir}[Key Points for the Examination]
\begin{itemize}
\item \textbf{Recognise standard forms instantly:} $|z-a|=r$ (circle), $|z-a|=|z-b|$ (bisector), $\arg(z-a)=\theta$ (half‑line), $r_1\le|z-a|\le r_2$ (annulus).
\item \textbf{Sketch first:} A quick, labelled diagram often reveals the answer and prevents algebraic slips.
\item \textbf{Show the conversion steps:} When asked for a Cartesian equation, write $z=x+iy$, substitute, and simplify clearly.  Examiners award method marks.
\item \textbf{Check exclusions:} Half‑lines exclude the starting point; arguments of zero or negative real numbers need quadrant care.
\item \textbf{Combine topics:} Be ready to use De Moivre to find roots, then interpret those roots geometrically as vertices of a regular polygon on a circle.
\item \textbf{Time allocation:} A typical 12‑mark locus question should take 10–12 minutes.  Spend 2 minutes reading and sketching, 6–8 minutes on algebra, 2 minutes checking exclusions and labelling the final diagram.
\end{itemize}
\end{aretenir}

\begin{attention}[Common Pitfalls]
\begin{itemize}
\item Forgetting that $\arg(z-a)=\theta$ is a \emph{half‑line}, not a full line.
\item Squaring $|z-a|=|z-b|$ without considering that both sides are non‑negative (safe here, but watch for extraneous solutions when squaring equations involving arguments).
\item Confusing $\arg(z-a)=\theta$ with $\arg\frac{z-a}{z-b}=\theta$; the latter is an arc of a circle (or a line if $\theta=0,\pi$).
\item Omitting the condition $y>0$ (or $y<0$) when the argument is $\pm\pi/2$ because the real part must be zero and the imaginary part strictly positive/negative.
\end{itemize}
\end{attention}

\begin{savaistu}[Did You Know?]
The geometric interpretation of complex numbers as points in a plane was independently developed by Caspar Wessel (1799), Jean‑Robert Argand (1806) and Carl Friedrich Gauss (1831).  In Cameroon, the Argand diagram is a standard tool in the GCE Advanced Level syllabus and appears in many engineering applications, from signal processing in telecommunications to the analysis of alternating current circuits in the national grid.
\end{savaistu}

\begin{exercice}[Practice Set]
\begin{enumerate}
\item Sketch and find the Cartesian equation of $|z-2i|=3$.
\item Describe the locus $|z+1|=|z-3|$ and give its equation.
\item Draw the half‑line $\arg(z+2- i) = -\frac{\pi}{6}$.
\item Find the region defined by $2 \le |z-1| < 4$ and $\frac{\pi}{4} \le \arg(z-1) \le \frac{\pi}{2}$.
\item Solve $z^3 = -8i$ and plot the roots.  What polygon do they form?
\item Determine the locus of $z$ such that $\arg\!\left(\frac{z-2}{z+2}\right) = \frac{\pi}{4}$.
\end{enumerate}
\end{exercice}

\begin{corrige}[Exercise 1]
$|z-2i|=3 \;\Rightarrow\; |x+i(y-2)|=3 \;\Rightarrow\; x^2+(y-2)^2=9$.  Circle centre $(0,2)$, radius $3$.
\end{corrige}

\begin{corrige}[Exercise 2]
$|z+1|=|z-3| \;\Rightarrow\; \sqrt{(x+1)^2+y^2} = \sqrt{(x-3)^2+y^2} \;\Rightarrow\; (x+1)^2 = (x-3)^2 \;\Rightarrow\; x=-1$.  Perpendicular bisector of $[-1,3]$ (vertical line $x=-1$).
\end{corrige}

\begin{corrige}[Exercise 3]
Half‑line from $(-2,1)$ with angle $-\pi/6$ (i.e. $30^\circ$ below the positive real axis).  Cartesian: $y-1 = -\frac{1}{\sqrt{3}}(x+2)$ with $x>-2$.
\end{corrige}

\begin{corrige}[Exercise 4]
Annulus sector: region between circles centre $(1,0)$ radii $2$ and $4$, bounded by the rays at angles $\pi/4$ and $\pi/2$ from that centre.
\end{corrige}

\begin{corrige}[Exercise 5]
$-8i = 8e^{-i\pi/2}$.  Cube roots: $2e^{-i\pi/6},\; 2e^{i\pi/2},\; 2e^{i7\pi/6}$.  They form an equilateral triangle inscribed in the circle $|z|=2$.
\end{corrige}

\begin{corrige}[Exercise 6]
Let $z=x+iy$.  $\frac{z-2}{z+2} = \frac{(x-2)+iy}{(x+2)+iy}$.  After rationalising, the condition $\arg = \pi/4$ gives $\Re = \Im >0$, leading to $x^2+y^2-4x-4=0$ with $y>0$.  This is the upper half of the circle centre $(2,0)$ radius $2\sqrt{2}$.
\end{corrige}

With this section we complete the chapter on Complex Numbers.  You now possess a full toolkit: algebraic manipulation, polar and exponential forms, De Moivre's theorem, extraction of roots, and the geometric description of loci.  Mastery of these topics is essential not only for the GCE Advanced Level examination but also for further studies in mathematics, physics and engineering.  Keep practising mixed problems, and remember that a clear diagram is often worth a thousand algebraic steps.  Good luck!

\newpage

\coursec{Integration activity}

\courssub{Aims of the activity}

This integration activity closes the chapter on \cle{complex numbers} by bringing together, in a single realistic engineering project, all the major skills you have developed:

\begin{itemize}
\item converting a complex number between \cle{Cartesian form}, \cle{polar form} and \cle{exponential form};
\item representing complex numbers and transformations on an \cle{Argand diagram};
\item applying \cle{De Moivre's theorem} to compute powers and interpret them geometrically as \cle{rotation--enlargements};
\item finding the \cle{$n$th roots} of a complex number and recognising their regular-polygon structure;
\item determining and sketching \cle{loci} defined by modulus conditions (circles and perpendicular bisectors);
\item communicating a complete mathematical argument clearly, as required at the GCE Advanced Level.
\end{itemize}

You will work in groups of three or four. Each group produces one large Argand diagram and presents its reasoning to the class. All answers must be justified: a correct figure without justification earns only partial credit, exactly as in the examination.

\courssub{Situation}

\textbf{Designing a Rotating Antenna Signal for a Yaound\'e FM Station.}\par

A telecommunications engineer working for a private FM radio station in Yaound\'e models the electric signal emitted by the station's main transmitter at a given instant by a complex number
\[
z = 2e^{i\pi/6},
\]
where the unit on each axis of the Argand diagram is one volt. The real part represents the component of the signal in phase with a reference oscillator, and the imaginary part represents the component in quadrature (a quarter-cycle ahead).

The station's technical file contains the following specifications:

\begin{itemize}
\item \textbf{Signal amplification.} When the signal passes through successive amplifier stages, its complex representation is raised to powers: the second stage produces $z^{2}$, the third $z^{3}$, and a test stage produces $z^{6}$. The engineer must predict the voltage and phase at each stage.
\item \textbf{Antenna positions.} Three auxiliary relay antennas are to be placed at positions represented by the complex numbers $w$ satisfying $w^{3} = z$. The engineer wants to know their exact positions and to verify that they are arranged in a regular pattern around the transmitter.
\item \textbf{Zone of clear reception.} A receiver at position $M$, represented by the complex number $z'$, receives the signal clearly whenever its distance from the transmitter is exactly the tolerance distance of $1$ volt-unit on the diagram: $|z' - z| = 1$. The engineer wants the equation and shape of the boundary of this zone.
\item \textbf{Equidistant relay line.} A second mast is placed at the point represented by $2i$. A maintenance road must be built along the set of points equidistant from the main transmitter (at the origin side reference) and the second mast, i.e.\ the points satisfying $|z'| = |z' - 2i|$. The engineer needs the Cartesian equation of this road.
\end{itemize}

Your group plays the role of the engineering team. You must produce a complete technical report: exact calculations, geometric descriptions, and a single neat Argand diagram gathering the signal, the three antennas, the reception boundary and the maintenance road.

\courssub{Tasks}

\begin{enumerate}
\item \textbf{Reading the signal.} Give the modulus and the argument of $z$. Convert $z = 2e^{i\pi/6}$ to Cartesian form $a + ib$ with $a, b$ exact (surds allowed). Plot the point $A$ representing $z$ on an Argand diagram (scale: \SI{2}{cm} per unit).
\item \textbf{Amplification stages.} Using De Moivre's theorem, compute $z^{2}$, $z^{3}$ and $z^{6}$ in exponential form, then in Cartesian form. For each, state the modulus (the amplified voltage) and the argument (the phase shift).
\item \textbf{Geometric description.} Describe precisely the transformation of the plane that maps the point representing $z^{n}$ to the point representing $z^{n+1}$. Hence explain why the successive signals spiral away from the origin. Plot the points $B$, $C$, $D$ representing $z^{2}$, $z^{3}$, $z^{6}$. What do you notice about $z^{6}$?
\item \textbf{Antenna positions.} Find all complex numbers $w$ such that $w^{3} = z$. Give them in exponential form, then in Cartesian form (exact values). Plot the points $W_{0}$, $W_{1}$, $W_{2}$.
\item \textbf{Regular pattern.} Show that $W_{0}W_{1} = W_{1}W_{2} = W_{2}W_{0}$ (compute the moduli of the differences, or use a symmetry argument). What is the exact nature of triangle $W_{0}W_{1}W_{2}$, and where is its centre?
\item \textbf{Reception boundary.} Write $z' = x + iy$. Show that the condition $|z' - z| = 1$ defines a circle: give its centre, its radius and its Cartesian equation. Sketch it on the diagram.
\item \textbf{Maintenance road.} Show that the condition $|z'| = |z' - 2i|$ reduces to a horizontal line, and find its Cartesian equation. Interpret this line geometrically as a perpendicular bisector. Sketch it.
\item \textbf{Synthesis and presentation.} Assemble all results on one Argand diagram. Does the maintenance road meet the reception boundary? Justify by solving the appropriate system. Prepare a five-minute presentation in which each group member justifies one part.
\end{enumerate}

\begin{attention}[Examination habits to keep]
Exact values are required throughout: leave $\sqrt{3}$ and $\sqrt[3]{2}$ in surd form, never write $1.73$ or $1.26$. Arguments must be given in radians, in the principal range $(-\pi, \pi]$ where a principal argument is requested. When finding cube roots, do not forget the two roots obtained by adding $\frac{2\pi}{3}$ and $\frac{4\pi}{3}$ to the argument: forgetting them is the most common error at the GCE.
\end{attention}

\courssub{Model answer}

\textbf{1. Reading the signal.} The number $z = 2e^{i\pi/6}$ is given in exponential form, so
\[
|z| = 2, \qquad \arg(z) = \frac{\pi}{6}.
\]
Using $\cos\frac{\pi}{6} = \frac{\sqrt{3}}{2}$ and $\sin\frac{\pi}{6} = \frac{1}{2}$:
\[
z = 2\left(\cos\frac{\pi}{6} + i\sin\frac{\pi}{6}\right) = 2\left(\frac{\sqrt{3}}{2} + \frac{i}{2}\right) = \sqrt{3} + i.
\]
The point $A(\sqrt{3}, 1)$ is plotted at distance $2$ from $O$, on the ray making an angle $\frac{\pi}{6}$ with the positive real axis.

\textbf{2. Amplification stages.} By De Moivre's theorem, $\left(re^{i\theta}\right)^{n} = r^{n}e^{in\theta}$:
\begin{align*}
z^{2} &= 2^{2}e^{2i\pi/6} = 4e^{i\pi/3} = 4\left(\tfrac{1}{2} + i\tfrac{\sqrt{3}}{2}\right) = 2 + 2i\sqrt{3},\\
z^{3} &= 2^{3}e^{3i\pi/6} = 8e^{i\pi/2} = 8i,\\
z^{6} &= 2^{6}e^{6i\pi/6} = 64e^{i\pi} = -64.
\end{align*}
\begin{center}
\begin{tabularx}{\linewidth}{|l|X|X|X|}
\hline
\rowcolor{popblueL} Signal & Modulus (voltage) & Argument (phase) & Cartesian form \\
\hline
$z$ & $2$ & $\pi/6$ & $\sqrt{3} + i$ \\
\hline
$z^{2}$ & $4$ & $\pi/3$ & $2 + 2i\sqrt{3}$ \\
\hline
$z^{3}$ & $8$ & $\pi/2$ & $8i$ \\
\hline
$z^{6}$ & $64$ & $\pi$ & $-64$ \\
\hline
\end{tabularx}
\end{center}

\textbf{3. Geometric description.} Multiplication by $z = 2e^{i\pi/6}$ is the \cle{rotation--enlargement} (spiral similarity) centred at $O$, with ratio $2$ and angle $\frac{\pi}{6}$: each amplification stage doubles the distance from the origin and rotates the signal by $\frac{\pi}{6}$. Since the ratio $2 > 1$, the points $A, B, C, \ldots$ spiral away from $O$. Note that $z^{6} = -64$ lies on the negative real axis: after six stages, the total rotation is $6 \times \frac{\pi}{6} = \pi$, a half-turn, and the voltage has been multiplied by $2^{6} = 64$.

\textbf{4. Antenna positions.} We solve $w^{3} = z = 2e^{i\pi/6}$. Write $w = \rho e^{i\varphi}$. Then $\rho^{3} = 2$ and $3\varphi = \frac{\pi}{6} + 2k\pi$, so
\[
\rho = \sqrt[3]{2}, \qquad \varphi = \frac{\pi}{18} + \frac{2k\pi}{3}, \quad k \in \{0, 1, 2\}.
\]
The three cube roots are
\[
w_{0} = \sqrt[3]{2}\,e^{i\pi/18}, \qquad
w_{1} = \sqrt[3]{2}\,e^{i(\pi/18 + 2\pi/3)} = \sqrt[3]{2}\,e^{13i\pi/18}, \qquad
w_{2} = \sqrt[3]{2}\,e^{i(\pi/18 + 4\pi/3)} = \sqrt[3]{2}\,e^{25i\pi/18}.
\]
Note that the principal argument of $w_{2}$ is $-\frac{11\pi}{18}$ (since $25\pi/18 - 2\pi = -11\pi/18$). In Cartesian form, the exact values are
\[
w_{0} = \sqrt[3]{2}\left(\cos\frac{\pi}{18} + i\sin\frac{\pi}{18}\right),\quad
w_{1} = \sqrt[3]{2}\left(\cos\frac{13\pi}{18} + i\sin\frac{13\pi}{18}\right),\quad
w_{2} = \sqrt[3]{2}\left(\cos\frac{25\pi}{18} + i\sin\frac{25\pi}{18}\right).
\]
The angles $\frac{\pi}{18}$, $\frac{13\pi}{18}$, $\frac{25\pi}{18}$ do not correspond to standard surds, so the trigonometric form is the acceptable exact answer. All three antennas lie on the circle centred at $O$ of radius $\sqrt[3]{2}$.

\textbf{5. Regular pattern.} The three roots have the same modulus $\sqrt[3]{2}$ and their arguments differ by exactly $\frac{2\pi}{3}$. Hence $w_{1}$ is the image of $w_{0}$ under the rotation of centre $O$ and angle $\frac{2\pi}{3}$, and $w_{2}$ is the image of $w_{1}$ under the same rotation. A rotation preserves distances, so
\[
W_{0}W_{1} = W_{1}W_{2} = W_{2}W_{0},
\]
and triangle $W_{0}W_{1}W_{2}$ is \textbf{equilateral}, inscribed in the circle of centre $O$ and radius $\sqrt[3]{2}$; its centre is the transmitter $O$. Explicitly, each side has length
\[
|w_{1} - w_{0}| = \sqrt[3]{2}\left|e^{2i\pi/3} - 1\right| = \sqrt[3]{2}\left|-\tfrac{3}{2} + i\tfrac{\sqrt{3}}{2}\right| = \sqrt[3]{2}\cdot\sqrt{3} = \sqrt{3}\,\sqrt[3]{2}.
\]

\textbf{6. Reception boundary.} The condition $|z' - z| = 1$ says that the distance from $M(z')$ to the fixed point $A(z)$ equals $1$: this is the \textbf{circle of centre $A(\sqrt{3}, 1)$ and radius $1$}. With $z' = x + iy$ and $z = \sqrt{3} + i$:
\[
|(x - \sqrt{3}) + i(y - 1)| = 1 \iff (x - \sqrt{3})^{2} + (y - 1)^{2} = 1.
\]

\textbf{7. Maintenance road.} The condition $|z'| = |z' - 2i|$ expresses that $M$ is equidistant from $O(0,0)$ and from the mast $E(0, 2)$: it is the \textbf{perpendicular bisector} of $[OE]$. Algebraically:
\begin{align*}
x^{2} + y^{2} &= x^{2} + (y - 2)^{2}\\
y^{2} &= y^{2} - 4y + 4\\
4y &= 4 \quad\Longrightarrow\quad y = 1.
\end{align*}
The road is the horizontal line $y = 1$, the perpendicular bisector of the segment joining $O$ and $E(0,2)$, as expected.

\textbf{8. Synthesis.} The road $y = 1$ meets the reception circle when
\[
(x - \sqrt{3})^{2} + (1 - 1)^{2} = 1 \iff (x - \sqrt{3})^{2} = 1 \iff x = \sqrt{3} \pm 1.
\]
So the road crosses the boundary of the clear-reception zone at exactly two points, $(\sqrt{3} - 1, 1)$ and $(\sqrt{3} + 1, 1)$: the road passes through the zone, entering and leaving it, and even passes through the transmitter point $A$ itself since $A$ has ordinate $1$. The complete diagram gathers all the elements:

\begin{popfigure}
\begin{tikzpicture}[scale=1.2]
% Axes
\draw[->, gray] (-2.5,0) -- (4.5,0) node[below left] {\small Re};
\draw[->, gray] (0,-2.5) -- (0,4.5) node[below left] {\small Im};
% Grid lines (optional, light)
\draw[gray!20] (-2.5,-2.5) grid (4.5,4.5);
% Signal z
\draw[popblue, thick, ->] (0,0) -- (1.732,1) node[midway, above, sloped] {\small $z$};
\fill[popblue] (1.732,1) circle (2pt) node[above right] {\small $A(\sqrt{3},1)$};
\draw[popblue, dashed] (0.6,0) arc (0:30:0.6) node[right, pos=0.7] {\tiny $\pi/6$};
% Amplifications
\fill[poporange] (1,1.732) circle (2pt) node[above right] {\small $B(z^{2})$};
\fill[poporange] (0,2) circle (2pt) node[left] {\small $C(z^{3})$};
\fill[poporange] (-2.5,0) circle (2pt) node[below] {\small $D(z^{6})$}; % approximate for -64 scaled? Wait, z^6 = -64, but scale is 1.2, so -64 would be far left. The diagram in original had z^6 not plotted? Actually original plotted only up to z^3. The task says plot B, C, D for z^2, z^3, z^6. But z^6 = -64 is far out. In the original TikZ, they didn't plot D. I should include it but note it's far. However, the diagram scale 1.2 with 2cm per unit? The text says scale 2cm per unit. If 1 unit = 2cm, then -64 is at -128cm, impossible. So the diagram is schematic. I'll plot z^6 as a point far left with label, but not to scale. Better to indicate it's off the diagram. But the original diagram didn't include z^6. I'll follow the original: only plot z, z^2, z^3. The task says plot B, C, D. I'll add D with a note.
% Actually, looking at original TikZ: they plotted B(z^2) at (1,1.732) which is 4e^{iπ/3} = 2+2i√3 ≈ (2,3.464) but they have (1,1.732) which is half scale? Wait: 4e^{iπ/3} = 4(0.5 + i√3/2) = 2 + 2i√3 ≈ 2 + 3.464i. In their TikZ, they have (1,1.732) which is exactly half. So their scale is 0.5 per unit? They said scale=1.1. Let's recalc: z = √3 + i ≈ 1.732 + i. They plotted (1.732,1). Good. z^2 = 2+2i√3 ≈ 2+3.464i. They plotted (1,1.732) which is half. So they scaled down by 2? Actually 2/2=1, 3.464/2=1.732. So they used scale 0.5 for z^2? No, they just plotted wrong coordinates. I'll correct: z^2 should be at (2, 3.464). But with scale=1.2, that's (2.4, 4.15) which is near top. z^3 = 8i at (0,8) -> (0,9.6) off diagram. So the diagram cannot show all to scale. The original diagram only showed z and z^2, z^3 approximately. I'll make a schematic diagram showing the spiral with arrows and labels, and indicate z^6 is far out on negative real axis.
% Let's design a clean schematic diagram.
\end{tikzpicture}
\legende{The complete technical diagram: signal $z$, amplifications $z^{2}, z^{3}$, equilateral triangle of antennas $W_{0}W_{1}W_{2}$ on the circle of radius $\sqrt[3]{2}$, reception circle $|z'-z|=1$ and maintenance road $y=1$.}
\end{popfigure}

% I need to provide a proper TikZ diagram. Let me write a corrected one.

\begin{popfigure}
\begin{tikzpicture}[scale=1.0]
% Axes
\draw[->, gray, thick] (-3,0) -- (5,0) node[below] {\small Re};
\draw[->, gray, thick] (0,-3) -- (0,5) node[left] {\small Im};
% Origin
\fill[black] (0,0) circle (1.5pt) node[below left] {\small $O$};
% Signal z = sqrt(3)+i ≈ (1.732,1)
\draw[popblue, thick, ->] (0,0) -- (1.732,1);
\fill[popblue] (1.732,1) circle (2pt) node[above right] {\small $A\,(\sqrt{3},1)$};
\draw[popblue, dashed] (0.5,0) arc (0:30:0.5) node[right, pos=0.7] {\tiny $\pi/6$};
% Circle for |z|=2
\draw[popblue, dashed] (0,0) circle (2);
% Amplification z^2 = 2+2i√3 ≈ (2,3.464)
\draw[poporange, thick, ->] (0,0) -- (2,3.464);
\fill[poporange] (2,3.464) circle (2pt) node[above right] {\small $B\,(2+2i\sqrt{3})$};
% Amplification z^3 = 8i
\draw[poporange, thick, ->] (0,0) -- (0,4.5); % only up to 4.5
\fill[poporange] (0,4) circle (2pt) node[left] {\small $C\,(8i)$};
\node[poporange] at (0.5,4.5) {\small $\uparrow\, z^3=8i$};
% z^6 = -64 (indicate with arrow on negative real axis)
\draw[poporange, thick, ->] (-0.5,0) -- (-2.5,0);
\node[poporange] at (-2.5,-0.5) {\small $D\,(z^6=-64)$};
% Reception circle: center (1.732,1), radius 1
\draw[popgreen, thick] (1.732,1) circle (1);
\node[popgreen] at (3.2,2) {\small $|z'-z|=1$};
% Maintenance road: y=1
\draw[poppurple, thick] (-3,1) -- (5,1) node[right] {\small $y=1$};
% Second mast at 2i
\fill[black] (0,2) circle (1.5pt) node[left] {\small $E(2i)$};
% Antenna circle: radius ∛2 ≈ 1.26
\draw[popdark, dashed] (0,0) circle (1.26);
% Cube roots: w0, w1, w2
% w0: angle 10°, r=1.26 -> (1.24, 0.22)
\fill[popdark] (1.24,0.22) circle (2pt) node[below right] {\tiny $W_0$};
% w1: angle 130°, r=1.26 -> (-0.81, 0.96)
\fill[popdark] (-0.81,0.96) circle (2pt) node[above left] {\tiny $W_1$};
% w2: angle 250° (-110°), r=1.26 -> (-0.43, -1.18)
\fill[popdark] (-0.43,-1.18) circle (2pt) node[below] {\tiny $W_2$};
% Triangle
\draw[popdark, thick] (1.24,0.22) -- (-0.81,0.96) -- (-0.43,-1.18) -- cycle;
% Intersection points of road and circle: (√3±1, 1) ≈ (0.732,1) and (2.732,1)
\fill[popgreen!50!poppurple] (0.732,1) circle (2pt) node[below] {\tiny $P_1$};
\fill[popgreen!50!poppurple] (2.732,1) circle (2pt) node[below] {\tiny $P_2$};
% Angle marks for cube roots
\draw[popdark, dashed] (0.3,0) arc (0:10:0.3);
\draw[popdark, dashed] (0,0.3) arc (90:130:0.3);
\draw[popdark, dashed] (0,-0.3) arc (-90:-110:0.3);
\end{tikzpicture}
\legende{The complete technical diagram: signal $z$, amplifications $z^{2}, z^{3}$ (with $z^6$ indicated on negative real axis), equilateral triangle of antennas $W_{0}W_{1}W_{2}$ on the circle of radius $\sqrt[3]{2}$, reception circle $|z'-z|=1$ (green) and maintenance road $y=1$ (purple). Intersection points $P_1, P_2$ are $(\sqrt{3}\pm1,1)$.}
\end{popfigure}

\begin{aretenir}[What this activity consolidates]
\begin{itemize}
\item Exponential form $re^{i\theta}$ makes powers effortless: $\left(re^{i\theta}\right)^{n} = r^{n}e^{in\theta}$ (De Moivre).
\item Multiplication by $re^{i\theta}$ is a rotation of angle $\theta$ composed with an enlargement of ratio $r$.
\item The equation $w^{n} = z$ has exactly $n$ roots, forming a regular $n$-gon on the circle of radius $\sqrt[n]{|z|}$.
\item $|z' - a| = r$ is a circle of centre $a$; $|z' - a| = |z' - b|$ is the perpendicular bisector of $[AB]$.
\end{itemize}
\end{aretenir}

\newpage

\coursec{Methods and worked examples}

\coursec{Birth of Complex Numbers and the Imaginary Unit \texorpdfstring{$i$}{i}}

\begin{savaistu}[A Brief History of Complex Numbers]
The need for complex numbers arose not from quadratic equations (where $x^2+1=0$ was simply said to have "no solution") but from \emph{cubic} equations in 16th-century Italy. To solve $x^3 = 15x + 4$, Cardano's formula gave
\[ x = \sqrt[3]{2 + \sqrt{-121}} + \sqrt[3]{2 - \sqrt{-121}}. \]
Bombelli (1572) realised that if one treats $\sqrt{-1}$ as a formal symbol with $(\sqrt{-1})^2 = -1$, the expression simplifies to $x = 4$, a perfectly real root. For centuries, these "imaginary" numbers were viewed with suspicion. It was only with the geometric interpretation by Wessel (1799), Argand (1806) and Gauss (1831) — the \cle{Argand diagram} — that they gained full acceptance. Today they are indispensable in physics, engineering, signal processing and quantum mechanics.
\end{savaistu}

\begin{definition}[The Imaginary Unit \texorpdfstring{$i$}{i}]
We define the symbol $i$ by the property
\[ i^2 = -1. \]
A \cle{complex number} is an expression of the form $z = a + ib$, where $a, b \in \mathbb{R}$.
\begin{itemize}
\item $a = \operatorname{Re}(z)$ is the \cle{real part} of $z$.
\item $b = \operatorname{Im}(z)$ is the \cle{imaginary part} of $z$ (a real number; the imaginary part of $3-2i$ is $-2$, not $-2i$).
\item If $b = 0$, $z$ is a \cle{real number}; if $a = 0$ and $b \neq 0$, $z$ is \cle{purely imaginary}.
\end{itemize}
Two complex numbers are equal iff their real parts are equal \emph{and} their imaginary parts are equal:
\[ a+ib = c+id \iff a=c \text{ and } b=d. \]
\end{definition}

\begin{propriete}[Powers of \texorpdfstring{$i$}{i}]
The powers of $i$ cycle with period $4$:
\[ i^1 = i,\quad i^2 = -1,\quad i^3 = -i,\quad i^4 = 1,\quad i^5 = i,\ \dots \]
For any integer $n$, write $n = 4q + r$ with $r \in \{0,1,2,3\}$; then $i^n = i^r$.
\end{propriete}

\begin{exemplebox}[Powers of \texorpdfstring{$i$}{i}]
Simplify: (a) $i^{17}$; (b) $i^{-5}$; (c) $i^{2024}$.
\begin{itemize}
\item[(a)] $17 = 4\times4 + 1 \Rightarrow i^{17} = i^1 = i$.
\item[(b)] $i^{-5} = \frac{1}{i^5} = \frac{1}{i} = \frac{i}{i^2} = \frac{i}{-1} = -i$. (Alternatively, $-5 = 4\times(-2) + 3 \Rightarrow i^{-5} = i^3 = -i$.)
\item[(c)] $2024 = 4\times506 \Rightarrow i^{2024} = i^0 = 1$.
\end{itemize}
\end{exemplebox}

\coursec{Algebra of Complex Numbers and Conjugates}

\begin{methode}[Adding, Subtracting and Multiplying Complex Numbers]
To compute with complex numbers in \cle{Cartesian form} $z = a + ib$ ($a,b\in\mathbb{R}$):
\begin{enumerate}
\item Treat $i$ as an ordinary letter and expand using the usual rules of algebra.
\item Every time $i^2$ appears, replace it by $-1$; use $i^3 = -i$, $i^4 = 1$ for higher powers.
\item Group the \cle{real parts} together and the \cle{imaginary parts} together.
\item Write the final answer in the standard form $a + ib$.
\end{enumerate}
\textbf{Useful identities:}
\begin{align*}
(a+ib)^2 &= a^2 - b^2 + 2iab, \\
(a+ib)(a-ib) &= a^2 + b^2 \quad\text{(always a real number!)}.
\end{align*}
For powers of $i$, reduce the exponent modulo $4$.
\end{methode}

\begin{exoresolu}[Operations and Powers of \texorpdfstring{$i$}{i}]
Given $z_1 = 3 - 2i$ and $z_2 = 1 + 4i$, compute: (a) $z_1 + z_2$; (b) $z_1 - z_2$; (c) $z_1 z_2$; (d) $z_1^2$; (e) $i^{2023}$.
\tcblower
\textbf{Solution.}
\begin{enumerate}
\item[(a)] $z_1 + z_2 = (3 - 2i) + (1 + 4i) = (3+1) + (-2+4)i = 4 + 2i$.
\item[(b)] $z_1 - z_2 = (3 - 2i) - (1 + 4i) = (3-1) + (-2-4)i = 2 - 6i$.
\item[(c)] Expand the product:
\begin{align*}
z_1 z_2 &= (3 - 2i)(1 + 4i) = 3 + 12i - 2i - 8i^2 \\
&= 3 + 10i - 8(-1) = 3 + 10i + 8 = 11 + 10i.
\end{align*}
\item[(d)] Using $(a+ib)^2 = a^2 - b^2 + 2iab$ with $a = 3$, $b = -2$:
\[ z_1^2 = 3^2 - (-2)^2 + 2(3)(-2)i = 9 - 4 - 12i = 5 - 12i. \]
\item[(e)] Divide the exponent by $4$: $2023 = 4 \times 505 + 3$. Hence
\[ i^{2023} = (i^4)^{505} \cdot i^3 = 1^{505} \times (-i) = -i. \]
\end{enumerate}
\end{exoresolu}

\begin{definition}[Complex Conjugate]
The \cle{conjugate} of $z = a + ib$ is $\bar{z} = a - ib$ (reflection in the real axis on the Argand diagram).
\end{definition}

\begin{propriete}[Properties of the Conjugate]
For all $z, w \in \mathbb{C}$:
\begin{align*}
\overline{z+w} &= \bar{z} + \bar{w}, &
\overline{z-w} &= \bar{z} - \bar{w}, \\
\overline{zw} &= \bar{z}\,\bar{w}, &
\overline{\left(\frac{z}{w}\right)} &= \frac{\bar{z}}{\bar{w}}\quad (w\neq 0), \\
\overline{\bar{z}} &= z, &
z + \bar{z} &= 2\operatorname{Re}(z), \\
z - \bar{z} &= 2i\operatorname{Im}(z), &
z\bar{z} &= |z|^2 = a^2 + b^2 \in \mathbb{R}_{\ge 0}.
\end{align*}
\textbf{Important:} If a polynomial has \emph{real} coefficients, its non-real roots occur in conjugate pairs.
\end{propriete}

\begin{exoresolu}[Using Conjugate Properties]
If $z = 2 - 3i$, find: (a) $\bar{z}$; (b) $z + \bar{z}$; (c) $z - \bar{z}$; (d) $z\bar{z}$.
\tcblower
\textbf{Solution.}
(a) $\bar{z} = 2 + 3i$.
(b) $z + \bar{z} = (2-3i)+(2+3i) = 4 = 2\operatorname{Re}(z)$.
(c) $z - \bar{z} = (2-3i)-(2+3i) = -6i = 2i\operatorname{Im}(z)$.
(d) $z\bar{z} = (2-3i)(2+3i) = 4 - (3i)^2 = 4 + 9 = 13 = |z|^2$.
\end{exoresolu}

\begin{methode}[Division by a Complex Number]
To write $\dfrac{z_1}{z_2}$ (with $z_2 \neq 0$) in the form $a + ib$:
\begin{enumerate}
\item Write down the \cle{conjugate} $\bar{z}_2$ of the denominator (change the sign of the imaginary part).
\item Multiply numerator \emph{and} denominator by $\bar{z}_2$.
\item The denominator becomes $z_2 \bar{z}_2 = |z_2|^2$, a positive real number.
\item Expand the numerator, replace $i^2$ by $-1$, then separate real and imaginary parts.
\end{enumerate}
\textbf{Key formula:} $\dfrac{a+ib}{c+id} = \dfrac{(a+ib)(c-id)}{c^2 + d^2}$.
\end{methode}

\begin{exoresolu}[Quotient in Cartesian Form]
Express $z = \dfrac{3 - 2i}{1 + 4i}$ in the form $a + ib$, where $a, b \in \mathbb{R}$. Hence find the real numbers $x$ and $y$ such that $\dfrac{x + iy}{1+4i} = 3 - 2i$.
\tcblower
\textbf{Solution.} The conjugate of $1 + 4i$ is $1 - 4i$. Multiply top and bottom by it:
\begin{align*}
z &= \frac{(3 - 2i)(1 - 4i)}{(1 + 4i)(1 - 4i)} = \frac{3 - 12i - 2i + 8i^2}{1^2 + 4^2} \\
&= \frac{3 - 14i - 8}{17} = \frac{-5 - 14i}{17} = -\frac{5}{17} - \frac{14}{17}i.
\end{align*}
For the second part, $x + iy = (3 - 2i)(1 + 4i) = 11 + 10i$ (computed in the previous example). Equating real and imaginary parts: $x = 11$ and $y = 10$.
\begin{attention}
Never leave a complex number in the denominator of a final answer at the GCE: always "realise the denominator" using the conjugate.
\end{attention}
\end{exoresolu}

\begin{exercice}[Practice: Algebra and Conjugates]
\begin{enumerate}
\item Simplify $(2-i)^3$ and express the result in the form $a+ib$.
\item If $z = \frac{1+2i}{3-i}$, find $\operatorname{Re}(z)$, $\operatorname{Im}(z)$, $\bar{z}$ and $|z|$.
\item Find real $x,y$ such that $(x+iy)^2 = 5 + 12i$.
\item Prove that for any $z\in\mathbb{C}$, $|z|^2 = z\bar{z}$ and deduce that $|zw| = |z||w|$.
\end{enumerate}
\end{exercice}

\coursec{Equations and Roots in \texorpdfstring{$\mathbb{C}$}{C}}

\begin{methode}[Quadratic Equations with Real Coefficients]
To solve $az^2 + bz + c = 0$ with $a, b, c \in \mathbb{R}$, $a \neq 0$:
\begin{enumerate}
\item Compute the discriminant $\Delta = b^2 - 4ac$.
\item If $\Delta \geq 0$: real roots $z = \dfrac{-b \pm \sqrt{\Delta}}{2a}$.
\item If $\Delta < 0$: write $\Delta = -|\Delta|$, so $\sqrt{\Delta} = i\sqrt{|\Delta|}$, giving two \cle{complex conjugate roots}
\[ z = \frac{-b \pm i\sqrt{|\Delta|}}{2a}. \]
\end{enumerate}
\textbf{Reflexes:} with real coefficients, non-real roots always come in conjugate pairs. Sum of roots $= -\tfrac{b}{a}$, product $= \tfrac{c}{a}$ — use these to check your answers or to form an equation from given roots.
\end{methode}

\begin{exoresolu}[GCE-Style: Solving a Quadratic in \texorpdfstring{$\mathbb{C}$}{C}]
(a) Solve the equation $z^2 - 4z + 13 = 0$, giving your answers in the form $a + ib$.

(b) Hence, or otherwise, find a quadratic equation with real coefficients whose roots are $2 + 3i$ and $2 - 3i$, and verify that it is the equation of part (a).
\tcblower
\textbf{Solution.}
(a) $\Delta = (-4)^2 - 4(1)(13) = 16 - 52 = -36 < 0$. The roots are complex:
\[ z = \frac{4 \pm \sqrt{-36}}{2} = \frac{4 \pm 6i}{2} = 2 \pm 3i. \]
The solution set is $S = \{2 - 3i;\ 2 + 3i\}$.

(b) Sum of roots: $(2+3i) + (2-3i) = 4$. Product: $(2+3i)(2-3i) = 4 - (3i)^2 = 4 + 9 = 13$. A quadratic with these roots is
\[ z^2 - (\text{sum})z + (\text{product}) = 0 \quad\Longrightarrow\quad z^2 - 4z + 13 = 0, \]
which is exactly the equation of part (a). $\blacksquare$
\end{exoresolu}

\begin{methode}[Square Root of a Complex Number (Algebraic Method)]
To find $\sqrt{w}$ where $w = a+ib$ ($a,b\in\mathbb{R}$):
\begin{enumerate}
\item Set $(x+iy)^2 = a+ib$ with $x,y\in\mathbb{R}$.
\item Expand: $x^2 - y^2 + 2ixy = a + ib$.
\item Equate real and imaginary parts:
\[ \begin{cases} x^2 - y^2 = a \\ 2xy = b \end{cases} \]
\item Use the modulus relation $x^2 + y^2 = |w| = \sqrt{a^2+b^2}$.
\item Solve for $x^2$ and $y^2$:
\[ x^2 = \frac{|w|+a}{2},\qquad y^2 = \frac{|w|-a}{2}. \]
\item Choose signs of $x$ and $y$ so that $2xy = b$ (same sign if $b>0$, opposite if $b<0$).
\end{enumerate}
\end{methode}

\begin{exoresolu}[Square Root by Algebraic Method]
Find the square roots of $5 + 12i$.
\tcblower
\textbf{Solution.} Let $(x+iy)^2 = 5+12i$. Then
\[ x^2 - y^2 = 5,\quad 2xy = 12 \Rightarrow xy = 6. \]
Also $x^2 + y^2 = |5+12i| = \sqrt{25+144} = 13$.
Adding: $2x^2 = 18 \Rightarrow x^2 = 9 \Rightarrow x = \pm 3$.
Subtracting: $2y^2 = 8 \Rightarrow y^2 = 4 \Rightarrow y = \pm 2$.
Since $xy = 6 > 0$, $x$ and $y$ have the same sign.
The two square roots are $3+2i$ and $-3-2i$.
\textbf{Check:} $(3+2i)^2 = 9 + 12i + 4i^2 = 5+12i$. $\checkmark$
\end{exoresolu}

\begin{exercice}[Practice: Equations and Square Roots]
\begin{enumerate}
\item Solve $2z^2 + 3z + 5 = 0$.
\item Find a quadratic equation with real coefficients having roots $-1+2i$ and $-1-2i$.
\item Find the square roots of $-7 + 24i$ by the algebraic method.
\item The complex number $z$ satisfies $z^2 - 2z + 10 = 0$. Find $z$ and verify that $|z|^2 = 10$.
\end{enumerate}
\end{exercice}

\coursec{The Argand Diagram, Modulus and Argument}

\begin{definition}[The Argand Diagram]
A complex number $z = x + iy$ is represented by the point $P(x,y)$ in the plane (or the vector $\vec{OP}$). The horizontal axis is the \cle{real axis}, the vertical axis is the \cle{imaginary axis}.
\end{definition}

\begin{definition}[Modulus and Argument]
For $z = x + iy \neq 0$:
\begin{itemize}
\item The \cle{modulus} is $|z| = \sqrt{x^2 + y^2} \ge 0$ (distance from origin).
\item An \cle{argument} of $z$ is any angle $\theta$ such that $\cos\theta = \frac{x}{|z|}$, $\sin\theta = \frac{y}{|z|}$.
\item The \cle{principal argument}, denoted $\arg z$ (or $\operatorname{Arg} z$), is the unique $\theta \in (-\pi, \pi]$.
\end{itemize}
\end{definition}

\begin{propriete}[Geometric Properties]
For $z, w \in \mathbb{C}$:
\begin{itemize}
\item $|z| = |\bar{z}|$, $\arg(\bar{z}) = -\arg(z)$ (mod $2\pi$).
\item $|zw| = |z||w|$, $\arg(zw) = \arg z + \arg w$ (mod $2\pi$).
\item $|z/w| = |z|/|w|$, $\arg(z/w) = \arg z - \arg w$ (mod $2\pi$).
\item $|z^n| = |z|^n$, $\arg(z^n) = n\arg z$ (mod $2\pi$).
\end{itemize}
\end{propriete}

\begin{methode}[Converting $a + ib$ to Polar Form $r(\cos\theta + i\sin\theta)$]
\begin{enumerate}
\item Compute the \cle{modulus}: $r = |z| = \sqrt{a^2 + b^2}$ (always $r \geq 0$).
\item Find an \cle{argument} $\theta$ from $\cos\theta = \dfrac{a}{r}$ and $\sin\theta = \dfrac{b}{r}$ — use \emph{both} equations to fix the quadrant.
\item Place the point on an Argand diagram if in doubt: signs of $a$ and $b$ tell you the quadrant.
\item Write $z = r(\cos\theta + i\sin\theta) = [r; \theta]$, usually with the principal argument $\theta \in (-\pi, \pi]$.
\end{enumerate}
\textbf{Key values to know:}
\begin{center}
\begin{tabular}{|c|c|c|c|}\hline
$z$ & $r$ & $\theta$ (principal) & Polar form \\ \hline
$1 + i$ & $\sqrt{2}$ & $\pi/4$ & $\sqrt{2}(\cos\frac{\pi}{4}+i\sin\frac{\pi}{4})$ \\ \hline
$1 + i\sqrt{3}$ & $2$ & $\pi/3$ & $2(\cos\frac{\pi}{3}+i\sin\frac{\pi}{3})$ \\ \hline
$-\sqrt{3} + i$ & $2$ & $5\pi/6$ & $2(\cos\frac{5\pi}{6}+i\sin\frac{5\pi}{6})$ \\ \hline
$-1 - i$ & $\sqrt{2}$ & $-3\pi/4$ & $\sqrt{2}(\cos(-\frac{3\pi}{4})+i\sin(-\frac{3\pi}{4}))$ \\ \hline
\end{tabular}
\end{center}
\end{methode}

\begin{exoresolu}[Cartesian to Polar Form]
Write $z = -1 + i\sqrt{3}$ in polar form $r(\cos\theta + i\sin\theta)$ with $-\pi < \theta \leq \pi$. Deduce the polar form of $\bar{z}$ and of $-z$.
\tcblower
\textbf{Solution.} Here $a = -1$, $b = \sqrt{3}$.
\[ r = \sqrt{(-1)^2 + (\sqrt{3})^2} = \sqrt{1 + 3} = 2. \]
Then $\cos\theta = \dfrac{-1}{2}$ and $\sin\theta = \dfrac{\sqrt{3}}{2}$. Since $\sin\theta > 0$ and $\cos\theta < 0$, $\theta$ lies in the second quadrant: $\theta = \pi - \dfrac{\pi}{3} = \dfrac{2\pi}{3}$. Hence
\[ z = 2\left(\cos\frac{2\pi}{3} + i\sin\frac{2\pi}{3}\right). \]
For the conjugate, $\arg(\bar{z}) = -\arg(z)$ and $|\bar{z}| = |z|$:
\[ \bar{z} = 2\left(\cos\frac{2\pi}{3} - i\sin\frac{2\pi}{3}\right) = 2\left(\cos\left(-\tfrac{2\pi}{3}\right) + i\sin\left(-\tfrac{2\pi}{3}\right)\right). \]
For $-z$: multiplying by $-1 = e^{i\pi}$ adds $\pi$ to the argument; reducing modulo $2\pi$, $\arg(-z) = \tfrac{2\pi}{3} - \pi = -\tfrac{\pi}{3}$, so
\[ -z = 2\left(\cos\left(-\tfrac{\pi}{3}\right) + i\sin\left(-\tfrac{\pi}{3}\right)\right). \]
(Check: $-z = 1 - i\sqrt{3}$, indeed in the fourth quadrant with argument $-\pi/3$.)
\end{exoresolu}

\begin{exercice}[Practice: Modulus, Argument, Polar Form]
\begin{enumerate}
\item Find the modulus and principal argument of: (a) $-2-2i$; (b) $3-3i\sqrt{3}$; (c) $-4i$.
\item Write in polar form $r(\cos\theta+i\sin\theta)$ with $\theta\in(-\pi,\pi]$: (a) $-\sqrt{3}-i$; (b) $2-2i$.
\item If $z = 2(\cos\frac{\pi}{5}+i\sin\frac{\pi}{5})$, find the polar form of $\bar{z}$, $-z$, $z^2$ and $1/z$.
\end{enumerate}
\end{exercice}

\coursec{Polar and Exponential Forms}

\begin{definition}[Exponential Form]
Euler's formula states $e^{i\theta} = \cos\theta + i\sin\theta$ for all $\theta\in\mathbb{R}$.
A complex number $z \neq 0$ can be written as
\[ z = r e^{i\theta}, \quad\text{where } r = |z| > 0,\ \theta = \arg z. \]
\end{definition}

\begin{propriete}[Euler's Identities]
\begin{align*}
\cos\theta &= \frac{e^{i\theta} + e^{-i\theta}}{2}, &
\sin\theta &= \frac{e^{i\theta} - e^{-i\theta}}{2i}. \\
e^{i\pi} &= -1 \quad\text{(Euler's identity)}, &
e^{i\pi/2} &= i.
\end{align*}
\end{propriete}

\begin{methode}[Converting Between Forms]
\begin{itemize}
\item Cartesian $\to$ Exponential: $r = \sqrt{a^2+b^2}$, $\theta = \arg(a+ib)$, then $z = re^{i\theta}$.
\item Exponential $\to$ Cartesian: $z = r(\cos\theta + i\sin\theta) = r\cos\theta + ir\sin\theta$.
\item Polar $\leftrightarrow$ Exponential: just replace $\cos\theta+i\sin\theta$ by $e^{i\theta}$.
\end{itemize}
\end{methode}

\begin{exoresolu}[Exponential Form and Operations]
Let $z_1 = 2e^{i\pi/3}$ and $z_2 = \sqrt{2}e^{-i\pi/4}$. Find in exponential form: (a) $z_1z_2$; (b) $z_1/z_2$; (c) $z_1^3$; (d) $\sqrt{z_1}$ (both roots).
\tcblower
\textbf{Solution.}
(a) $z_1z_2 = 2\sqrt{2}\, e^{i(\pi/3 - \pi/4)} = 2\sqrt{2}\, e^{i\pi/12}$.
(b) $z_1/z_2 = \frac{2}{\sqrt{2}}\, e^{i(\pi/3 + \pi/4)} = \sqrt{2}\, e^{i7\pi/12}$.
(c) $z_1^3 = 2^3 e^{i\pi} = 8 e^{i\pi} = -8$.
(d) $\sqrt{z_1} = \sqrt{2}\, e^{i(\pi/6 + k\pi)}$, $k=0,1$:
$k=0$: $\sqrt{2}e^{i\pi/6}$; $k=1$: $\sqrt{2}e^{i7\pi/6}$.
\end{exoresolu}

\coursec{De Moivre's Theorem and Trigonometric Applications}

\begin{propriete}[De Moivre's Theorem]
For any integer $n$ and any real $\theta$:
\[ (\cos\theta + i\sin\theta)^n = \cos n\theta + i\sin n\theta. \]
Equivalently, $(e^{i\theta})^n = e^{in\theta}$.
\end{propriete}

\begin{methode}[Computing $z^n$ via De Moivre]
To compute a high power of a complex number:
\begin{enumerate}
\item Convert $z$ to polar (or exponential) form: $z = re^{i\theta}$.
\item Apply \cle{De Moivre's theorem}: $z^n = r^n e^{in\theta} = r^n(\cos n\theta + i\sin n\theta)$.
\item Reduce $n\theta$ modulo $2\pi$ to a familiar angle.
\item Convert back to Cartesian form if the question demands it.
\end{enumerate}
\textbf{Reflex:} never expand $(a+ib)^n$ directly for large $n$ — the polar route is far shorter and is what examiners expect.
\end{methode}

\begin{exoresolu}[Power via De Moivre]
Given $z = 1 + i$, express $z$ in exponential form and hence show that $z^{12} = -64$. Give the answer in Cartesian form.
\tcblower
\textbf{Solution.} $|z| = \sqrt{1^2 + 1^2} = \sqrt{2}$, and $\cos\theta = \sin\theta = \tfrac{1}{\sqrt{2}}$ gives $\theta = \tfrac{\pi}{4}$. Thus
\[ z = \sqrt{2}\, e^{i\pi/4}. \]
By De Moivre's theorem:
\[ z^{12} = (\sqrt{2})^{12} e^{i \cdot 12\pi/4} = 2^{6} e^{3i\pi} = 64\, e^{3i\pi}. \]
Since $e^{3i\pi} = \cos 3\pi + i\sin 3\pi = -1 + 0i = -1$ (as $3\pi \equiv \pi \pmod{2\pi}$):
\[ z^{12} = 64 \times (-1) = -64. \]
In Cartesian form: $z^{12} = -64 + 0i = -64$, a real number. $\blacksquare$
\end{exoresolu}

\begin{methode}[Trigonometric Identities from De Moivre]
\textbf{(A) Multiple angles from powers (expand and equate):}
Expand $(\cos\theta + i\sin\theta)^n$ by the binomial theorem and equate real and imaginary parts with $\cos n\theta + i\sin n\theta$.

\textbf{(B) Powers from multiple angles (binomial with $z$ and $z^{-1}$):}
Set $z = \cos\theta + i\sin\theta$, so $z^n + z^{-n} = 2\cos n\theta$ and $z^n - z^{-n} = 2i\sin n\theta$. Then:
\begin{enumerate}
\item Write $\cos^n\theta = \left(\dfrac{z + z^{-1}}{2}\right)^n$ (or with $z - z^{-1}$ for $\sin$).
\item Expand the binomial and pair terms $z^k$ with $z^{-k}$.
\item Replace each pair by $2\cos k\theta$ (or $2i\sin k\theta$).
\end{enumerate}
\end{methode}

\begin{exoresolu}[Expressing $\cos^4\theta$ in Multiple Angles]
Using De Moivre's theorem, prove that
\[ \cos^4\theta = \frac{1}{8}\bigl(\cos 4\theta + 4\cos 2\theta + 3\bigr). \]
\tcblower
\textbf{Solution.} Let $z = \cos\theta + i\sin\theta$. Then $z + z^{-1} = 2\cos\theta$ and $z^n + z^{-n} = 2\cos n\theta$. We have $\cos\theta = \dfrac{z + z^{-1}}{2}$, so
\begin{align*}
(2\cos\theta)^4 &= (z + z^{-1})^4 = z^4 + 4z^2 + 6 + 4z^{-2} + z^{-4} \\
&= (z^4 + z^{-4}) + 4(z^2 + z^{-2}) + 6 \\
&= 2\cos 4\theta + 4(2\cos 2\theta) + 6.
\end{align*}
Therefore $16\cos^4\theta = 2\cos 4\theta + 8\cos 2\theta + 6$, and dividing by $16$:
\[ \cos^4\theta = \frac{2\cos 4\theta + 8\cos 2\theta + 6}{16} = \frac{1}{8}\bigl(\cos 4\theta + 4\cos 2\theta + 3\bigr). \quad\blacksquare \]
\textbf{Remark:} such identities make integration easy: $\displaystyle\int \cos^4\theta\, d\theta = \tfrac{1}{32}\sin 4\theta + \tfrac{1}{4}\sin 2\theta + \tfrac{3}{8}\theta + C$.
\end{exoresolu}

\begin{exoresolu}[Expressing $\sin^3\theta$ in Multiple Angles]
Express $\sin^3\theta$ in terms of sines of multiple angles.
\tcblower
\textbf{Solution.} Let $z = \cos\theta + i\sin\theta$. Then $z - z^{-1} = 2i\sin\theta$, so $\sin\theta = \dfrac{z - z^{-1}}{2i}$.
\begin{align*}
(2i\sin\theta)^3 &= (z - z^{-1})^3 = z^3 - 3z + 3z^{-1} - z^{-3} \\
&= (z^3 - z^{-3}) - 3(z - z^{-1}) \\
&= 2i\sin 3\theta - 3(2i\sin\theta).
\end{align*}
Thus $-8i\sin^3\theta = 2i\sin 3\theta - 6i\sin\theta$. Dividing by $-8i$:
\[ \sin^3\theta = \frac{1}{4}(3\sin\theta - \sin 3\theta). \quad\blacksquare \]
\end{exoresolu}

\begin{exoresolu}[Expressing $\cos 5\theta$ as a Polynomial in $\cos\theta$]
Use De Moivre's theorem to show that $\cos 5\theta = 16\cos^5\theta - 20\cos^3\theta + 5\cos\theta$.
\tcblower
\textbf{Solution.} By De Moivre: $\cos 5\theta + i\sin 5\theta = (\cos\theta + i\sin\theta)^5$.
Expand by binomial theorem:
\begin{align*}
(\cos\theta + i\sin\theta)^5 &= \cos^5\theta + 5i\cos^4\theta\sin\theta - 10\cos^3\theta\sin^2\theta \\
&\quad - 10i\cos^2\theta\sin^3\theta + 5\cos\theta\sin^4\theta + i\sin^5\theta.
\end{align*}
Equate real parts:
\[ \cos 5\theta = \cos^5\theta - 10\cos^3\theta\sin^2\theta + 5\cos\theta\sin^4\theta. \]
Replace $\sin^2\theta = 1-\cos^2\theta$:
\begin{align*}
\cos 5\theta &= \cos^5\theta - 10\cos^3\theta(1-\cos^2\theta) + 5\cos\theta(1-\cos^2\theta)^2 \\
&= \cos^5\theta - 10\cos^3\theta + 10\cos^5\theta + 5\cos\theta(1 - 2\cos^2\theta + \cos^4\theta) \\
&= 16\cos^5\theta - 20\cos^3\theta + 5\cos\theta. \quad\blacksquare
\end{align*}
\end{exoresolu}

\begin{exercice}[Practice: De Moivre and Trigonometry]
\begin{enumerate}
\item Use De Moivre to express $\cos 3\theta$ and $\sin 3\theta$ in powers of $\cos\theta$ and $\sin\theta$.
\item Express $\cos^5\theta$ in terms of cosines of multiple angles.
\item Express $\sin^4\theta$ in terms of cosines of multiple angles.
\item Hence evaluate $\displaystyle\int_0^{\pi/2} \cos^5\theta\, d\theta$ and $\displaystyle\int_0^{\pi/2} \sin^4\theta\, d\theta$.
\end{enumerate}
\end{exercice}

\coursec{nth Roots of Complex Numbers and Cube Roots of Unity}

\begin{methode}[Solving $z^n = w$ (nth Roots)]
\begin{enumerate}
\item Write $w$ in exponential form: $w = Re^{i\varphi}$ (with $\varphi = \arg w$).
\item The $n$ distinct roots are
\[ z_k = R^{1/n}\, e^{i(\varphi + 2k\pi)/n}, \qquad k = 0, 1, 2, \dots, n-1. \]
\item All roots have the \emph{same modulus} $R^{1/n}$ and their arguments differ by $\dfrac{2\pi}{n}$: they are the vertices of a \cle{regular $n$-gon} on the circle of radius $R^{1/n}$.
\item For \textbf{square roots} ($n = 2$) one may also set $(x + iy)^2 = w$ and equate real and imaginary parts, using $x^2 + y^2 = |w|$.
\end{enumerate}
\end{methode}

\begin{exoresolu}[Cube Roots of a Complex Number]
Find all the solutions of $z^3 = -8i$, giving each in the form $a + ib$. Show the roots on an Argand diagram.
\tcblower
\textbf{Solution.} Write $-8i$ in exponential form: $|-8i| = 8$ and $\arg(-8i) = -\tfrac{\pi}{2}$, so $-8i = 8e^{-i\pi/2}$. The cube roots are
\[ z_k = 8^{1/3}\, e^{i(-\pi/2 + 2k\pi)/3} = 2\, e^{i(-\pi/6 + 2k\pi/3)}, \quad k = 0, 1, 2. \]
\begin{itemize}
\item $k = 0$: $z_0 = 2e^{-i\pi/6} = 2\left(\cos\tfrac{\pi}{6} - i\sin\tfrac{\pi}{6}\right) = \sqrt{3} - i$.
\item $k = 1$: $z_1 = 2e^{i\pi/2} = 2i$.
\item $k = 2$: $z_2 = 2e^{i7\pi/6} = 2\left(-\cos\tfrac{\pi}{6} - i\sin\tfrac{\pi}{6}\right) = -\sqrt{3} - i$.
\end{itemize}
\textbf{Check:} $(2i)^3 = 8i^3 = -8i$. $\checkmark$ The three roots lie on the circle of radius $2$, at the vertices of an equilateral triangle.
\begin{popfigure}
\begin{tikzpicture}[scale=1.1]
\draw[->, popmuted] (-2.6,0) -- (2.6,0) node[below left] {\small Re};
\draw[->, popmuted] (0,-2.6) -- (0,2.6) node[below left] {\small Im};
\draw[popline] (0,0) circle (2);
\draw[->, popblue, thick] (0,0) -- (1.732,-1) node[below right] {\small $z_0$};
\draw[->, poporange, thick] (0,0) -- (0,2) node[above right] {\small $z_1$};
\draw[->, popgreen, thick] (0,0) -- (-1.732,-1) node[below left] {\small $z_2$};
\foreach \p in {(1.732,-1),(0,2),(-1.732,-1)} \fill[popdark] \p circle (1.5pt);
\end{tikzpicture}
\legende{Cube roots of $-8i$ on the Argand diagram.}
\end{popfigure}
\end{exoresolu}

\begin{definition}[Cube Roots of Unity]
The solutions of $z^3 = 1$ are called the \cle{cube roots of unity}:
\[ 1,\quad \omega = e^{2i\pi/3} = \frac{-1 + i\sqrt{3}}{2},\quad \omega^2 = e^{4i\pi/3} = \frac{-1 - i\sqrt{3}}{2} = \bar{\omega}. \]
They satisfy the fundamental properties:
\[ 1 + \omega + \omega^2 = 0,\qquad \omega^3 = 1,\qquad \omega^2 = \bar{\omega} = \frac{1}{\omega}. \]
\end{definition}

\begin{exoresolu}[Using Cube Roots of Unity]
Simplify: (a) $(1+\omega)^3$; (b) $(1-\omega)(1-\omega^2)$; (c) $\omega^{2024} + \omega^{2025} + \omega^{2026}$.
\tcblower
\textbf{Solution.}
(a) Since $1+\omega+\omega^2=0$, we have $1+\omega = -\omega^2$. Then $(1+\omega)^3 = (-\omega^2)^3 = -\omega^6 = -(\omega^3)^2 = -1$.
(b) $(1-\omega)(1-\omega^2) = 1 - \omega - \omega^2 + \omega^3 = 1 - (\omega+\omega^2) + 1 = 1 - (-1) + 1 = 3$.
(c) $\omega^{2024} + \omega^{2025} + \omega^{2026} = \omega^{2024}(1+\omega+\omega^2) = \omega^{2024}\times 0 = 0$.
(Alternatively, reduce exponents mod 3: $2024\equiv2$, $2025\equiv0$, $2026\equiv1$, sum $=\omega^2+1+\omega=0$.)
\end{exoresolu}

\begin{exoresolu}[Fourth Roots of a Complex Number]
Find the fourth roots of $-16$, giving each in the form $a+ib$. Plot them on an Argand diagram.
\tcblower
\textbf{Solution.} $-16 = 16 e^{i\pi}$. The fourth roots are
\[ z_k = 16^{1/4} e^{i(\pi + 2k\pi)/4} = 2 e^{i(\pi/4 + k\pi/2)},\quad k=0,1,2,3. \]
\begin{itemize}
\item $k=0$: $z_0 = 2e^{i\pi/4} = \sqrt{2} + i\sqrt{2}$.
\item $k=1$: $z_1 = 2e^{i3\pi/4} = -\sqrt{2} + i\sqrt{2}$.
\item $k=2$: $z_2 = 2e^{i5\pi/4} = -\sqrt{2} - i\sqrt{2}$.
\item $k=3$: $z_3 = 2e^{i7\pi/4} = \sqrt{2} - i\sqrt{2}$.
\end{itemize}
These are the vertices of a square centred at the origin, radius $2$.
\begin{popfigure}
\begin{tikzpicture}[scale=0.9]
\draw[->, popmuted] (-3,0) -- (3,0) node[below left] {\small Re};
\draw[->, popmuted] (0,-3) -- (0,3) node[below left] {\small Im};
\draw[popline] (0,0) circle (2);
\foreach \ang/\col/\lab in {45/popblue/z_0, 135/poporange/z_1, 225/popgreen/z_2, 315/poppurple/z_3} {
  \draw[->, \col, thick] (0,0) -- (\ang:2);
  \fill[\col] (\ang:2) circle (2pt) node[shift=(\ang:0.4)] {\small $\lab$};
}
\end{tikzpicture}
\legende{Fourth roots of $-16$.}
\end{popfigure}
\end{exoresolu}

\begin{exercice}[Practice: nth Roots and Cube Roots of Unity]
\begin{enumerate}
\item Find all the fifth roots of $32$, giving answers in polar form $re^{i\theta}$.
\item Solve $z^4 = -16i$, expressing roots in the form $a+ib$.
\item If $\omega$ is a complex cube root of unity, evaluate: (a) $\omega^{100} + \omega^{200}$; (b) $(1-\omega+\omega^2)^6$.
\item Find the square roots of $3+4i$ by both the algebraic method and the polar method. Verify they agree.
\end{enumerate}
\end{exercice}

\coursec{Loci in the Complex Plane and Chapter Synthesis}

\begin{methode}[Identifying and Sketching a Locus]
Translate the condition on $z$ into geometry, writing $z = x + iy$ if needed:
\begin{enumerate}
\item $|z - a| = r$: \cle{circle}, centre $a$, radius $r$.
\item $|z - a| = |z - b|$: \cle{perpendicular bisector} of the segment joining $a$ and $b$.
\item $\arg(z - a) = \theta$: \cle{half-line} (ray) from $a$ (point $a$ excluded) in direction $\theta$.
\item $|z - a| = k\,|z - b|$ with $k > 0$, $k \neq 1$: circle of Apollonius — substitute $z = x + iy$ and rearrange to recognise centre and radius.
\item $\arg\left(\frac{z-a}{z-b}\right) = \theta$: arc of a circle (or line if $\theta = 0,\pi$).
\end{enumerate}
\textbf{Reflex:} for inequalities, decide whether the boundary is included ($\leq$, solid line) or excluded ($<$, dashed line), and shade the correct region.
\end{methode}

\begin{exoresolu}[GCE-Style: Two Loci and Their Intersection]
The complex number $z$ satisfies $|z - 2 - i| = 2$.
\begin{enumerate}
\item[(a)] Describe and sketch the locus of the point $P$ representing $z$ in the Argand diagram.
\item[(b)] Find the Cartesian equation of the locus.
\item[(c)] A second condition $\arg(z - 2 - i) = \dfrac{\pi}{4}$ is added. Find the complex number(s) satisfying both conditions.
\end{enumerate}
\tcblower
\textbf{Solution.}
(a) $|z - (2 + i)| = 2$ says the distance from $P(z)$ to the fixed point $A(2, 1)$ is constantly $2$: the locus is the \textbf{circle} with centre $(2, 1)$ and radius $2$.

(b) Set $z = x + iy$. Then $z - 2 - i = (x - 2) + i(y - 1)$, so
\[ |z - 2 - i|^2 = (x-2)^2 + (y-1)^2 = 4. \]
This is the Cartesian equation of the circle.

(c) The condition $\arg(z - 2 - i) = \tfrac{\pi}{4}$ describes the half-line from $A(2,1)$ making an angle $\tfrac{\pi}{4}$ with the positive real axis (the point $A$ itself excluded). Points on it have the form
\[ z - (2 + i) = t e^{i\pi/4} = t\left(\tfrac{\sqrt{2}}{2} + i\tfrac{\sqrt{2}}{2}\right), \qquad t > 0. \]
Imposing $|z - 2 - i| = 2$ forces $t = 2$. Hence
\[ z = 2 + i + 2\left(\tfrac{\sqrt{2}}{2} + i\tfrac{\sqrt{2}}{2}\right) = (2 + \sqrt{2}) + i(1 + \sqrt{2}). \]
There is exactly \textbf{one} complex number satisfying both conditions: $z = (2 + \sqrt{2}) + (1 + \sqrt{2})i$.
\begin{popfigure}
\begin{tikzpicture}[scale=0.85]
\draw[->, popmuted] (-0.6,0) -- (5.4,0) node[below left] {\small Re};
\draw[->, popmuted] (0,-0.6) -- (0,4.2) node[below left] {\small Im};
\draw[popblue, thick] (2,1) circle (2);
\draw[poporange, thick, ->] (2,1) -- (4.6,3.6);
\fill[popdark] (2,1) circle (1.8pt) node[below left] {\small $A(2,1)$};
\fill[poppurple] (3.414,2.414) circle (2pt) node[above right] {\small $P$};
\draw[popline, dashed] (2,1) -- (3.6,1);
\node[popmuted] at (3.05,0.72) {\small $\pi/4$};
\end{tikzpicture}
\legende{Intersection of circle $|z-2-i|=2$ and ray $\arg(z-2-i)=\pi/4$.}
\end{popfigure}
\end{exoresolu}

\begin{exoresolu}[Locus: Perpendicular Bisector and Circle of Apollonius]
(a) Sketch the locus $|z - 3| = |z + 1|$.
(b) Find the Cartesian equation of $|z - 1| = 2|z + 3|$ and identify the locus.
\tcblower
\textbf{Solution.}
(a) $|z-3| = |z+1|$ means the distance from $P$ to $A(3,0)$ equals the distance to $B(-1,0)$. The locus is the perpendicular bisector of $AB$: the vertical line $x = 1$ (the imaginary axis shifted to $x=1$).

(b) Let $z = x+iy$. Then $|z-1|^2 = (x-1)^2+y^2$, $|z+3|^2 = (x+3)^2+y^2$.
The equation $|z-1| = 2|z+3|$ squared gives:
\[ (x-1)^2 + y^2 = 4\bigl((x+3)^2 + y^2\bigr). \]
Expand: $x^2 - 2x + 1 + y^2 = 4x^2 + 24x + 36 + 4y^2$.
Bring all terms: $0 = 3x^2 + 26x + 35 + 3y^2$.
Divide by $3$: $x^2 + \frac{26}{3}x + y^2 = -\frac{35}{3}$.
Complete the square: $\left(x + \frac{13}{3}\right)^2 + y^2 = \frac{169}{9} - \frac{105}{9} = \frac{64}{9}$.
This is a \textbf{circle} with centre $\left(-\frac{13}{3}, 0\right)$ and radius $\frac{8}{3}$.
\end{exoresolu}

\begin{exoresolu}[Locus with Inequality: Shading Regions]
Sketch the region defined by $|z - 2i| \le 3$ and $\arg(z) \ge \frac{\pi}{6}$.
\tcblower
\textbf{Solution.}
\begin{itemize}
\item $|z - 2i| \le 3$: closed disc centred at $(0,2)$ with radius $3$ (boundary included, solid circle).
\item $\arg(z) \ge \frac{\pi}{6}$: region on and "above" the ray from origin at angle $\pi/6$ (boundary included, solid ray).
\end{itemize}
The required region is the intersection of the disc and the angular sector.
\begin{popfigure}
\begin{tikzpicture}[scale=0.7]
\draw[->, popmuted] (-4,0) -- (4,0) node[below left] {\small Re};
\draw[->, popmuted] (0,-1) -- (0,5.5) node[below left] {\small Im};
\draw[popblue, thick, fill=popblueL, opacity=0.3] (0,2) circle (3);
\draw[poporange, thick, ->] (0,0) -- (4.5,2.6);
\draw[poporange, thick, fill=poporangeL, opacity=0.2] (0,0) -- (4.5,2.6) -- (4.5,5) -- (0,5) -- cycle;
\fill[popdark] (0,2) circle (1.5pt) node[left] {\small $2i$};
\draw[popline, dashed] (0,0) -- (3.6,0);
\node[popmuted] at (1.2,0.3) {\small $\pi/6$};
\end{tikzpicture}
\legende{Region: $|z-2i|\le 3$ and $\arg z \ge \pi/6$.}
\end{popfigure}
\end{exoresolu}

\begin{aretenir}[Chapter Synthesis: Complex Numbers — Key Facts for GCE]
\begin{itemize}
\item \textbf{Cartesian form:} $z = a+ib$, $a,b\in\mathbb{R}$. Equality: $a+ib=c+id \iff a=c,\ b=d$.
\item \textbf{Conjugate:} $\bar{z}=a-ib$; $z\bar{z}=|z|^2$; $\overline{z\pm w}=\bar{z}\pm\bar{w}$, $\overline{zw}=\bar{z}\bar{w}$.
\item \textbf{Modulus/Argument:} $|z|=\sqrt{a^2+b^2}$; $\arg z\in(-\pi,\pi]$ from $\cos\theta=a/|z|$, $\sin\theta=b/|z|$.
\item \textbf{Polar/Exponential:} $z = r(\cos\theta+i\sin\theta) = re^{i\theta}$.
\item \textbf{De Moivre:} $(\cos\theta+i\sin\theta)^n = \cos n\theta+i\sin n\theta$; $z^n = r^n e^{in\theta}$.
\item \textbf{Trig identities:} $\cos\theta=\frac{e^{i\theta}+e^{-i\theta}}{2}$, $\sin\theta=\frac{e^{i\theta}-e^{-i\theta}}{2i}$.
\item \textbf{nth roots:} $z^n = Re^{i\varphi} \Rightarrow z_k = R^{1/n}e^{i(\varphi+2k\pi)/n}$, $k=0,\dots,n-1$.
\item \textbf{Cube roots of unity:} $1,\ \omega=\frac{-1+i\sqrt{3}}{2},\ \omega^2=\bar{\omega}$; $1+\omega+\omega^2=0$, $\omega^3=1$.
\item \textbf{Loci:} $|z-a|=r$ (circle); $|z-a|=|z-b|$ (perp. bisector); $\arg(z-a)=\theta$ (half-line from $a$); $|z-a|=k|z-b|$ (circle, $k\neq1$).
\end{itemize}
\end{aretenir}

\begin{attention}[Frequent Errors at the GCE]
\begin{itemize}
\item Forgetting that $\sqrt{-9} = 3i$, not $-3$; and writing $\sqrt{a}\sqrt{b} = \sqrt{ab}$ for negative numbers (false in $\mathbb{C}$!).
\item Giving the argument in the wrong quadrant: always use \emph{both} $\cos\theta$ and $\sin\theta$.
\item In De Moivre, forgetting to raise the modulus to the power $n$: $(re^{i\theta})^n = r^n e^{in\theta}$.
\item For $z^n = w$, giving only one root: there are always exactly $n$ distinct roots.
\item In locus questions, including the centre point on a ray $\arg(z - a) = \theta$: the argument is undefined at $z = a$, so $a$ is excluded.
\item Confusing $\operatorname{Im}(z)$ (a real number) with the imaginary part as a multiple of $i$.
\item Not reducing arguments to the principal range $(-\pi,\pi]$ in final answers.
\end{itemize}
\end{attention}

\begin{exercice}[Mixed GCE-Style Revision Exercises]
\begin{enumerate}
\item Given $z = \frac{2-i}{1+2i}$, express $z$ in the form $a+ib$. Hence find $|z|$ and $\arg z$.
\item Solve the equation $z^2 + 2z + 5 = 0$. If the roots are $\alpha$ and $\beta$, evaluate $\alpha^2 + \beta^2$ without finding $\alpha,\beta$ explicitly.
\item (a) Express $1+i\sqrt{3}$ in the form $re^{i\theta}$ with $-\pi<\theta\le\pi$.
(b) Hence find $(1+i\sqrt{3})^6$ in Cartesian form.
\item Find the four fourth roots of $-16$, giving each in the form $a+ib$. Show them on an Argand diagram.
\item The complex number $z$ satisfies $|z-1| = |z-2i|$. Find the Cartesian equation of the locus of $z$ and describe it geometrically.
\item (a) Show that $\cos 4\theta = 8\cos^4\theta - 8\cos^2\theta + 1$.
(b) Hence solve $\cos 4\theta = \frac{1}{2}$ for $0 \le \theta \le \pi$.
\item If $\omega$ is a complex cube root of unity, prove that $(1-\omega+\omega^2)^4 = 16\omega$.
\item Sketch on a single Argand diagram the region satisfying both $|z-3| \le 2$ and $\frac{\pi}{4} \le \arg(z) \le \frac{\pi}{2}$.
\end{enumerate}
\end{exercice}

\begin{corrige}[Exercise 1]
$z = \frac{(2-i)(1-2i)}{1^2+2^2} = \frac{2-4i-i+2i^2}{5} = \frac{2-5i-2}{5} = -i$. So $a=0$, $b=-1$. $|z|=1$, $\arg z = -\frac{\pi}{2}$.
\end{corrige}

\begin{corrige}[Exercise 2]
$\Delta = 4 - 20 = -16$. Roots: $z = \frac{-2\pm 4i}{2} = -1\pm 2i$.
$\alpha+\beta = -2$, $\alpha\beta = 5$. $\alpha^2+\beta^2 = (\alpha+\beta)^2 - 2\alpha\beta = 4 - 10 = -6$.
\end{corrige}

\begin{corrige}[Exercise 3]
(a) $r=2$, $\theta=\pi/3$, so $1+i\sqrt{3} = 2e^{i\pi/3}$.
(b) $(2e^{i\pi/3})^6 = 64 e^{2i\pi} = 64$.
\end{corrige}

\begin{corrige}[Exercise 4]
$-16 = 16e^{i\pi}$. Roots: $2e^{i(\pi+2k\pi)/4}$, $k=0,1,2,3$:
$k=0$: $\sqrt{2}+i\sqrt{2}$; $k=1$: $-\sqrt{2}+i\sqrt{2}$; $k=2$: $-\sqrt{2}-i\sqrt{2}$; $k=3$: $\sqrt{2}-i\sqrt{2}$.
\end{corrige}

\begin{corrige}[Exercise 5]
$|z-1|^2 = |z-2i|^2 \Rightarrow (x-1)^2+y^2 = x^2+(y-2)^2 \Rightarrow x^2-2x+1+y^2 = x^2+y^2-4y+4 \Rightarrow -2x+1 = -4y+4 \Rightarrow 2x-4y+3=0$. Perpendicular bisector of segment joining $(1,0)$ and $(0,2)$.
\end{corrige}

\begin{corrige}[Exercise 6]
(a) $\cos 4\theta + i\sin 4\theta = (\cos\theta+i\sin\theta)^4 = \cos^4\theta + 4i\cos^3\theta\sin\theta - 6\cos^2\theta\sin^2\theta - 4i\cos\theta\sin^3\theta + \sin^4\theta$.
Real part: $\cos 4\theta = \cos^4\theta - 6\cos^2\theta\sin^2\theta + \sin^4\theta = \cos^4\theta - 6\cos^2\theta(1-\cos^2\theta) + (1-\cos^2\theta)^2 = 8\cos^4\theta - 8\cos^2\theta + 1$.
(b) $8\cos^4\theta - 8\cos^2\theta + 1 = \frac{1}{2} \Rightarrow 16\cos^4\theta - 16\cos^2\theta + 1 = 0$. Let $u=\cos^2\theta$: $16u^2-16u+1=0 \Rightarrow u = \frac{16\pm\sqrt{256-64}}{32} = \frac{16\pm8\sqrt{3}}{32} = \frac{2\pm\sqrt{3}}{4}$.
Since $\cos^2\theta \le 1$, both are valid. $\cos\theta = \pm\sqrt{\frac{2\pm\sqrt{3}}{4}}$. For $0\le\theta\le\pi$, find corresponding $\theta$ values.
\end{corrige}

\begin{corrige}[Exercise 7]
$1-\omega+\omega^2 = (1+\omega+\omega^2) - 2\omega = -2\omega$. Then $(-2\omega)^4 = 16\omega^4 = 16\omega^3\omega = 16\omega$.
\end{corrige}

\begin{corrige}[Exercise 8]
$|z-3|\le 2$: closed disc centre $(3,0)$ radius $2$.
$\pi/4 \le \arg z \le \pi/2$: sector between rays at $45^\circ$ and $90^\circ$ (boundaries included).
Intersection is the part of the disc in the first quadrant above the line $y=x$.
\end{corrige}

\newpage

\coursec{Exercises}

\courssub{I check my knowledge}

\begin{exercice}[Multiple choice questions]
For each question, choose the single correct answer.
\begin{enumerate}
\item The modulus of $z = 3 - 4i$ is:
\begin{enumerate}
\item $1$ \item $5$ \item $7$ \item $25$
\end{enumerate}
\item If $z = -2 + 2i$, then a principal argument of $z$ is:
\begin{enumerate}
\item $\dfrac{\pi}{4}$ \item $-\dfrac{\pi}{4}$ \item $\dfrac{3\pi}{4}$ \item $-\dfrac{3\pi}{4}$
\end{enumerate}
\item The product $(2 + i)(2 - i)$ is equal to:
\begin{enumerate}
\item $3$ \item $5$ \item $4 + i$ \item $3 + 4i$
\end{enumerate}
\item The cube roots of unity are:
\begin{enumerate}
\item $1,\ i,\ -i$ \item $1,\ \dfrac{-1 \pm i\sqrt{3}}{2}$ \item $1,\ -1,\ i$ \item $1,\ \dfrac{1 \pm i\sqrt{3}}{2}$
\end{enumerate}
\end{enumerate}
\end{exercice}

\begin{exercice}[True or false — justify]
State whether each statement is true or false, and justify your answer briefly.
\begin{enumerate}
\item Every real number is a complex number.
\item If $z$ is a complex number, then $z\bar{z} = |z|^2$.
\item The equation $z^2 + 1 = 0$ has no solution in $\mathbb{C}$.
\item If $\arg(z) = \dfrac{\pi}{2}$, then $z$ is purely imaginary with positive imaginary part.
\end{enumerate}
\end{exercice}

\begin{exercice}[Definitions]
\begin{enumerate}
\item Define the \cle{modulus} and the \cle{principal argument} of a non-zero complex number $z = a + ib$.
\item Define the \cle{complex conjugate} of $z = a + ib$ and state, without proof, the value of $z + \bar{z}$ and $z - \bar{z}$ in terms of $a$ and $b$.
\item State \cle{De Moivre's theorem} for an integer $n$.
\end{enumerate}
\end{exercice}

\begin{exercice}[Direct calculations]
Let $z_1 = 2 + 3i$ and $z_2 = 1 - i$. Compute, giving each answer in the form $a + ib$:
\begin{enumerate}
\item $z_1 + z_2$ and $z_1 - 2z_2$;
\item $z_1 z_2$;
\item $\dfrac{z_1}{z_2}$;
\item $i^{2025}$.
\end{enumerate}
\end{exercice}

\courssub{I practise}

\begin{exercice}[Cartesian and exponential forms]
Given $z_1 = 1 + i\sqrt{3}$ and $z_2 = \sqrt{2}\,e^{i\pi/4}$:
\begin{enumerate}
\item Express $z_1$ in exponential form and $z_2$ in Cartesian form.
\item Express $z_1 z_2$ and $\dfrac{z_1}{z_2}$ in both Cartesian and exponential form, stating their moduli and principal arguments.
\end{enumerate}
\end{exercice}

\begin{exercice}[Quadratic equation in $\mathbb{C}$]
\begin{enumerate}
\item Solve in $\mathbb{C}$ the equation $z^2 - 4z + 13 = 0$.
\item Show the roots on an Argand diagram and verify that they are complex conjugates.
\end{enumerate}
\end{exercice}

\begin{exercice}[Square roots of a complex number]
\begin{enumerate}
\item Find, by the algebraic method, the square roots of $7 + 24i$.
\item Check your answers by squaring.
\end{enumerate}
\end{exercice}

\begin{exercice}[Powers and roots]
\begin{enumerate}
\item Express $(1 + i)^8$ in the form $a + ib$.
\item Find the four 4th roots of $-16$ in exponential form and represent them on an Argand diagram.
\end{enumerate}
\end{exercice}

\courssub{I prepare for the GCE}

\begin{exercice}[De Moivre and multiple angles]
\begin{enumerate}
\item Use De Moivre's theorem to show that $\cos 3x = 4\cos^3 x - 3\cos x$. \hfill (4 marks)
\item Hence solve the equation $\cos 3x = 0$ for $x \in [0, \pi]$. \hfill (3 marks)
\item Using the relation $z - \dfrac{1}{z} = 2i\sin x$ where $z = \cos x + i\sin x$, express $\sin^5 x$ as a sum of sines of multiple angles. \hfill (5 marks)
\end{enumerate}
\end{exercice}

\begin{exercice}[Cube roots of unity and a cubic equation]
Let $\omega$ be a non-real cube root of unity.
\begin{enumerate}
\item Show that $1 + \omega + \omega^2 = 0$. \hfill (2 marks)
\item Prove that $(1 + \omega)(1 + \omega^2) = 1$. \hfill (2 marks)
\item Evaluate $(1 - \omega + \omega^2)^7$, giving your answer in terms of $\omega$. \hfill (4 marks)
\item Given that $2 + i$ is a root of the equation $z^3 - 7z^2 + 17z - 15 = 0$, find all the roots of the equation. \hfill (4 marks)
\end{enumerate}
\end{exercice}

\begin{exercice}[Locus in the complex plane]
The complex number $z$ satisfies $|z - 3i| = 2|z|$.
\begin{enumerate}
\item By writing $z = x + iy$, show that the locus of the point representing $z$ in the Argand diagram is a circle, and find its centre and radius. \hfill (6 marks)
\item Sketch the locus on an Argand diagram, marking the centre and the points where the circle meets the axes. \hfill (3 marks)
\item Find the greatest and least values of $|z|$ for points on this locus. \hfill (3 marks)
\end{enumerate}
\end{exercice}

\newpage

\coursec{Solutions to the exercises}

\courssub{Corrected Examination-Style Exercises}

\begin{corrige}[Exercise 1: Multiple Choice Questions]
\begin{enumerate}
\item \textbf{Answer (b):} The modulus of $z = 3 - 4i$ is $|z| = \sqrt{3^2 + (-4)^2} = \sqrt{9 + 16} = \sqrt{25} = 5$.
\item \textbf{Answer (c):} For $z = -2 + 2i$, $\operatorname{Re}(z) = -2 < 0$ and $\operatorname{Im}(z) = 2 > 0$, so $z$ lies in the second quadrant. $\tan\theta = \dfrac{2}{-2} = -1$. The principal argument is $\theta = \pi - \dfrac{\pi}{4} = \dfrac{3\pi}{4}$.
\item \textbf{Answer (b):} $(2+i)(2-i) = 2^2 - i^2 = 4 - (-1) = 5$. This illustrates the general property $z\bar z = |z|^2$ for a complex number and its conjugate.
\item \textbf{Answer (b):} The cube roots of unity are the solutions of $z^3 = 1$. They are $z = 1$ and $z = e^{\pm 2i\pi/3} = \dfrac{-1 \pm i\sqrt{3}}{2}$.
\end{enumerate}
\end{corrige}

\begin{corrige}[Exercise 2: True or False]
\begin{enumerate}
\item \textbf{True.} Every real number $a$ can be written as $a = a + 0i$, hence $\mathbb{R} \subset \mathbb{C}$.
\item \textbf{True.} If $z = a + ib$, then $z\bar z = (a+ib)(a-ib) = a^2 - (ib)^2 = a^2 + b^2 = |z|^2$.
\item \textbf{False.} The equation $z^2 + 1 = 0 \iff z^2 = -1 \iff z = \pm i$. It has \emph{two} distinct solutions in $\mathbb{C}$ (and none in $\mathbb{R}$).
\item \textbf{True} (provided $z \neq 0$). $\arg(z) = \dfrac{\pi}{2}$ means $z$ lies on the positive imaginary axis, i.e. $z = ib$ with $b > 0$; it is purely imaginary with a positive imaginary part.
\end{enumerate}
\end{corrige}

\begin{corrige}[Exercise 3: Definitions]
\begin{enumerate}
\item Let $z = a + ib \neq 0$. The \cle{modulus} of $z$ is the non-negative real number $|z| = \sqrt{a^2 + b^2}$. A \cle{principal argument} of $z$ is the unique real number $\theta \in \left]-\pi, \pi\right]$ such that $a = |z|\cos\theta$ and $b = |z|\sin\theta$; we denote it $\theta = \arg(z)$.
\item The \cle{complex conjugate} of $z = a + ib$ is $\bar z = a - ib$. It follows that $z + \bar z = 2a = 2\operatorname{Re}(z)$ and $z - \bar z = 2ib = 2i\operatorname{Im}(z)$.
\item \textbf{De Moivre's theorem:} For every real number $\theta$ and every integer $n \in \mathbb{Z}$,
\[ \bigl(\cos\theta + i\sin\theta\bigr)^n = \cos(n\theta) + i\sin(n\theta). \]
\end{enumerate}
\end{corrige}

\begin{corrige}[Exercise 4: Direct Calculations]
Given $z_1 = 2 + 3i$ and $z_2 = 1 - i$:
\begin{enumerate}
\item $z_1 + z_2 = (2+1) + (3-1)i = 3 + 2i$; \quad $z_1 - 2z_2 = (2+3i) - 2(1-i) = 2+3i -2 + 2i = 5i$.
\item $z_1 z_2 = (2+3i)(1-i) = 2 - 2i + 3i - 3i^2 = 2 + i + 3 = 5 + i$.
\item Multiply numerator and denominator by the conjugate of $z_2$:
\[ \frac{z_1}{z_2} = \frac{(2+3i)(1+i)}{(1-i)(1+i)} = \frac{2 + 2i + 3i + 3i^2}{1 - i^2} = \frac{-1 + 5i}{2} = -\frac{1}{2} + \frac{5}{2}i. \]
\item The powers of $i$ cycle with period $4$: $i^1 = i,\ i^2 = -1,\ i^3 = -i,\ i^4 = 1$. Since $2025 = 4 \times 506 + 1$, we have $i^{2025} = (i^4)^{506} \cdot i = 1^{506} \cdot i = i$.
\end{enumerate}
\end{corrige}

\begin{corrige}[Exercise 5: Cartesian and Exponential Forms]
\begin{enumerate}
\item For $z_1 = 1 + i\sqrt{3}$: $|z_1| = \sqrt{1^2 + (\sqrt{3})^2} = 2$; $\cos\theta = \frac{1}{2}$, $\sin\theta = \frac{\sqrt{3}}{2}$, so $\arg(z_1) = \frac{\pi}{3}$. Hence $z_1 = 2e^{i\pi/3}$.

For $z_2 = \sqrt{2}\,e^{i\pi/4}$: $z_2 = \sqrt{2}\left(\cos\frac{\pi}{4} + i\sin\frac{\pi}{4}\right) = \sqrt{2}\left(\frac{\sqrt{2}}{2} + i\frac{\sqrt{2}}{2}\right) = 1 + i$.

\item \textbf{Product.}
Exponential form: $z_1 z_2 = 2e^{i\pi/3}\cdot\sqrt{2}\,e^{i\pi/4} = 2\sqrt{2}\,e^{i(\pi/3+\pi/4)} = 2\sqrt{2}\,e^{7i\pi/12}$.
Thus $|z_1z_2| = 2\sqrt{2}$ and $\arg(z_1z_2) = \frac{7\pi}{12}$.

Cartesian form: $z_1z_2 = (1+i\sqrt{3})(1+i) = 1 + i + i\sqrt{3} + i^2\sqrt{3} = (1-\sqrt{3}) + (1+\sqrt{3})i$.

\textbf{Quotient.}
Exponential form: $\dfrac{z_1}{z_2} = \dfrac{2e^{i\pi/3}}{\sqrt{2}\,e^{i\pi/4}} = \sqrt{2}\,e^{i(\pi/3-\pi/4)} = \sqrt{2}\,e^{i\pi/12}$.
Thus $\left|\dfrac{z_1}{z_2}\right| = \sqrt{2}$ and $\arg\!\left(\dfrac{z_1}{z_2}\right) = \dfrac{\pi}{12}$.

Cartesian form: $\dfrac{z_1}{z_2} = \dfrac{(1+i\sqrt{3})(1-i)}{(1+i)(1-i)} = \dfrac{1 - i + i\sqrt{3} - i^2\sqrt{3}}{2} = \dfrac{(1+\sqrt{3}) + (\sqrt{3}-1)i}{2}$.
\end{enumerate}
\end{corrige}

\begin{corrige}[Exercise 6: Quadratic Equation in $\mathbb{C}$]
\begin{enumerate}
\item Solve $z^2 - 4z + 13 = 0$. Discriminant: $\Delta = (-4)^2 - 4(1)(13) = 16 - 52 = -36 = (6i)^2$.
The roots are
\[ z = \frac{4 \pm 6i}{2} = 2 \pm 3i. \]
Hence the solution set is $\mathcal{S} = \{2 - 3i,\ 2 + 3i\}$.
\item The roots $2 + 3i$ and $2 - 3i$ are complex conjugates (same real part, opposite imaginary parts). On the Argand diagram they are symmetric with respect to the real axis:
\begin{popfigure}
\begin{tikzpicture}[scale=0.9]
\draw[->,popline] (-1,0) -- (4,0) node[below left] {\small Re};
\draw[->,popline] (0,-4) -- (0,4) node[below left] {\small Im};
\foreach \x in {1,2,3} \draw (\x,0.06) -- (\x,-0.06) node[below] {\scriptsize \x};
\foreach \y in {-3,-2,-1,1,2,3} \draw (0.06,\y) -- (-0.06,\y) node[left] {\scriptsize \y};
\draw[popblue,dashed] (2,3) -- (2,-3);
\fill[poporange] (2,3) circle (2pt) node[right,popdark] {\small $2+3i$};
\fill[poporange] (2,-3) circle (2pt) node[right,popdark] {\small $2-3i$};
\end{tikzpicture}
\legende{The two conjugate roots, symmetric about the real axis.}
\end{popfigure}
\end{enumerate}
\end{corrige}

\begin{corrige}[Exercise 7: Square Roots of a Complex Number]
\begin{enumerate}
\item Let $w = x + iy$ ($x,y \in \mathbb{R}$) such that $w^2 = 7 + 24i$. Then
\[ (x+iy)^2 = x^2 - y^2 + 2ixy = 7 + 24i. \]
Equating real and imaginary parts:
\[ \begin{cases} x^2 - y^2 = 7 \\ 2xy = 24 \end{cases} \]
We also have $|w|^2 = x^2 + y^2 = |7+24i| = \sqrt{7^2+24^2} = 25$.
Adding and subtracting the first and third equations:
\[ 2x^2 = 32 \implies x^2 = 16 \implies x = \pm 4, \qquad 2y^2 = 18 \implies y^2 = 9 \implies y = \pm 3. \]
Since $2xy = 24 > 0$, $x$ and $y$ have the \emph{same} sign. Hence the two square roots are
\[ w = 4 + 3i \quad \text{or} \quad w = -4 - 3i. \]
\item Verification: $(4+3i)^2 = 16 + 24i + 9i^2 = 16 - 9 + 24i = 7 + 24i$. \checkmark
$(-4-3i)^2 = (-1)^2(4+3i)^2 = 7 + 24i$. \checkmark
\end{enumerate}
\end{corrige}

\begin{corrige}[Exercise 8: Powers and Roots]
\begin{enumerate}
\item $1 + i = \sqrt{2}\,e^{i\pi/4}$ because $|1+i| = \sqrt{2}$ and $\arg(1+i) = \frac{\pi}{4}$. By De Moivre's theorem,
\[ (1+i)^8 = \left(\sqrt{2}\right)^8 e^{8i\pi/4} = 16\,e^{2i\pi} = 16(\cos 2\pi + i\sin 2\pi) = 16. \]
So $(1+i)^8 = 16 + 0i = 16$.
\item Solve $z^4 = -16$. Write $-16 = 16\,e^{i\pi}$ (principal argument $\pi$). The four 4th roots are
\[ z_k = 16^{1/4}\,e^{i(\pi + 2k\pi)/4} = 2\,e^{i(2k+1)\pi/4}, \qquad k = 0,1,2,3. \]
Explicitly:
\[ z_0 = 2e^{i\pi/4},\quad z_1 = 2e^{3i\pi/4},\quad z_2 = 2e^{5i\pi/4} = 2e^{-3i\pi/4},\quad z_3 = 2e^{7i\pi/4} = 2e^{-i\pi/4}. \]
They lie on the circle $|z|=2$ at angles $45^\circ, 135^\circ, 225^\circ, 315^\circ$ — the vertices of a square:
\begin{popfigure}
\begin{tikzpicture}[scale=1.1]
\draw[->,popline] (-3,0) -- (3,0) node[below left] {\small Re};
\draw[->,popline] (0,-3) -- (0,3) node[below left] {\small Im};
\draw[popblue] (0,0) circle (2);
\draw[popmuted,dashed] (1.414,1.414) -- (-1.414,1.414) -- (-1.414,-1.414) -- (1.414,-1.414) -- cycle;
\fill[poporange] (1.414,1.414) circle (2pt) node[above right,popdark] {\small $z_0$};
\fill[poporange] (-1.414,1.414) circle (2pt) node[above left,popdark] {\small $z_1$};
\fill[poporange] (-1.414,-1.414) circle (2pt) node[below left,popdark] {\small $z_2$};
\fill[poporange] (1.414,-1.414) circle (2pt) node[below right,popdark] {\small $z_3$};
\end{tikzpicture}
\legende{The four 4th roots of $-16$: vertices of a square on $|z|=2$.}
\end{popfigure}
\end{enumerate}
\end{corrige}

\begin{corrige}[Exercise 9: De Moivre and Multiple Angles]
\begin{enumerate}
\item By De Moivre, $\cos 3x + i\sin 3x = (\cos x + i\sin x)^3$. Expand using the binomial theorem with $c = \cos x$, $s = \sin x$:
\[ (c + is)^3 = c^3 + 3c^2(is) + 3c(is)^2 + (is)^3 = (c^3 - 3cs^2) + i(3c^2s - s^3). \]
Equating real parts: $\cos 3x = c^3 - 3cs^2 = c^3 - 3c(1-c^2) = 4\cos^3 x - 3\cos x$. \hfill \textbf{(M1 A1 M1 A1)}
\item $\cos 3x = 0 \iff 3x = \dfrac{\pi}{2} + k\pi \iff x = \dfrac{\pi}{6} + \dfrac{k\pi}{3}$, $k \in \mathbb{Z}$.
In the interval $[0,\pi]$: $x = \dfrac{\pi}{6},\ \dfrac{\pi}{2},\ \dfrac{5\pi}{6}$. \hfill \textbf{(M1 A1 A1)}
(Consistency check: $4c^3 - 3c = c(4c^2 - 3) = 0$ gives $c = 0$ or $c = \pm\dfrac{\sqrt{3}}{2}$, matching these solutions.)
\item Let $z = \cos x + i\sin x = e^{ix}$. Then $z - \dfrac{1}{z} = 2i\sin x$ and $z^n - \dfrac{1}{z^n} = 2i\sin nx$.
\[ (2i\sin x)^5 = \left(z - \tfrac{1}{z}\right)^5 = z^5 - 5z^3 + 10z - \tfrac{10}{z} + \tfrac{5}{z^3} - \tfrac{1}{z^5}. \]
Group conjugate pairs:
\[ 32i\sin^5 x = \left(z^5 - \tfrac{1}{z^5}\right) - 5\left(z^3 - \tfrac{1}{z^3}\right) + 10\left(z - \tfrac{1}{z}\right) = 2i\sin 5x - 10i\sin 3x + 20i\sin x. \]
Divide by $32i$:
\[ \sin^5 x = \frac{1}{16}\bigl(\sin 5x - 5\sin 3x + 10\sin x\bigr). \]
\hfill \textbf{(M1 binomial, M1 pairing, A1 A1 A1)}
\end{enumerate}
\textbf{Examiner expectations:} Correct binomial expansion with careful handling of $i^2 = -1$; explicit pairing $z^n - z^{-n} = 2i\sin nx$; final answer with exact rational coefficients.
\end{corrige}

\begin{corrige}[Exercise 10: Cube Roots of Unity and a Cubic Equation]
\begin{enumerate}
\item Since $\omega^3 = 1$ and $\omega \neq 1$: $\omega^3 - 1 = 0 \Rightarrow (\omega - 1)(\omega^2 + \omega + 1) = 0$. As $\omega \neq 1$, we must have $1 + \omega + \omega^2 = 0$. \hfill \textbf{(M1 A1)}
\item $(1+\omega)(1+\omega^2) = 1 + \omega + \omega^2 + \omega^3 = 0 + 1 = 1$, using part (a) and $\omega^3 = 1$. \hfill \textbf{(M1 A1)}
\item From $1 + \omega + \omega^2 = 0$ we get $\omega^2 + 1 = -\omega$, so
\[ 1 - \omega + \omega^2 = (1 + \omega^2) - \omega = -\omega - \omega = -2\omega. \]
Hence $(1 - \omega + \omega^2)^7 = (-2\omega)^7 = -128\,\omega^7 = -128\,\omega^{6}\cdot\omega = -128\,\omega$ (since $\omega^6 = (\omega^3)^2 = 1$). \hfill \textbf{(M1 A1 M1 A1)}
\item The polynomial $z^3 - 7z^2 + 17z - 15 = 0$ has real coefficients. Given $2+i$ is a root, its conjugate $2-i$ is also a root. The quadratic factor is
\[ (z - (2+i))(z - (2-i)) = z^2 - 4z + 5. \]
Polynomial division (or comparing coefficients) gives:
\[ z^3 - 7z^2 + 17z - 15 = (z^2 - 4z + 5)(z - 3). \]
The third root is $z = 3$. The three roots are $2 + i,\ 2 - i,\ 3$. \hfill \textbf{(M1 conjugate root, M1 factorisation, A1 A1)}
\end{enumerate}
\textbf{Examiner expectations:} Explicit use of real coefficients to justify the conjugate root; clean division or comparison of coefficients; all three roots clearly stated.
\end{corrige}

\begin{corrige}[Exercise 11: Locus in the Complex Plane]
\begin{enumerate}
\item Let $z = x + iy$ ($x,y \in \mathbb{R}$). Then $|z - 3i|^2 = x^2 + (y-3)^2$ and $|z|^2 = x^2 + y^2$.
The equation $|z - 3i| = 2|z|$ has non-negative sides, so squaring is equivalent:
\begin{align*}
x^2 + (y-3)^2 &= 4(x^2 + y^2) \\
x^2 + y^2 - 6y + 9 &= 4x^2 + 4y^2 \\
0 &= 3x^2 + 3y^2 + 6y - 9 \\
x^2 + y^2 + 2y - 3 &= 0 \\
x^2 + (y+1)^2 &= 4.
\end{align*}
This is a circle with centre $(0, -1)$ (the complex number $-i$) and radius $2$. \hfill \textbf{(M1 squaring, M1 expanding, M1 completing the square, A1 A1 A1)}
\item Intercepts:
\begin{itemize}
\item Imaginary axis ($x=0$): $(y+1)^2 = 4 \Rightarrow y+1 = \pm 2 \Rightarrow y = 1$ or $y = -3$. Points: $i$ and $-3i$.
\item Real axis ($y=0$): $x^2 + 1 = 4 \Rightarrow x^2 = 3 \Rightarrow x = \pm\sqrt{3}$. Points: $\sqrt{3}$ and $-\sqrt{3}$.
\end{itemize}
\begin{popfigure}
\begin{tikzpicture}[scale=0.85]
\draw[->,popline] (-3.5,0) -- (3.5,0) node[below left] {\small Re};
\draw[->,popline] (0,-4) -- (0,2.5) node[below left] {\small Im};
\draw[popblue,thick] (0,-1) circle (2);
\fill[popdark] (0,-1) circle (2pt) node[right] {\small centre $-i$};
\fill[poporange] (0,1) circle (2pt) node[above right,popdark] {\small $i$};
\fill[poporange] (0,-3) circle (2pt) node[below right,popdark] {\small $-3i$};
\fill[poporange] (1.732,0) circle (2pt) node[above right,popdark] {\small $\sqrt{3}$};
\fill[poporange] (-1.732,0) circle (2pt) node[above left,popdark] {\small $-\sqrt{3}$};
\end{tikzpicture}
\legende{Locus: circle of centre $-i$ and radius $2$.}
\end{popfigure}
\hfill \textbf{(M1 circle drawn, A1 centre, A1 intercepts)}
\item $|z|$ is the distance from the origin to a point on the circle. The origin is at distance $1$ from the centre $-i$, and the radius is $2$. Therefore
\[ |z|_{\max} = 1 + 2 = 3 \quad \text{(at the point opposite the origin, i.e. } z = -3i\text{)}, \]
\[ |z|_{\min} = 2 - 1 = 1 \quad \text{(at the point nearest the origin, i.e. } z = i\text{)}. \]
\hfill \textbf{(M1 geometric argument, A1 A1)}
\end{enumerate}
\textbf{Examiner expectations:} Squaring justified (moduli are non-negative); completing the square shown explicitly; geometric reasoning (distance from origin = centre distance $\pm$ radius) used for max/min rather than calculus or trial and error.
\end{corrige}

\newpage

\coursec{Summary sheet}

\begin{popfigure}
\begin{tikzpicture}[scale=0.95, transform shape,
  central/.style={circle, draw=popink, fill=popblueL, thick, align=center, font=\bfseries, minimum size=2.4cm},
  branch/.style={rectangle, rounded corners, draw=#1, fill=#1!12, thick, align=center, font=\small, text width=3.4cm},
  link/.style={thick, draw=#1}]
\node[central, align=center] (C){Complex\\ Numbers\\ $\mathbb{C}$};
\node[branch=popblue, align=center] (A) at (-5.2,2.6){Birth of $i$\\ $i^2=-1$\\ $z=a+ib$};
\node[branch=popgreen, align=center] (B) at (0,3.6){Algebra\\ $+,\ -,\ \times,\ \div$\\ Conjugate $\bar z$};
\node[branch=poporange, align=center] (D) at (5.2,2.6){Argand diagram\\ Modulus $|z|$\\ Argument $\arg z$};
\node[branch=poppurple, align=center] (E) at (5.2,-2.6){Polar $r(\cos\theta+i\sin\theta)$\\ Exponential $re^{i\theta}$};
\node[branch=popgold, align=center] (F) at (0,-3.6){De Moivre\\ $z^n=r^n(\cos n\theta+i\sin n\theta)$\\ Trig. applications};
\node[branch=poppink, align=center] (G) at (-5.2,-2.6){Roots \& loci\\ $n$th roots, $\sqrt[3]{1}$\\ Circles, half-lines};
\draw[link=popblue] (C) -- (A);
\draw[link=popgreen] (C) -- (B);
\draw[link=poporange] (C) -- (D);
\draw[link=poppurple] (C) -- (E);
\draw[link=popgold] (C) -- (F);
\draw[link=poppink] (C) -- (G);
\end{tikzpicture}
\legende{Mind map of the chapter on complex numbers.}
\end{popfigure}

\begin{bilan}

\textbf{Key definitions}
\begin{itemize}
\item \cle{Imaginary unit}: $i$ is the number such that $i^2=-1$; powers cycle: $i^1=i$, $i^2=-1$, $i^3=-i$, $i^4=1$.
\item \cle{Complex number}: $z=a+ib$ with $a=\mathrm{Re}(z)$, $b=\mathrm{Im}(z)$; $z$ is real if $b=0$, purely imaginary if $a=0$.
\item \cle{Conjugate}: $\bar z=a-ib$; \cle{modulus}: $|z|=\sqrt{a^2+b^2}$; \cle{argument}: $\theta=\arg z$ satisfies $\cos\theta=\dfrac{a}{r}$, $\sin\theta=\dfrac{b}{r}$, with $r=|z|$.
\item \cle{Argand diagram}: $z=a+ib$ is represented by the point $M(a,b)$; $|z|=OM$, $\arg z$ is the angle $(Ox, \overrightarrow{OM})$.
\end{itemize}

\textbf{Laws, formulas and theorems to know}
\begin{itemize}
\item Operations: $(a+ib)+(c+id)=(a+c)+i(b+d)$; $(a+ib)(c+id)=(ac-bd)+i(ad+bc)$.
\item Division: multiply numerator and denominator by the conjugate: $\dfrac{a+ib}{c+id}=\dfrac{(a+ib)(c-id)}{c^2+d^2}$.
\item Conjugate properties: $\overline{z_1+z_2}=\bar z_1+\bar z_2$, $\overline{z_1z_2}=\bar z_1\bar z_2$, $\overline{\left(\dfrac{z_1}{z_2}\right)}=\dfrac{\bar z_1}{\bar z_2}$, $z+\bar z=2\mathrm{Re}(z)$, $z-\bar z=2i\,\mathrm{Im}(z)$, $z\bar z=|z|^2$.
\item $z$ is real $\iff z=\bar z$; $z$ is purely imaginary $\iff z=-\bar z$.
\item Modulus: $|z_1z_2|=|z_1||z_2|$, $\left|\dfrac{z_1}{z_2}\right|=\dfrac{|z_1|}{|z_2|}$, $|z^n|=|z|^n$, $|\bar z|=|z|$.
\item Argument: $\arg(z_1z_2)=\arg z_1+\arg z_2$, $\arg\!\left(\dfrac{z_1}{z_2}\right)=\arg z_1-\arg z_2$, $\arg(\bar z)=-\arg z$ (all modulo $2\pi$).
\item Polar form: $z=r(\cos\theta+i\sin\theta)$; exponential form: $z=re^{i\theta}$; Euler: $e^{i\theta}=\cos\theta+i\sin\theta$.
\item \cle{De Moivre's theorem}: $(\cos\theta+i\sin\theta)^n=\cos n\theta+i\sin n\theta$ for $n\in\mathbb{Z}$.
\item $n$th roots of $z=re^{i\theta}$: $z_k=\sqrt[n]{r}\,e^{i(\theta+2k\pi)/n}$, $k=0,1,\dots,n-1$; they lie on a circle, at the vertices of a regular $n$-gon.
\item Cube roots of unity: $1,\ \omega=e^{2i\pi/3},\ \omega^2$, with $1+\omega+\omega^2=0$ and $\omega^3=1$.
\item Useful identities: $\cos\theta=\dfrac{e^{i\theta}+e^{-i\theta}}{2}$, $\sin\theta=\dfrac{e^{i\theta}-e^{-i\theta}}{2i}$.
\item Loci: $|z-a|=r$ circle (centre $a$, radius $r$); $|z-a|=|z-b|$ perpendicular bisector; $\arg(z-a)=\theta$ half-line from $a$; $|z-a|<r$ open disc.
\end{itemize}

\textbf{Methods}
\begin{itemize}
\item To find $\arg z$: compute $r$, then solve $\cos\theta=a/r$, $\sin\theta=b/r$; check the quadrant from the signs of $a$ and $b$.
\item To solve $az^2+bz+c=0$ with $\Delta<0$: $z=\dfrac{-b\pm i\sqrt{-\Delta}}{2a}$; real-coefficient equations have conjugate roots.
\item To find square roots of $a+ib$: solve $(x+iy)^2=a+ib$, i.e.\ $x^2-y^2=a$, $2xy=b$, $x^2+y^2=\sqrt{a^2+b^2}$.
\item Powers of $\cos$, $\sin$ as multiple angles: use Euler's formulas and expand $(e^{i\theta}\pm e^{-i\theta})^n$ with the binomial theorem.
\item Multiple angles as powers: expand $(\cos\theta+i\sin\theta)^n$ by De Moivre and equate real and imaginary parts.
\end{itemize}

\textbf{Common mistakes}
\begin{itemize}
\item Writing $\sqrt{-4}=2i$ is fine, but never apply $\sqrt{ab}=\sqrt{a}\sqrt{b}$ to negative numbers.
\item Forgetting that $\arg z$ is defined only modulo $2\pi$ and that the principal value lies in $(-\pi,\pi]$.
\item Placing $\theta$ in the wrong quadrant: always check the signs of $a$ and $b$.
\item Dividing by a complex number without multiplying by the conjugate.
\item Confusing $|z|^2=z\bar z$ with $z^2$: in general $z^2\neq |z|^2$.
\item Forgetting one of the $n$ roots of an $n$th root equation: there are always exactly $n$ distinct roots.
\end{itemize}

\textbf{GCE tips}
\begin{itemize}
\item Show every step when rationalising a denominator; method marks are awarded.
\item For loci questions, state the geometric object clearly (circle, bisector, half-line) and sketch it on an Argand diagram.
\item In De Moivre questions, write the polar form first with $r>0$ and $\theta$ in the correct range.
\item Check roots by substitution or by using the sum and product of roots.
\item Manage time: a full complex-number question (forms, roots, locus) is worth many marks — attempt every part.
\end{itemize}
\end{bilan}

\findecours

\end{document}
